%! Author = crist
%! Date = 03/02/2026
% Document

\spacing{1.3}
\section{Consideraciones previas}

\subsection{Conceptos básicos}
\begin{itemize}
    \item \negsubro{Fórmula química}: es como la ``receta'' de un compuesto. Nos dice qué elementos lo forman y en qué cantidad. Usamos el símbolo químico de cada elemento y unos números pequeños (subíndices) que indican cuántos átomos hay de cada uno. Por ejemplo: \ch{Fe2O3} significa 2 átomos de hierro y 3 átomos de oxígeno.

    \item \negsubro{Formulación}: es el proceso de escribir la fórmula química cuando nos dan el nombre del compuesto. Es decir, pasar del nombre a la fórmula.

    \item \negsubro{Nomenclatura}: es el proceso inverso. Dada la fórmula química, debemos encontrar el nombre correcto del compuesto.

    \item \negsubro{Compuestos binarios}: están formados por dos elementos químicos diferentes (aunque algunos casos especiales como \ch{H2} también se consideran binarios).

    \item \negsubro{Compuestos ternarios}: están formados por tres elementos químicos diferentes.

    \item \negsubro{Metales}: son los elementos de los grupos del 1 al 12 de la tabla periódica. Ejemplos: sodio (Na), hierro (Fe), cobre (Cu), calcio (Ca).

    \item \negsubro{No metales}: son los elementos de los grupos 13 al 18. Ejemplos: oxígeno (O), nitrógeno (N), cloro (Cl), azufre (S).

    \item \negsubro{Halógenos}: son los elementos del grupo 17 (flúor, cloro, bromo, yodo).
\end{itemize}

\subsection{Estados de oxidación: la clave de la formulación}

Cuando dos elementos se unen para formar un compuesto, se produce un intercambio de electrones: uno de los elementos \textbf{cede} electrones y el otro los \textbf{recibe}.

\begin{definition}[Estados de oxidación]
        Los estados de oxidación son como una ``puntuación'' que indica cuántos electrones cede (+) o recibe (-) un elemento al unirse con otro para formar un compuesto. Un mismo elemento puede tener varios estados de oxidación dependiendo del compuesto que forme.

        \textbf{Regla importante}: El elemento con estado de oxidación \textcolor{blue}{positivo} se escribe a la \textcolor{blue}{izquierda}, y el que tiene estado de oxidación \textcolor{red}{negativo} se escribe a la \textcolor{red}{derecha}.
\end{definition}
\vspace{8pt}

\begin{example}
    El cobre (\ch{Cu}) puede unirse con el oxígeno (\ch{O}) de dos formas diferentes:
    \begin{enumerate}[label=\letraminuscula]
            \item \textbf{Formando \ch{CuO}}:
            \begin{itemize}
                \item El \ch{Cu} cede 2 electrones $\rightarrow$ estado de oxidación \textcolor{blue}{+2}
                \item El \ch{O} recibe 2 electrones $\rightarrow$ estado de oxidación \textcolor{red}{-2}
            \end{itemize}

            \item \textbf{Formando \ch{Cu2O}}:
            \begin{itemize}
                \item El \ch{Cu} cede 1 electrón $\rightarrow$ estado de oxidación \textcolor{blue}{+1}
                \item El \ch{O} recibe 2 electrones $\rightarrow$ estado de oxidación \textcolor{red}{-2}
            \end{itemize}
    \end{enumerate}
\end{example}

\subsection{Cómo usar los estados de oxidación para formular}

\textbf{Durante este curso trabajaremos con los estados de oxidación en valor absoluto} (sin el signo) para intercambiarlos y crear las fórmulas. El proceso es el siguiente:

\begin{example}
    Vamos a formar el compuesto entre el sodio (\ch{Na}) y el oxígeno (\ch{O}):

    \begin{enumerate}[label=\letraminuscula]
            \item Escribimos el \ch{Na} (estado de oxidación +1) a la izquierda y el \ch{O} (estado de oxidación -2) a la derecha:

            \eo{blue}{1}{Na} \eo{red}{2}{O}

            \item \textbf{Intercambiamos los estados de oxidación} (en valor absoluto): el 1 del Na pasa como subíndice al O, y el 2 del O pasa como subíndice al Na:

            \ch{Na2O1}

            \item Como el subíndice 1 no se escribe, la fórmula final es: \ch{Na2O}

            \item \textbf{Comprobación}: $2 \times (+1) + 1 \times (-2) = +2 - 2 = 0$ $\checkmark$ (la suma debe dar siempre cero)
    \end{enumerate}
\end{example}

\begin{remark}
    \textbf{\textcolor{red}{$\triangleright$ Regla de oro}}: Los estados de oxidación se intercambian en valor absoluto y pasan a ser los subíndices del otro elemento.
\end{remark}


%    \noindent\fcolorbox{red}{yellow!20}{%
%        \begin{minipage}{0.95\textwidth}
%            \textbf{\textcolor{red}{$\triangleright$ Regla de oro}}: Los estados de oxidación se intercambian en valor absoluto y pasan a ser los subíndices del otro elemento.
%        \end{minipage}
%    }

\textbf{Nota importante}: Cuando un elemento tiene varios estados de oxidación posibles, el ejercicio nos indicará cuál usar, o lo deduciremos de los subíndices de la fórmula.

\subsection{Prefijos multiplicadores}

Los prefijos son palabras que ponemos delante del nombre del elemento para indicar el número de átomos de ese elemento en la fórmula.

\begin{multicols}{2}
    \begin{itemize}
        \item \textbf{Mono}: 1 átomo. \textit{Normalmente se omite}.
        \item \textbf{Di}: 2 átomos
        \item \textbf{Tri}: 3 átomos
        \item \textbf{Tetra}: 4 átomos
        \item \textbf{Penta}: 5 átomos
        \item \textbf{Hexa}: 6 átomos
        \item \textbf{Hepta}: 7 átomos
        \item \textbf{Octa}: 8 átomos
    \end{itemize}
\end{multicols}

\begin{example}
    Aquí tenemos varios ejemplos, el prefijo -mono se suele omitir pero en el caso del \ch{CO} no se omite, es una de las excepciones.
    \begin{itemize}
        \item \ch{CO}: monóxido de carbono (1 oxígeno)
        \item \ch{CO2}: dióxido de carbono (2 oxígenos)
        \item \ch{N2O3}: trióxido de dinitrógeno (3 oxígenos, 2 nitrógenos)
    \end{itemize}
\end{example}

\textbf{Nota}: Los nombres de los prefijos son iguales en inglés y en español.

\subsection{Tipos de nomenclatura}

La IUPAC (Unión Internacional de Química Pura y Aplicada) permite nombrar los compuestos de varias formas. En este curso aprenderemos dos nomenclaturas principales:

\begin{enumerate}[label=\letraminuscula]
        \item \textbf{Nomenclatura de multiplicadores} (o nomenclatura sistemática): usa los prefijos (mono, di, tri...) para indicar el número de átomos.

        \item \textbf{Nomenclatura con estados de oxidación}: indica el estado de oxidación del elemento mediante números romanos entre paréntesis.
\end{enumerate}

\begin{example}
    El compuesto \ch{Fe2O3} se puede nombrar de dos formas:
    \begin{itemize}
        \item \textbf{Nomenclatura de multiplicadores}: trióxido de dihierro
        \item \textbf{Nomenclatura con estados de oxidación}: óxido de hierro (III)
    \end{itemize}
\end{example}

%\subsection{Diagrama para identificar tipos de compuestos}
%
%Antes de formular o nombrar, es importante identificar qué tipo de compuesto tenemos. Usa esta guía paso a paso:
%
%\vspace{0.5cm}
%
%\begin{center}
%    \fcolorbox{black}{blue!10}{%
%        \begin{minipage}{0.85\textwidth}
%            \textbf{\large Paso 1:} ¿Tiene grupo \ch{OH}?
%            \begin{itemize}
%                \item \textbf{SÍ} → Es un \textcolor{blue}{\textbf{HIDRÓXIDO}}
%                \item \textbf{NO} → Pasa al Paso 2
%            \end{itemize}
%        \end{minipage}
%    }
%\end{center}
%
%\vspace{0.3cm}
%
%\begin{center}
%    \fcolorbox{black}{green!10}{%
%        \begin{minipage}{0.85\textwidth}
%            \textbf{\large Paso 2:} ¿Tiene oxígeno?
%            \begin{itemize}
%                \item \textbf{SÍ} → Pasa al Paso 3
%                \item \textbf{NO} → Pasa al Paso 4
%            \end{itemize}
%        \end{minipage}
%    }
%\end{center}
%
%\vspace{0.3cm}
%
%\begin{center}
%    \fcolorbox{black}{yellow!20}{%
%        \begin{minipage}{0.85\textwidth}
%            \textbf{\large Paso 3:} ¿El oxígeno lleva subíndice 2 que no se puede simplificar?
%            \begin{itemize}
%                \item \textbf{SÍ} (ej: \ch{Na2O2}, \ch{CaO2}) → Es un \textcolor{red}{\textbf{PERÓXIDO}}
%                \item \textbf{NO} (ej: \ch{CO2}, \ch{FeO}) → Es un \textcolor{blue}{\textbf{ÓXIDO}}
%            \end{itemize}
%        \end{minipage}
%    }
%\end{center}
%
%\vspace{0.3cm}
%
%\begin{center}
%    \fcolorbox{black}{orange!20}{%
%        \begin{minipage}{0.85\textwidth}
%            \textbf{\large Paso 4:} ¿Tiene hidrógeno?
%            \begin{itemize}
%                \item \textbf{SÍ} → Es un \textcolor{purple}{\textbf{HIDRURO}} (o \textbf{HIDRÁCIDO} si tiene (ac) o (aq))
%                \item \textbf{NO} → Es una \textcolor{green!50!black}{\textbf{SAL BINARIA}} (metal + no metal)
%            \end{itemize}
%        \end{minipage}
%    }
%\end{center}
%
%\vspace{0.5cm}

\section{Óxidos}

\subsection{¿Qué son los óxidos?}

Los óxidos son compuestos formados por un elemento (metal o no metal) y oxígeno.

\begin{itemize}
    \item El oxígeno \textbf{siempre} va a la derecha en la fórmula
    \item El oxígeno tiene estado de oxidación \textbf{-2} (en óxidos normales)
    \item Fórmula general: \ch{E_xO_y} (donde E es cualquier elemento)
\end{itemize}

\subsection{Nomenclatura de multiplicadores}

\begin{itemize}
    \item \noindent\worldflag[width=3.9mm,stretch=1]{ES}: \textcolor{blue}{prefijo + óxido} + de + \textcolor{red}{prefijo + nombre del elemento}

    \textbf{Ejemplos}:

    \ballnumber{1}: \ch{Na2O}: óxido de disodio (o monóxido de disodio)

    \ballnumber{2}: \ch{SO3}: trióxido de azufre

    \ballnumber{3}: \ch{Al2O3}: trióxido de dialuminio


    \item \noindent\worldflag[width=3.0mm,stretch=1]{GB}: \textcolor{red}{prefix + name of the element} + \textcolor{blue}{prefix + oxide}

    \textbf{ejemplos}:

    \ballnumber{1}: \ch{CaO}: calcium oxide (or calcium monoxide)

    \ballnumber{2}: \ch{CO2}: carbon dioxide
\end{itemize}

\subsection{Nomenclatura con estados de oxidación}

\begin{itemize}
    \item \noindent\worldflag[width=3.9mm,stretch=1]{ES}: \textcolor{blue}{óxido} + de + \textcolor{red}{nombre del elemento} $\Bigl($\textcolor{green!50!black}{estado de oxidación}$\Bigr)$

    \textbf{Regla importante}: Si el elemento solo tiene un estado de oxidación, \textbf{NO} se pone el número romano.

    \textbf{Ejemplos}:

    \ballnumber{1}: \ch{MgO}: óxido de magnesio (el Mg solo tiene +2, no hace falta indicarlo)

    \ballnumber{2}: \ch{FeO}: óxido de hierro (II)

    \ballnumber{3}: \ch{Fe2O3}: óxido de hierro (III)


    \item \noindent\worldflag[width=3.0mm,stretch=1]{GB}: \textcolor{red}{name of the element} $\Bigl($\textcolor{green!50!black}{oxidation state}$\Bigr)$ + \textcolor{blue}{oxide}

    \textbf{ejemplos}:

    \ballnumber{1}: \ch{CuO}: copper (II) oxide

    \ballnumber{2}: \ch{Cu2O}: copper (I) oxide
\end{itemize}

\subsection{Cómo formular óxidos}

\textbf{Pasos para formular}:

\begin{enumerate}
    \item Escribe el símbolo del elemento a la izquierda
    \item Escribe el símbolo del oxígeno (O) a la derecha
    \item Intercambia los estados de oxidación como subíndices
    \item Simplifica los subíndices si es posible (divide por el máximo común divisor)
\end{enumerate}

\begin{example}
    \textbf{Formular}: Óxido de aluminio

    \begin{itemize}
        \item Al tiene estado +3, O tiene estado -2
        \item Escribimos: \eo{blue}{3}{Al} \eo{red}{2}{O}
        \item Intercambiamos: \ch{Al2O3}
        \item No se puede simplificar → \textbf{Fórmula final}: \ch{Al2O3}
    \end{itemize}
\end{example}


\begin{remark}
    \textbf{\textcolor{red}{¡Atención!}} Cuando simplifiquemos, debemos dividir \textbf{todos} los subíndices por el mismo número (siempre será 2).

    \textbf{Ejemplo}: \ch{Ca2O2} se simplifica a \ch{CaO} (dividimos ambos subíndices entre 2)
\end{remark}

%    \noindent\fcolorbox{red}{yellow!20}{%
%        \begin{minipage}{0.95\textwidth}
%            \textbf{\textcolor{red}{¡Atención!}} Cuando simplifiquemos, debemos dividir \textbf{todos} los subíndices por el mismo número.
%
%            \textbf{Ejemplo}: \ch{Ca2O2} se simplifica a \ch{CaO} (dividimos ambos subíndices entre 2)
%        \end{minipage}
%    }


%\begin{exercise}
%    Nombra los siguientes óxidos con la nomenclatura de multiplicadores y estados de oxidación.
%    \begin{multicols}{4}
%        \begin{enumerate}[label=\letraminuscula]
%            \item
%        \end{enumerate}
%    \end{multicols}
%\end{exercise}
%
%\begin{exercise}
%    Escribe la fórmula de los siguientes compuestos..
%    \begin{multicols}{4}
%        \begin{enumerate}[label=\letraminuscula]
%            \item
%        \end{enumerate}
%    \end{multicols}
%\end{exercise}


\section{Hidruros}

\subsection{¿Qué son los hidruros?}

Los hidruros son compuestos formados por hidrógeno y otro elemento.

\begin{itemize}
    \item Cuando el hidrógeno se une a un \textbf{metal}, tiene estado de oxidación \textbf{-1} → Hidruros metálicos
    \item Cuando el hidrógeno se une a un \textbf{no metal}, tiene estado de oxidación \textbf{+1} → Hidruros no metálicos
\end{itemize}

\subsection{Hidruros metálicos}

En los hidruros metálicos:
\begin{itemize}
    \item El metal va a la \textbf{izquierda}
    \item El hidrógeno va a la \textbf{derecha}
    \item El hidrógeno tiene estado de oxidación \textbf{-1}
\end{itemize}

\subsubsection{Nomenclatura de multiplicadores}

\begin{itemize}
    \item \noindent\worldflag[width=3.9mm,stretch=1]{ES}: \textcolor{blue}{prefijo + hidruro} + de + \textcolor{red}{prefijo + nombre del metal}

    \textbf{Ejemplos}:

    \ballnumber{1}: \ch{NaH}: hidruro de sodio

    \ballnumber{2}: \ch{CaH2}: dihidruro de calcio

    \ballnumber{3}: \ch{AlH3}: trihidruro de aluminio


    \item \noindent\worldflag[width=3.0mm,stretch=1]{GB}: \textcolor{red}{prefix + name of the metal} + \textcolor{blue}{prefix + hydride}

    \textbf{ejemplos}:

    \ballnumber{1}: \ch{LiH}: lithium hydride

    \ballnumber{2}: \ch{FeH3}: iron trihydride
\end{itemize}

\subsubsection{Nomenclatura con estados de oxidación}

\begin{itemize}
    \item \noindent\worldflag[width=3.9mm,stretch=1]{ES}: \textcolor{blue}{hidruro} + de + \textcolor{red}{nombre del metal} $\Bigl($\textcolor{green!50!black}{estado de oxidación}$\Bigr)$

    \textbf{Ejemplos}:

    \ballnumber{1}: \ch{KH}: hidruro de potasio (el K solo tiene +1)

    \ballnumber{2}: \ch{FeH2}: hidruro de hierro (II)

    \ballnumber{3}: \ch{FeH3}: hidruro de hierro (III)


    \item \noindent\worldflag[width=3.0mm,stretch=1]{GB}: \textcolor{red}{name of the metal} $\Bigl($\textcolor{green!50!black}{oxidation state}$\Bigr)$ + \textcolor{blue}{hydride}

    \textbf{ejemplos}:

    \ballnumber{1}: \ch{CuH}: copper (I) hydride

    \ballnumber{2}: \ch{CuH2}: copper (II) hydride
\end{itemize}

\subsubsection{Cómo formular hidruros metálicos}

\begin{enumerate}
    \item Escribe el símbolo del metal a la izquierda
    \item Escribe H a la derecha
    \item Intercambia los estados de oxidación como subíndices
    \item Simplifica si es posible
\end{enumerate}

\begin{example}
    \textbf{Formular}: Dihidruro de magnesio

    \begin{itemize}
        \item Mg tiene estado +2, H tiene estado -1
        \item Escribimos: \eo{blue}{2}{Mg} \eo{red}{1}{H}
        \item Intercambiamos: \ch{MgH2}
        \item No se puede simplificar → \textbf{Fórmula final}: \ch{MgH2}
    \end{itemize}
\end{example}

\subsection{Hidruros no metálicos (Hidrácidos)}

Los hidruros no metálicos son compuestos del hidrógeno con elementos del grupo 16 (S, Se, Te) y del grupo 17 (F, Cl, Br, I).

\begin{itemize}
    \item El hidrógeno va a la \textbf{izquierda}
    \item El no metal va a la \textbf{derecha}
    \item El hidrógeno tiene estado de oxidación \textbf{+1}
\end{itemize}

\textbf{Característica especial}: Cuando estos compuestos se disuelven en agua, forman \textbf{ácidos} y entonces se llaman \textbf{hidrácidos}. Para indicar que están en disolución acuosa, añadimos \textbf{(ac)} en español o \textbf{(aq)} en inglés después de la fórmula.

\subsubsection{Nomenclatura de multiplicadores}

\textbf{Como compuesto puro (sin disolver en agua)}:

\begin{itemize}
    \item \noindent\worldflag[width=3.9mm,stretch=1]{ES}: \textcolor{blue}{prefijo + nombre del no metal terminado en -uro} + de + hidrógeno

    \textbf{Ejemplos}:

    \ballnumber{1}: \ch{HF}: fluoruro de hidrógeno

    \ballnumber{2}: \ch{H2S}: sulfuro de hidrógeno (o sulfuro de dihidrógeno)

    \ballnumber{3}: \ch{HCl}: cloruro de hidrógeno


    \item \noindent\worldflag[width=3.0mm,stretch=1]{GB}: \textcolor{blue}{hydrogen + prefix + name of nonmetal ending in -ide}

    \textbf{ejemplos}:

    \ballnumber{1}: \ch{HBr}: hydrogen bromide

    \ballnumber{2}: \ch{H2Se}: hydrogen selenide
\end{itemize}

\textbf{Como ácido (disuelto en agua - hidrácido)}:

\begin{itemize}
    \item \noindent\worldflag[width=3.9mm,stretch=1]{ES}: \textcolor{blue}{ácido + raíz del no metal + -hídrico}

    \textbf{Ejemplos}:

    \ballnumber{1}: \ch{HCl(ac)}: ácido clorhídrico

    \ballnumber{2}: \ch{H2S(ac)}: ácido sulfhídrico

    \ballnumber{3}: \ch{HF(ac)}: ácido fluorhídrico


    \item \noindent\worldflag[width=3.0mm,stretch=1]{GB}: \textcolor{blue}{hydro + root of nonmetal + -ic} + acid

    \textbf{ejemplos}:

    \ballnumber{1}: \ch{HCl(aq)}: hydrochloric acid

    \ballnumber{2}: \ch{H2S(aq)}: hydrosulfuric acid
\end{itemize}

\vspace{0.3cm}
\singlespacing
\begin{remark}
    \textbf{\textcolor{red}{¡No te confundas!}}
    \begin{itemize}
        \item \ch{HCl} (gas) = cloruro de hidrógeno
        \item \ch{HCl(ac)} (en agua) = ácido clorhídrico
    \end{itemize}
    ¡Es el mismo compuesto, pero el nombre cambia según esté disuelto o no!
\end{remark}
\spacing{1.3}
%    \noindent\fcolorbox{red}{yellow!20}{%
%        \begin{minipage}{0.95\textwidth}
%            \textbf{\textcolor{red}{¡No te confundas!}}
%            \begin{itemize}
%                \item \ch{HCl} (gas) = cloruro de hidrógeno
%                \item \ch{HCl(ac)} (en agua) = ácido clorhídrico
%            \end{itemize}
%            ¡Es el mismo compuesto, pero el nombre cambia según esté disuelto o no!
%        \end{minipage}
%    }
\vspace{0.3cm}

\subsubsection{Nomenclatura con estados de oxidación}

Los hidrácidos \textbf{NO} se nombran con estados de oxidación. Solo se usa la nomenclatura de multiplicadores.

\subsubsection{Cómo formular hidruros no metálicos}

\begin{enumerate}
    \item Escribe H a la izquierda
    \item Escribe el símbolo del no metal a la derecha
    \item Intercambia los estados de oxidación como subíndices
    \item Simplifica si es posible
    \item Si el compuesto está en disolución acuosa, añade (ac) o (aq)
\end{enumerate}

\begin{example}
    \textbf{Formular}: Ácido bromhídrico

    \begin{itemize}
        \item Es un hidrácido → H a la izquierda, Br a la derecha
        \item H tiene estado +1, Br tiene estado -1
        \item Intercambiamos: \ch{H1Br1} → \ch{HBr}
        \item Como es un ácido (está en disolución), añadimos (ac)
        \item \textbf{Fórmula final}: \ch{HBr(ac)}
    \end{itemize}
\end{example}


\section{Sales binarias (combinaciones de metal y no metal)}

\subsection{¿Qué son las sales binarias?}

Las sales binarias son compuestos formados por un metal y un no metal que no sean los casos anteriores (óxidos, hidruros, hidróxidos o peróxidos).

\begin{itemize}
    \item El metal va a la \textbf{izquierda}
    \item El no metal va a la \textbf{derecha}
    \item El no metal usa su estado de oxidación negativo
\end{itemize}

\subsection{Nomenclatura de multiplicadores}

\begin{itemize}
    \item \noindent\worldflag[width=3.9mm,stretch=1]{ES}: \textcolor{blue}{prefijo + raíz del no metal + -uro} + de + \textcolor{red}{prefijo + nombre del metal}

    \textbf{Ejemplos}:

    \ballnumber{1}: \ch{NaCl}: cloruro de sodio (la sal común)

    \ballnumber{2}: \ch{Al4C3}: tricarburo de tetraaluminio

    \ballnumber{3}: \ch{CaF2}: difluoruro de calcio


    \item \noindent\worldflag[width=3.0mm,stretch=1]{GB}: \textcolor{red}{prefix + name of the metal} + \textcolor{blue}{prefix + root of nonmetal + -ide}

    \textbf{ejemplos}:

    \ballnumber{1}: \ch{NaF}: sodium fluoride

    \ballnumber{2}: \ch{CaCl2}: calcium dichloride

    \ballnumber{3}: \ch{FeS}: iron sulfide
\end{itemize}

\subsection{Nomenclatura con estados de oxidación}

\begin{itemize}
    \item \noindent\worldflag[width=3.9mm,stretch=1]{ES}: \textcolor{blue}{raíz del no metal + -uro} + de + \textcolor{red}{nombre del metal} $\Bigl($\textcolor{green!50!black}{estado de oxidación}$\Bigr)$

    \textbf{Ejemplos}:

    \ballnumber{1}: \ch{NaCl}: cloruro de sodio (el Na solo tiene +1)

    \ballnumber{2}: \ch{FeCl2}: cloruro de hierro (II)

    \ballnumber{3}: \ch{FeCl3}: cloruro de hierro (III)


    \item \noindent\worldflag[width=3.0mm,stretch=1]{GB}: \textcolor{red}{name of the metal} $\Bigl($\textcolor{green!50!black}{oxidation state}$\Bigr)$ + \textcolor{blue}{root of nonmetal + -ide}

    \textbf{ejemplos}:

    \ballnumber{1}: \ch{CuCl}: copper (I) chloride

    \ballnumber{2}: \ch{CuCl2}: copper (II) chloride
\end{itemize}

\subsection{Cómo formular sales binarias}

\begin{enumerate}
    \item Escribe el símbolo del metal a la izquierda
    \item Escribe el símbolo del no metal a la derecha
    \item Intercambia los estados de oxidación como subíndices
    \item Simplifica si es posible
\end{enumerate}

\begin{example}
    \textbf{Formular}: Sulfuro de aluminio

    \begin{itemize}
        \item Al tiene estado +3, S tiene estado -2
        \item Escribimos: \eo{blue}{3}{Al} \eo{red}{2}{S}
        \item Intercambiamos: \ch{Al2S3}
        \item No se puede simplificar → \textbf{Fórmula final}: \ch{Al2S3}
    \end{itemize}
\end{example}


\section{Hidróxidos}

\subsection{¿Qué son los hidróxidos?}

Los hidróxidos son compuestos ternarios (tres elementos) formados por:
\begin{itemize}
    \item Un metal (o catión)
    \item El grupo hidróxido: \ch{OH-}
\end{itemize}

\textbf{Reglas importantes}:
\begin{itemize}
    \item El grupo \ch{OH} es \textbf{indivisible}, funciona como una unidad
    \item El grupo \ch{OH} tiene estado de oxidación \textbf{-1}
    \item Fórmula general: \ch{M(OH)_n}, donde M es el metal y $n$ es el estado de oxidación del metal
\end{itemize}

\subsection{Nomenclatura de multiplicadores}

\begin{itemize}
    \item \noindent\worldflag[width=3.9mm,stretch=1]{ES}: \textcolor{blue}{prefijo + hidróxido} + de + \textcolor{red}{nombre del metal}

    \textbf{Ejemplos}:

    \ballnumber{1}: \ch{LiOH}: hidróxido de litio

    \ballnumber{2}: \ch{Ca(OH)2}: dihidróxido de calcio

    \ballnumber{3}: \ch{Al(OH)3}: trihidróxido de aluminio


    \item \noindent\worldflag[width=3.0mm,stretch=1]{GB}: \textcolor{red}{name of the metal} + \textcolor{blue}{prefix + hydroxide}

    \textbf{ejemplos}:

    \ballnumber{1}: \ch{NaOH}: sodium hydroxide

    \ballnumber{2}: \ch{Ca(OH)2}: calcium dihydroxide
\end{itemize}

\subsection{Nomenclatura con estados de oxidación}

\begin{itemize}
    \item \noindent\worldflag[width=3.9mm,stretch=1]{ES}: \textcolor{blue}{hidróxido} + de + \textcolor{red}{nombre del metal} $\Bigl($\textcolor{green!50!black}{estado de oxidación}$\Bigr)$

    \textbf{Ejemplos}:

    \ballnumber{1}: \ch{NaOH}: hidróxido de sodio (el Na solo tiene +1, se omite)

    \ballnumber{2}: \ch{Fe(OH)2}: hidróxido de hierro (II)

    \ballnumber{3}: \ch{Fe(OH)3}: hidróxido de hierro (III)


    \item \noindent\worldflag[width=3.0mm,stretch=1]{GB}: \textcolor{red}{name of the metal} $\Bigl($\textcolor{green!50!black}{oxidation state}$\Bigr)$ + \textcolor{blue}{hydroxide}

    \textbf{ejemplos}:

    \ballnumber{1}: \ch{KOH}: potassium hydroxide (potassium (I) hydroxide)

    \ballnumber{2}: \ch{Fe(OH)2}: iron (II) hydroxide
\end{itemize}

\subsection{Cómo formular hidróxidos}

\begin{enumerate}
    \item Escribe el símbolo del metal a la izquierda
    \item Escribe \ch{OH} a la derecha
    \item El estado de oxidación del metal será el subíndice del grupo \ch{OH}
    \item Si el subíndice es mayor que 1, el \ch{OH} debe ir entre paréntesis: \ch{(OH)_n}
    \item Si el subíndice es 1, no hace falta paréntesis
\end{enumerate}

\begin{example}
    \textbf{Formular}: Trihidróxido de hierro

    \begin{itemize}
        \item ``Tri'' = 3 grupos OH
        \item Escribimos: \ch{Fe(OH)3}
        \item \textbf{Fórmula final}: \ch{Fe(OH)3}
    \end{itemize}

    \vspace{0.3cm}

    \textbf{Formular}: Hidróxido de cesio

    \begin{itemize}
        \item No hay prefijo = 1 grupo OH
        \item Escribimos: \ch{CsOH} (sin paréntesis)
        \item \textbf{Fórmula final}: \ch{CsOH}
    \end{itemize}
\end{example}

\vspace{0.3cm}
\singlespacing
\begin{remark}
    \textbf{\textcolor{red}{¡Atención con los paréntesis!}}
    \begin{itemize}
        \item 1 grupo OH → \ch{NaOH} (sin paréntesis)
        \item 2 o más grupos OH → \ch{Ca(OH)2}, \ch{Al(OH)3} (con paréntesis)
    \end{itemize}
\end{remark}
\spacing{1.3}

\vspace{0.3cm}


\section{Peróxidos}

\subsection{¿Qué son los peróxidos?}

Los peróxidos son compuestos muy parecidos a los óxidos, pero con una característica especial:

\begin{itemize}
    \item Están formados por un metal y oxígeno
    \item El oxígeno \textbf{siempre} lleva el subíndice 2: \ch{O2}
    \item Este subíndice \textbf{NUNCA} se puede simplificar
    \item Se forman principalmente con metales alcalinos (grupo 1) y alcalinotérreos (grupo 2)
\end{itemize}

\textbf{Fórmulas generales}:
\begin{itemize}
    \item \textbf{Grupo 1 (Alcalinos)}: \ch{M2O2}
    \item \textbf{Grupo 2 (Alcalinotérreos)}: \ch{MO2}
\end{itemize}

\subsection{Nomenclatura de peróxidos}

Aunque se pueden nombrar con la nomenclatura de multiplicadores, lo más común y recomendado es usar la \textbf{nomenclatura específica de peróxidos}:

\begin{itemize}
    \item \noindent\worldflag[width=3.9mm,stretch=1]{ES}: \textcolor{blue}{peróxido} + de + \textcolor{red}{nombre del metal}

    \textbf{Ejemplos}:

    \ballnumber{1}: \ch{Na2O2}: peróxido de sodio

    \ballnumber{2}: \ch{CaO2}: peróxido de calcio

    \ballnumber{3}: \ch{K2O2}: peróxido de potasio


    \item \noindent\worldflag[width=3.0mm,stretch=1]{GB}: \textcolor{red}{name of the metal} + \textcolor{blue}{peroxide}

    \textbf{Examples}:

    \ballnumber{1}: \ch{Li2O2}: lithium peroxide

    \ballnumber{2}: \ch{MgO2}: magnesium peroxide

    \ballnumber{3}: \ch{BaO2}: barium peroxide
\end{itemize}

\subsection{Cómo formular peróxidos}

\begin{enumerate}
    \item Identifica el grupo del metal (grupo 1 o grupo 2)
    \item Escribe \ch{O2} a la derecha (siempre)
    \item Escribe el metal a la izquierda:
    \begin{itemize}
        \item Si es del \textbf{grupo 1}: \ch{M2O2}
        \item Si es del \textbf{grupo 2}: \ch{MO2}
    \end{itemize}
\end{enumerate}

\begin{example}
    \textbf{Formular}: Peróxido de litio

    \begin{itemize}
        \item El litio (Li) está en el grupo 1
        \item Escribimos: \ch{Li2O2}
        \item \textbf{¡No simplificamos!} → \textbf{Fórmula final}: \ch{Li2O2}
    \end{itemize}

    \vspace{0.3cm}

    \textbf{Formular}: Peróxido de berilio

    \begin{itemize}
        \item El berilio (Be) está en el grupo 2
        \item Escribimos: \ch{BeO2}
        \item \textbf{Fórmula final}: \ch{BeO2}
    \end{itemize}
\end{example}

\vspace{0.3cm}
\begin{remark}
    \textbf{\textcolor{red}{¡Cuidado!}} No confundas:
    \begin{itemize}
        \item \ch{Na2O} = óxido de sodio (se puede simplificar si fuera necesario)
        \item \ch{Na2O2} = peróxido de sodio (¡NUNCA se simplifica!)
    \end{itemize}
\end{remark}
%    \noindent\fcolorbox{red}{yellow!20}{%
%        \begin{minipage}{0.95\textwidth}
%            \textbf{\textcolor{red}{¡Cuidado!}} No confundas:
%            \begin{itemize}
%                \item \ch{Na2O} = óxido de sodio (se puede simplificar si fuera necesario)
%                \item \ch{Na2O2} = peróxido de sodio (¡NUNCA se simplifica!)
%            \end{itemize}
%        \end{minipage}
%    }
\vspace{0.3cm}


\section{Trucos y consejos para no equivocarte}

\subsection{Errores más comunes}

\begin{enumerate}
    \item \textbf{Confundir óxido con peróxido}:
    \begin{itemize}
        \item Si tiene \ch{O2} que no se puede simplificar → peróxido
        \item Si el oxígeno tiene cualquier otro subíndice → óxido
    \end{itemize}

    \item \textbf{Olvidar los paréntesis en hidróxidos}:
    \begin{itemize}
        \item \ch{Ca(OH)2} $\checkmark$ correcto
        \item \ch{CaOH2} $\times$ incorrecto
    \end{itemize}

    \item \textbf{No simplificar cuando se puede}:
    \begin{itemize}
        \item \ch{Ca2O2} → simplificar a \ch{CaO} $\checkmark$
        \item \ch{Na2O2} → NO simplificar (es peróxido) $\checkmark$
    \end{itemize}

    \item \textbf{Poner (ac) en compuestos que no son ácidos}:
    \begin{itemize}
        \item \ch{HCl(ac)} $\checkmark$ (es un hidrácido)
        \item \ch{NaCl(ac)} $\times$ (no es un ácido)
    \end{itemize}

    \item \textbf{Confundir hidruro con hidróxido}:
    \begin{itemize}
        \item Hidruro: solo metal + H → \ch{NaH}
        \item Hidróxido: metal + OH → \ch{NaOH}
    \end{itemize}
\end{enumerate}

\subsection{Reglas nemotécnicas}

\begin{itemize}
    \item \textbf{``URO'' para no metales}: cuando combinamos metal + no metal → termina en -uro (cloruro, sulfuro, bromuro...)

    \item \textbf{``IDO'' con oxígeno}: cuando hay oxígeno en el compuesto → termina en -ido (óxido, peróxido, hidróxido)

    \item \textbf{El oxígeno es mandón}: siempre va a la derecha (excepto con el flúor, pero eso no lo veremos este año)

    \item \textbf{Los peróxidos son gemelos}: el \ch{O2} en peróxidos es como gemelos unidos, no se pueden separar (no se simplifica)

    \item \textbf{El grupo OH es un equipo}: siempre van juntos, por eso los ponemos entre paréntesis cuando hay más de uno
\end{itemize}

\newpage


\section{Tablas de interés}

\subsection{Tabla de estados de oxidación más comunes}

\begin{tcolorbox}[enhanced,
    fonttitle=\bfseries,
%fontupper=\normalsize\sffamily,
    colback=white, colframe=blue2, colbacktitle=blue5, coltitle=white,
    center title, clip upper,
    title=Estados de oxidación más comunes,
    tabularray={
        colspec={X|[0.25mm]X|[0.25mm]X|[0.25mm]},
        column{2-Z}={halign=c},
        row{1}={font=\bfseries,bg=blue9},
        row{1-Z}={halign=c},
    }
]
    \textbf{Elemento} & \textbf{Símbolo} & \textbf{Estados de oxidación}         \\
    \hline
    \SetCell[c=3]{c, font=\bfseries, bg=blue!5} METALES   \\
    \hline
    Sodio             & Na               & +1                                    \\
    Potasio           & K                & +1                                    \\
    Litio             & Li               & +1                                    \\
    Plata             & Ag               & +1                                    \\
    \hline
    Calcio            & Ca               & +2                                    \\
    Magnesio          & Mg               & +2                                    \\
    Bario             & Ba               & +2                                    \\
    Zinc              & Zn               & +2                                    \\
    \hline
    Aluminio          & Al               & +3                                    \\
    \hline
    Hierro            & Fe               & +2, +3                                \\
    Cobre             & Cu               & +1, +2                                \\
    Plomo             & Pb               & +2, +4                                \\
    Estaño            & Sn               & +2, +4                                \\
    \hline
    \SetCell[c=3]{c, font=\bfseries, bg=blue!5} NO METALES \\
    \hline
    Oxígeno           & O                & -2 (casi siempre)                     \\
    Hidrógeno         & H                & +1 (con no metales), -1 (con metales) \\
    Flúor             & F                & -1                                    \\
    Cloro             & Cl               & -1, +1, +3, +5, +7                    \\
    Bromo             & Br               & -1, +1, +3, +5                        \\
    Yodo              & I                & -1, +1, +3, +5, +7                    \\
    Azufre            & S                & -2, +2, +4, +6                        \\
    Nitrógeno         & N                & -3, +1, +2, +3, +4, +5                \\
    \hline
\end{tcolorbox}
%
%    \begin{center}
%        \begin{tabular}{|c|c|c|}
%            \hline
%            \textbf{Elemento} & \textbf{Símbolo} & \textbf{Estados de oxidación}         \\
%            \hline
%            \multicolumn{3}{|c|}{\textbf{METALES}} \\
%            \hline
%            Sodio             & Na               & +1                                    \\
%            Potasio           & K                & +1                                    \\
%            Litio             & Li               & +1                                    \\
%            Plata             & Ag               & +1                                    \\
%            \hline
%            Calcio            & Ca               & +2                                    \\
%            Magnesio          & Mg               & +2                                    \\
%            Bario             & Ba               & +2                                    \\
%            Zinc              & Zn               & +2                                    \\
%            \hline
%            Aluminio          & Al               & +3                                    \\
%            \hline
%            Hierro            & Fe               & +2, +3                                \\
%            Cobre             & Cu               & +1, +2                                \\
%            Plomo             & Pb               & +2, +4                                \\
%            Estaño            & Sn               & +2, +4                                \\
%            \hline
%            \multicolumn{3}{|c|}{\textbf{NO METALES}} \\
%            \hline
%            Oxígeno           & O                & -2 (casi siempre)                     \\
%            Hidrógeno         & H                & +1 (con no metales), -1 (con metales) \\
%            Flúor             & F                & -1                                    \\
%            Cloro             & Cl               & -1, +1, +3, +5, +7                    \\
%            Bromo             & Br               & -1, +1, +3, +5                        \\
%            Yodo              & I                & -1, +1, +3, +5, +7                    \\
%            Azufre            & S                & -2, +2, +4, +6                        \\
%            Nitrógeno         & N                & -3, +1, +2, +3, +4, +5                \\
%            \hline
%        \end{tabular}
%    \end{center}

\newpage

\subsection{Tabla resumen de todos los compuestos}

%\begin{tcolorbox}[enhanced,
%    fonttitle=\bfseries,
%%fontupper=\normalsize\sffamily,% para poner el tipo de letra en el resto de la table
%    colback=white, colframe=blue2, colbacktitle=blue5, coltitle=white,
%    center title, clip upper,
%    title=Tabla resumen,
%    tabularray={
%        colspec={X|[0.25mm]X|[0.25mm]X|[0.25mm]X|[0.25mm]X|[0.25mm]X},
%        column{2-Z}={halign=c},
%        row{1}={font=\bfseries,bg=blue9},
%        row{1-Z}={halign=c},
%    }
%]
%    \textbf{Tipo} & \textbf{Fórmula general} & \textbf{Ejemplo} & \textbf{Nombre} \\
%    \hline
%    Óxido & \ch{E_xO_y} & \ch{Fe2O3} & Trióxido de dihierro \\
%    &  &  & Óxido de hierro (III) \\
%    \hline
%    Hidruro metálico & \ch{MH_n} & \ch{CaH2} & Dihidruro de calcio \\
%    &  &  & Hidruro de calcio \\
%    \hline
%    Hidruro no metálico & \ch{H_nX} & \ch{HCl} & Cloruro de hidrógeno \\
%    (compuesto puro) &  &  &  \\
%    \hline
%    Hidrácido & \ch{H_nX(ac)} & \ch{HCl(ac)} & Ácido clorhídrico \\
%    (en disolución) &  &  &  \\
%    \hline
%    Sal binaria & \ch{M_xX_y} & \ch{NaCl} & Cloruro de sodio \\
%    &  & \ch{FeCl3} & Tricloruro de hierro \\
%    &  &  & Cloruro de hierro (III) \\
%    \hline
%    Hidróxido & \ch{M(OH)_n} & \ch{Ca(OH)2} & Dihidróxido de calcio \\
%    &  &  & Hidróxido de calcio \\
%    \hline
%    Peróxido & \ch{M2O2} o \ch{MO2} & \ch{Na2O2} & Peróxido de sodio \\
%    &  & \ch{CaO2} & Peróxido de calcio \\
%    \hline
%\end{tcolorbox}

\begin{tcolorbox}[enhanced,
    fonttitle=\bfseries,
%fontupper=\normalsize\sffamily,% para poner el tipo de letra en el resto de la table
    colback=white, colframe=ocre, colbacktitle=ocre, coltitle=white,
    center title, clip upper,
    title=Tabla resumen,
    tabularray={
        colspec={X|[0.25mm]X|[0.25mm]X|[0.25mm]X|[0.25mm]X|[0.25mm]X},
        column{2-Z}={halign=c},
        row{1}={font=\bfseries,bg=ocre!10},
        row{1-Z}={halign=c},
    }
]
    \textbf{Tipo} & \textbf{Fórmula general} & \textbf{Ejemplo} & \textbf{Nombre} \\
    \hline
    Óxido & \ch{E_xO_y} & \ch{Fe2O3} & Trióxido de dihierro \\
    &  &  & Óxido de hierro (III) \\
    \hline
    Hidruro metálico & \ch{MH_n} & \ch{CaH2} & Dihidruro de calcio \\
    &  &  & Hidruro de calcio \\
    \hline
    Hidruro no metálico & \ch{H_nX} & \ch{HCl} & Cloruro de hidrógeno \\
    (compuesto puro) &  &  &  \\
    \hline
    Hidrácido & \ch{H_nX(ac)} & \ch{HCl(ac)} & Ácido clorhídrico \\
    (en disolución) &  &  &  \\
    \hline
    Sal binaria & \ch{M_xX_y} & \ch{NaCl} & Cloruro de sodio \\
    &  & \ch{FeCl3} & Tricloruro de hierro \\
    &  &  & Cloruro de hierro (III) \\
    \hline
    Hidróxido & \ch{M(OH)_n} & \ch{Ca(OH)2} & Dihidróxido de calcio \\
    &  &  & Hidróxido de calcio \\
    \hline
    Peróxido & \ch{M2O2} o \ch{MO2} & \ch{Na2O2} & Peróxido de sodio \\
    &  & \ch{CaO2} & Peróxido de calcio \\
    \hline
\end{tcolorbox}

\textbf{Leyenda}:
\begin{itemize}
    \item E = cualquier elemento
    \item M = metal
    \item X = no metal
    \item n = número que depende del estado de oxidación
\end{itemize}

\newpage

\pagestyle{empty}


\newgeometry{
    top=0.75cm,
    bottom=0.5cm,
    left=0.5cm,
    right=0.5cm
}


\DefTblrTemplate{contfoot-text}{normal}{Continúa en la siguiente página}
    \SetTblrTemplate{contfoot-text}{normal}
    \DefTblrTemplate{conthead-text}{normal}{(Continuación)}
    \SetTblrTemplate{conthead-text}{normal}

    \setcounter{numeroformula}{0}

    \begin{landscape}

        % --- tabla antigua ---
        % \begin{longtblr}[caption=Compuestos binarios]{hlines,vlines, colspec={Q[c,m] Q[c,m, 2cm] Q[c,m,12cm] Q[c,m,12cm] }, rowsep=5pt, rowhead=1}
            % \setcounter{numeroformula}{0}
            % \SetRow{gray!50} nº & Fórmula         & Multiplicadores            & Estado de Oxidación            \\
            % \SetCell[c=4]{c} \textbf{Óxidos} \\
            % \myorder{}          & \ch{Na2O}       &    \rojo{}                       &             \rojo{}                            \\
            % \myorder{}          & \ch{Co2O3}      &           \rojo{}                 &   \rojo{}                                      \\
            % \myorder{}          & \ch{CuO}        &            \rojo{}                &   \rojo{}                                     \\
            % \myorder{}          & \ch{Ag2O}       &       \rojo{}                     &          \rojo{}                               \\
            % \myorder{}          & \ch{CdO}        &              \rojo{}              &       \rojo{}                                  \\
            % \myorder{}          & \ch{NO}         &           \rojo{}                 &               \rojo{}                        \\
            % \myorder{}          & \ch{CO2}        &              \rojo{}              &      \rojo{}                                   \\
            % \myorder{}          & \ch{As2O5}      &            \rojo{}                &         \rojo{}                                \\
            % \myorder{}          & \ch{SO3}        &               \rojo{}             &        \rojo{}                                 \\
            % \myorder{}          & \ch{O5I2}       &               \rojo{}             &     \rojo{}                                    \\
            % \myorder{}          & \ch{O7Br2}      &                    \rojo{}        &        \rojo{}                                 \\
            % \myorder{}          &      \rojo{\ch{}}          & Óxido de dilitio           &         \rojo{}                                \\
            % \myorder{}          &   \rojo{\ch{}}              &        \rojo{}                    & Óxido de magnesio                       \\
            % \myorder{}          &   \rojo{\ch{}}              &        \rojo{}                    & Óxido de potasio                        \\
            % \myorder{}          &   \rojo{\ch{}}              &         \rojo{}                   & Óxido de estroncio                      \\
            % \myorder{}          &   \rojo{\ch{}}              &         \rojo{}                   & óxido de hierro (II)                    \\
            % \myorder{}          &    \rojo{\ch{}}             & Dióxido de plomo           &         \rojo{}                                \\
            % \myorder{}          &    \rojo{\ch{}}             &       \rojo{}                     & óxido de oro (III)                      \\
            % \myorder{}          &    \rojo{\ch{}}             & trióxido de dialuminio     &           \rojo{}                             \\
            % \myorder{}          &    \rojo{\ch{}}             &       \rojo{}                     & óxido de estaño (II)                    \\
            % \myorder{}          &    \rojo{\ch{}}             &       \rojo{}                     & óxido de cinc                           \\
            % \myorder{}          &    \rojo{\ch{}}             &          \rojo{}                  & óxido de mercurio (II)                  \\
            % \myorder{}          &    \rojo{\ch{}}             & tetraóxido de dinitrógeno  &       \rojo{}                                  \\
            % \myorder{}          &    \rojo{\ch{}}             &               \rojo{}             & óxido de cobalto (III)                  \\
            % \myorder{}          &     \rojo{\ch{}}            & Pentóxido de dinitrógeno   &          \rojo{}                               \\
            % \myorder{}          &     \rojo{\ch{}}            &                 \rojo{}          & óxido de bario                          \\
            % \myorder{}          &     \rojo{\ch{}}            & trióxido de azufre         &      \rojo{}                                   \\
            % \myorder{}          &     \rojo{\ch{}}            & trióxido de dicromo        &       \rojo{}                                  \\
            % \myorder{}          &     \rojo{\ch{}}            &                 \rojo{}          & óxido de manganeso (IV)                 \\
            % \myorder{}          &     \rojo{\ch{}}            & monóxido de dinitrógeno    &        \rojo{}                                 \\
            % \myorder{}          &     \rojo{\ch{}}            &                   \rojo{}        & óxido de platino (IV)                   \\
            % \myorder{}          &     \rojo{\ch{}}            & dibromuro de trioxígeno    &       \rojo{}                                 \\
            % \SetCell[c=4]{c} \textbf{Hidruros} \\
            % \myorder{}          & \ch{NaH}        &              \rojo{}              &      \rojo{}                                   \\
            % \myorder{}          & \ch{CaH2}       &                   \rojo{}         &     \rojo{}                                    \\
            % \myorder{}          & \ch{CoH2}       &                          \rojo{}  &      \rojo{}                                   \\
            % \myorder{}          & \ch{CuH2}       &                 \rojo{}           &  \rojo{}                                       \\
            % \myorder{}          & \ch{AgH}        &                    \rojo{}        &      \rojo{}                                   \\
            % \myorder{}          & \ch{AlH3}       &                  \rojo{}          &                \rojo{}                         \\
            % \myorder{}          & \ch{NH3}        &                   \rojo{}         &                                         \\
            % \myorder{}          & \ch{CdH2}       &                   \rojo{}         &               \rojo{}                          \\
            % \myorder{}          & \ch{AuH3}       &                  \rojo{}          &                \rojo{}                         \\
            % \myorder{}          & \ch{MnH3}       &                   \rojo{}         &              \rojo{}                           \\
            % \myorder{}          & \ch{PtH2}       &                   \rojo{}         &               \rojo{}                          \\
            % \myorder{}          &       \rojo{\ch{}}          & dihidruro de níquel        &                  \rojo{}                       \\
            % \myorder{}          &      \rojo{\ch{}}          &                    \rojo{}        & hidruro de níquel  (III)                \\
            % \myorder{}          &        \rojo{\ch{}}         &                   \rojo{}         & hidruro de sodio                        \\
            % \myorder{}          &       \rojo{\ch{}}          &                 \rojo{}           & hidruro de calcio                       \\
            % \myorder{}          &          \rojo{\ch{}}       & dihidruro de magnesio      &             \rojo{}                            \\
            % \myorder{}          &        \rojo{\ch{}}         &                  \rojo{}          & hidruro de mercurio (II)                \\
            % \myorder{}          &            \rojo{\ch{}}     & dihidruro de hierro        &             \rojo{}                            \\
            % \myorder{}          &             \rojo{\ch{}}    &                  \rojo{}          & hidruro de cinc                         \\
            % \myorder{}          &           \rojo{\ch{}}      & dihidruro de cobre         &            \rojo{}                             \\
            % \myorder{}          &            \rojo{\ch{}}     &                 \rojo{}           & hidruro de plomo (IV)                   \\
            % \myorder{}          &            \rojo{\ch{}}     &                  \rojo{}          & hidruro de estaño (II)                  \\
            % \myorder{}          &           \rojo{\ch{}}     &                  \rojo{}           & hidruro de níquel (III)                 \\
            % \myorder{}          &            \rojo{\ch{}}     & trihidruro de cromo        &              \rojo{}                           \\
            % \myorder{}          &           \rojo{\ch{}}      &                  \rojo{}          & hidruro de cromo (II)                   \\
            % \myorder{}          & \ch{PH3}        &                                   &                                         \\
            % \myorder{}          & \ch{AsH3}       &                                   &                                         \\
            % \myorder{}          & \ch{SbH3}       &                                   &                                         \\
            % \myorder{}          & \ch{BH3}        &                                   &                                         \\
            % \myorder{}          & \ch{SiH4}       &                                   &                                         \\
            % \myorder{}          & \ch{CH4}        &                                   &                                         \\
            % \myorder{}          & \ch{H2O}        &                                   &                                         \\
            % \myorder{}          & \ch{H2S}        &                                   &                                         \\
            % \myorder{}          & \ch{H2Se}       &                                   &                                         \\
            % \myorder{}          & \ch{H2Te}       &                                   &                                         \\
            % \myorder{}          & \ch{HF}         &                                   &                                         \\
            % \myorder{}          & \ch{HCl}        &                                   &                                         \\
            % \myorder{}          & \ch{HBr}        &                                   &                                         \\
            % \myorder{}          & \ch{HI}         &                                   &                                         \\
            % \SetCell[c=4]{c} \textbf{Sales binarias} \\
            % \myorder{}          & \ch{NaCl}       &                            &             \rojo{}                            \\
            % \myorder{}          & \ch{MgI2}       &                            &                \rojo{}                         \\
            % \myorder{}          & \ch{BeBr2}      &                            &              \rojo{}                           \\
            % \myorder{}          & \ch{K2S}        &                            &             \rojo{}                            \\
            % \myorder{}          & \ch{FeF2}       &                            &             \rojo{}                            \\
            % \myorder{}          & \ch{PbCl4}      &                            &             \rojo{}                            \\
            % \myorder{}          & \ch{Au2S3}      &                            &             \rojo{}                            \\
            % \myorder{}          & \ch{SnS}        &                            &             \rojo{}                            \\
            % \myorder{}          & \ch{ZnBr2}      &                            &             \rojo{}                            \\
            % \myorder{}          & \ch{HgTe}       &                            &             \rojo{}                            \\
            % \myorder{}          & \ch{CoI2}       &                            &             \rojo{}                            \\
            % \myorder{}          & \ch{CaCl2}      &                            &             \rojo{}                            \\
            % \myorder{}          & \ch{Al2S3}      &                            &             \rojo{}                            \\
            % \myorder{}          & \ch{Cr2S3}      &                            &             \rojo{}                            \\
            % \myorder{}          & \ch{MnBr2}      &                            &             \rojo{}                            \\
            % \myorder{}          & \ch{Na2Te}      &                            &             \rojo{}                            \\
            % \myorder{}          & \ch{ZnCl2}      &                            &             \rojo{}                            \\
            % \myorder{}          & \ch{Ag2S}       &                            &             \rojo{}                            \\
            % \myorder{}          & \ch{CI4}        &                            &             \rojo{}                            \\
            % \myorder{}          & \ch{NF3}        &                            &               \rojo{}                          \\
            % \myorder{}          & \ch{IBr}        &                      \rojo{}      &                                         \\
            % \myorder{}          & \ch{SeBr4}      &                            &               \rojo{}                          \\
            % \myorder{}          & \ch{SeCl6}      &                            &               \rojo{}                          \\
            % \myorder{}          & \ch{AsF5}       &                            &                                         \\
            % \myorder{}          & \ch{CCl4}       &                            &                \rojo{}                         \\
            % \myorder{}          & \ch{SiF4}       &                            &                \rojo{}                         \\
            % \myorder{}          &       \rojo{\ch{}}          &                            & cloruro de litio                        \\
            % \myorder{}          &      \rojo{\ch{}}          & diyoduro de magnesio       &              \rojo{}                          \\
            % \myorder{}          &        \rojo{\ch{}}         &                 \rojo{}           & bromuro de berilio                      \\
            % \myorder{}          &       \rojo{\ch{}}          & sulfuro de dipotasio       &             \rojo{}                            \\
            % \myorder{}          &        \rojo{\ch{}}         &                \rojo{}            & fluroro de hierro (II)                  \\
            % \myorder{}          &         \rojo{\ch{}}        & tetracloruro de plomo      &              \rojo{}                           \\
            % \myorder{}          &         \rojo{\ch{}}        &       \rojo{}                     & sulfuro de oro (III)                    \\
            % \myorder{}          &         \rojo{\ch{}}        &      \rojo{}                      & sulfuro de estaño (II)                  \\
            % \myorder{}          &          \rojo{\ch{}}       & dibromuro de cinc          &              \rojo{}                           \\
            % \myorder{}          &           \rojo{\ch{}}      &     \rojo{}                     & teluro de mercurio (II)                 \\
            % \myorder{}          &           \rojo{\ch{}}      &      \rojo{}                & yoduro de cobalto (II)                  \\
            % \myorder{}          &            \rojo{\ch{}}     &  \rojo{}                       & cloruro de calcio                       \\
            % \myorder{}          &           \rojo{\ch{}}      & trisulfuro de dialuminio   &              \rojo{}                           \\
            % \myorder{}          &           \rojo{\ch{}}      & trisulfuro de dicromo      &              \rojo{}                           \\
            % \myorder{}          &          \rojo{\ch{}}       &   \rojo{}                         & bromuro de manganeso (II)               \\
            % \myorder{}          &            \rojo{\ch{}}     & teluro de disodio          &               \rojo{}                          \\
            % \myorder{}          &           \rojo{\ch{}}     &     \rojo{}                      & cloruro de cinc                         \\
            % \myorder{}          &            \rojo{\ch{}}     & sulfuro de diplata         &               \rojo{}                          \\
            % \myorder{}          &            \rojo{\ch{}}     & trifluoruro de bromo       &               \rojo{}                          \\
            % \myorder{}          &           \rojo{\ch{}}      &     \rojo{}                      & cloruro de bromo (I)                    \\
            % \myorder{}          &          \rojo{\ch{}}      & hexafluoruro de azufre     &                \rojo{}                         \\
            % \myorder{}          &          \rojo{\ch{}}       &   \rojo{}                       & sulfuro de carbono (IV)                 \\
            % \SetCell[c=4]{c} \textbf{Hidróxidos} \\
            % \myorder{}          & \ch{CsOH}       &              \rojo{}              &               \rojo{}                          \\
            % \myorder{}          & \ch{KOH}        &               \rojo{}             &             \rojo{}                            \\
            % \myorder{}          & \ch{Be(OH)2}    &                 \rojo{}           &               \rojo{}                          \\
            % \myorder{}          & \ch{Fe(OH)3}    &                \rojo{}            &               \rojo{}                          \\
            % \myorder{}          & \ch{AgOH}       &               \rojo{}             &               \rojo{}                          \\
            % \myorder{}       & \ch{Al(OH)3}    &              \rojo{}              &                 \rojo{}                        \\
            % \myorder{}       & \ch{NH4OH}      &               \rojo{}             &                  \rojo{}                       \\
            % \myorder{}       & \ch{Cd(OH)2}    &               \rojo{}             &                \rojo{}                         \\
            % \myorder{}       & \ch{Au(OH)3}    &               \rojo{}             &                 \rojo{}                       \\
            % \myorder{}       & \ch{Mn(OH)3}    &               \rojo{}             &                 \rojo{}                        \\
            % \myorder{}       & \ch{Pt(OH)2}    &               \rojo{}             &                  \rojo{}                       \\
            % \myorder{}       &        \rojo{\ch{}}         & dihidróxido de níquel  &                  \rojo{}                       \\
            % \myorder{}       &       \rojo{\ch{}}          &     \rojo{}                      & hidróxido de níquel  (III)                   \\
            % \myorder{}       &         \rojo{\ch{}}        &       \rojo{}                    & hidróxido de sodio                      \\
            % \myorder{}       &        \rojo{\ch{}}         &     \rojo{}                        & hidróxido de calcio                     \\
            % \myorder{}       &          \rojo{\ch{}}       & dihidróxido de magnesio    &                   \rojo{}                      \\
            % \myorder{}       &        \rojo{\ch{}}         &     \rojo{}                     & hidróxido de mercurio (II)              \\
            % \myorder{}       &            \rojo{\ch{}}     & dihidróxido de hierro      &                 \rojo{}                        \\
            % \myorder{}       &             \rojo{\ch{}}    &     \rojo{}                   & hidróxido de cinc                       \\
            % \myorder{}       &           \rojo{\ch{}}      & dihidróxido de cobre       &              \rojo{}                           \\
            % \myorder{}       &            \rojo{\ch{}}     &    \rojo{}                       & hidróxido de plomo (IV)                 \\
            % \myorder{}       &            \rojo{\ch{}}     &   \rojo{}                         & hidróxido de estaño (II)                \\
            % \myorder{}       &           \rojo{\ch{}}     &    \rojo{}                        & hidróxido de níquel (III)               \\
            % \myorder{}       &            \rojo{\ch{}}     & trihidróxido de cromo      &              \rojo{}                           \\
            % \myorder{}       &           \rojo{\ch{}}      &   \rojo{}                         & hidróxido de cromo (II)                 \\
            % \SetCell[c=4]{c} { \textbf{Peróxidos}. \\ En este caso solo hemos dado la nomenclatura de los peróxidos.       \\ Podéis poner su nombre en la casilla de multiplicadores dejando el resto en blanco}       \\
            % \myorder{}       & \ch{Na2O2}       &         \rojo{}                   & --                                      \\
            % \myorder{}       & \ch{K2O2}       &     \rojo{}                       & --                                      \\
            % \myorder{}       & \ch{SrO2}      &      \rojo{}                      & --                                      \\
            % \myorder{}       & \ch{CaO2}       &      \rojo{}                      & --                                     \\
            % \myorder{}       & \ch{Rb2O2}       &     \rojo{}                       & --                                    \\
            % \myorder{} & \ch{CuO2} & \rojo{} & --  \\
            % \myorder{}           & \ch{MgO2} & \rojo{} & -- \\
            % \myorder{}       & \ch{BaO2} & \rojo{} & --                            \\
            % \SetCell[c=4]{c} {\textbf{Hidruros grupos 16 y 17 en disolución acuosa}. Solo hemos dado una nomenclatura.       \\ Podéis poner su nombre en la casilla de multiplicadores dejando el resto en blanco }       \\
            % \myorder{}       & \ch{HCl_{(aq)}}    & \rojo{}            &                \\
            % \myorder{}       & \ch{HBr_\acu}   & \rojo{}               &             \\
            % \myorder{}       & \ch{HI_\acu}  & \rojo{}                &            \\
            % \myorder{}       & \ch{HF_\acu}  & \rojo{}                &              \\
            % \myorder{}       & \ch{H2S_\acu}  & \rojo{}                &                \\
            % \myorder{}       & \ch{H2Te_\acu}  & \rojo{}                &             \\
            % \myorder{}       & \ch{H2Se_\acu}  & \rojo{}                &              \\
            % \hline
        % \end{longtblr}
        % --- tabla nueva ---
        % Ejercicio con las soluciones de la tabla: cada libro carga las suyas
        % (\LibroENG se define en Libro_3_ESO_ENG.tex)
        \begin{exercise}
            \ifdefined\LibroENG
                Complete Table~\ref{tab:compuestos_binarios}: write the formula or the missing names (multipliers and oxidation state) of each compound.
            \else
                Completa la Tabla~\ref{tab:compuestos_binarios}: escribe la fórmula o los nombres que faltan (multiplicadores y estado de oxidación) de cada compuesto.
            \fi
            \begin{solu}
                \ifdefined\LibroENG% Soluciones de la tabla «Compuestos binarios» del capítulo de formulación (inorganica_apuntes_3_ESO_claude_Libro.tex).
% Se carga desde el fichero de soluciones (ejercicio de la tabla, entorno solu) en el libro en inglés.
% En negro, los datos del enunciado; en rojo, las respuestas. Óxidos, hidróxidos y peróxidos coinciden con la tabla «Binary compounds».

\providecommand{\cabeceraSolFormEN}[1]{%
    \toprule
    \textbf{No.} & \textbf{Formula} & \textbf{Multipliers} & \textbf{Oxidation state} \\
    \midrule
    \multicolumn{4}{c}{\textbf{#1}} \\
    \midrule
}

{\small
\begin{center}
    \begin{tabular}{@{}r C{1.9cm} L{5cm} L{5cm}@{}}
        \cabeceraSolFormEN{Oxides}
        1 & \ch{Na2O} & \rojo{Disodium oxide} & \rojo{Sodium oxide} \\[2pt]
        2 & \ch{Co2O3} & \rojo{Dicobalt trioxide} & \rojo{Cobalt (III) oxide} \\[2pt]
        3 & \ch{CuO} & \rojo{Copper oxide (copper monoxide)} & \rojo{Copper (II) oxide} \\[2pt]
        4 & \ch{Ag2O} & \rojo{Disilver oxide} & \rojo{Silver oxide} \\[2pt]
        5 & \ch{CdO} & \rojo{Cadmium oxide} & \rojo{Cadmium oxide} \\[2pt]
        6 & \ch{NO} & \rojo{Nitrogen monoxide} & \rojo{Nitrogen (II) oxide} \\[2pt]
        7 & \ch{CO2} & \rojo{Carbon dioxide} & \rojo{Carbon (IV) oxide} \\[2pt]
        8 & \ch{As2O5} & \rojo{Diarsenic pentoxide} & \rojo{Arsenic (V) oxide} \\[2pt]
        9 & \ch{SO3} & \rojo{Sulfur trioxide} & \rojo{Sulfur (VI) oxide} \\[2pt]
        10 & \ch{O5I2} & \rojo{Pentaoxygen diiodide} & -- \\[2pt]
        11 & \ch{O7Br2} & \rojo{Heptaoxygen dibromide} & -- \\[2pt]
        12 & \rojo{\ch{Li2O}} & Dilithium oxide & \rojo{Lithium oxide} \\[2pt]
        13 & \rojo{\ch{MgO}} & \rojo{Magnesium oxide} & Magnesium oxide \\[2pt]
        14 & \rojo{\ch{K2O}} & \rojo{Dipotassium oxide} & Potassium oxide \\[2pt]
        15 & \rojo{\ch{SrO}} & \rojo{Strontium oxide} & Strontium oxide \\[2pt]
        16 & \rojo{\ch{FeO}} & \rojo{Iron oxide (iron monoxide)} & Iron (II) oxide \\[2pt]
        17 & \rojo{\ch{PbO2}} & Lead dioxide & \rojo{Lead (IV) oxide} \\[2pt]
        18 & \rojo{\ch{Au2O3}} & \rojo{Digold trioxide} & Gold (III) oxide \\[2pt]
        19 & \rojo{\ch{Al2O3}} & Dialuminum trioxide & \rojo{Aluminum oxide} \\[2pt]
        20 & \rojo{\ch{SnO}} & \rojo{Tin oxide (tin monoxide)} & Tin (II) oxide \\[2pt]
        21 & \rojo{\ch{ZnO}} & \rojo{Zinc oxide} & Zinc oxide \\[2pt]
        22 & \rojo{\ch{HgO}} & \rojo{Mercury oxide (mercury monoxide)} & Mercury (II) oxide \\[2pt]
        23 & \rojo{\ch{N2O4}} & Dinitrogen tetraoxide & \rojo{Nitrogen (IV) oxide} \\[2pt]
        24 & \rojo{\ch{Co2O3}} & \rojo{Dicobalt trioxide} & Cobalt (III) oxide \\[2pt]
        25 & \rojo{\ch{N2O5}} & Dinitrogen pentoxide & \rojo{Nitrogen (V) oxide} \\[2pt]
        26 & \rojo{\ch{BaO}} & \rojo{Barium oxide} & Barium oxide \\[2pt]
        27 & \rojo{\ch{SO3}} & Sulfur trioxide & \rojo{Sulfur (VI) oxide} \\[2pt]
        28 & \rojo{\ch{Cr2O3}} & Dichromium trioxide & \rojo{Chromium (III) oxide} \\[2pt]
        29 & \rojo{\ch{MnO2}} & \rojo{Manganese dioxide} & Manganese (IV) oxide \\[2pt]
        30 & \rojo{\ch{N2O}} & Dinitrogen monoxide & \rojo{Nitrogen (I) oxide} \\[2pt]
        31 & \rojo{\ch{PtO2}} & \rojo{Platinum dioxide} & Platinum (IV) oxide \\[2pt]
        32 & \rojo{\ch{O3Br2}} & Trioxygen dibromide & -- \\
        \bottomrule
    \end{tabular}
\end{center}

\noindent In 10, 11 and 32 the halogen is written on the right (it is more electronegative than oxygen), so they are named only with multipliers.

\begin{center}
    \begin{tabular}{@{}r C{1.9cm} L{5cm} L{5cm}@{}}
        \cabeceraSolFormEN{Hydrides}
        33 & \ch{NaH} & \rojo{Sodium hydride} & \rojo{Sodium hydride} \\[2pt]
        34 & \ch{CaH2} & \rojo{Calcium dihydride} & \rojo{Calcium hydride} \\[2pt]
        35 & \ch{CoH2} & \rojo{Cobalt dihydride} & \rojo{Cobalt (II) hydride} \\[2pt]
        36 & \ch{CuH2} & \rojo{Copper dihydride} & \rojo{Copper (II) hydride} \\[2pt]
        37 & \ch{AgH} & \rojo{Silver hydride} & \rojo{Silver hydride} \\[2pt]
        38 & \ch{AlH3} & \rojo{Aluminum trihydride} & \rojo{Aluminum hydride} \\[2pt]
        39 & \ch{NH3} & \rojo{Ammonia (nitrogen trihydride)} & -- \\[2pt]
        40 & \ch{CdH2} & \rojo{Cadmium dihydride} & \rojo{Cadmium hydride} \\[2pt]
        41 & \ch{AuH3} & \rojo{Gold trihydride} & \rojo{Gold (III) hydride} \\[2pt]
        42 & \ch{MnH3} & \rojo{Manganese trihydride} & \rojo{Manganese (III) hydride} \\[2pt]
        43 & \ch{PtH2} & \rojo{Platinum dihydride} & \rojo{Platinum (II) hydride} \\[2pt]
        44 & \rojo{\ch{NiH2}} & Nickel dihydride & \rojo{Nickel (II) hydride} \\[2pt]
        45 & \rojo{\ch{NiH3}} & \rojo{Nickel trihydride} & Nickel (III) hydride \\[2pt]
        46 & \rojo{\ch{NaH}} & \rojo{Sodium hydride} & Sodium hydride \\[2pt]
        47 & \rojo{\ch{CaH2}} & \rojo{Calcium dihydride} & Calcium hydride \\[2pt]
        48 & \rojo{\ch{MgH2}} & Magnesium dihydride & \rojo{Magnesium hydride} \\[2pt]
        49 & \rojo{\ch{HgH2}} & \rojo{Mercury dihydride} & Mercury (II) hydride \\[2pt]
        50 & \rojo{\ch{FeH2}} & Iron dihydride & \rojo{Iron (II) hydride} \\[2pt]
        51 & \rojo{\ch{ZnH2}} & \rojo{Zinc dihydride} & Zinc hydride \\[2pt]
        52 & \rojo{\ch{CuH2}} & Copper dihydride & \rojo{Copper (II) hydride} \\[2pt]
        53 & \rojo{\ch{PbH4}} & \rojo{Lead tetrahydride} & Lead (IV) hydride \\[2pt]
        54 & \rojo{\ch{SnH2}} & \rojo{Tin dihydride} & Tin (II) hydride \\[2pt]
        55 & \rojo{\ch{NiH3}} & \rojo{Nickel trihydride} & Nickel (III) hydride \\[2pt]
        56 & \rojo{\ch{CrH3}} & Chromium trihydride & \rojo{Chromium (III) hydride} \\[2pt]
        57 & \rojo{\ch{CrH2}} & \rojo{Chromium dihydride} & Chromium (II) hydride \\
        \midrule
        \multicolumn{4}{c}{\textbf{Hydrides of groups 13, 14 and 15 and water (names accepted by IUPAC)}} \\
        \midrule
        58 & \ch{PH3} & \rojo{Phosphane (phosphorus trihydride)} & -- \\[2pt]
        59 & \ch{AsH3} & \rojo{Arsane (arsenic trihydride)} & -- \\[2pt]
        60 & \ch{SbH3} & \rojo{Stibane (antimony trihydride)} & -- \\[2pt]
        61 & \ch{BH3} & \rojo{Borane (boron trihydride)} & -- \\[2pt]
        62 & \ch{SiH4} & \rojo{Silane (silicon tetrahydride)} & -- \\[2pt]
        63 & \ch{CH4} & \rojo{Methane (carbon tetrahydride)} & -- \\[2pt]
        64 & \ch{H2O} & \rojo{Water (oxidane)} & -- \\
        \bottomrule
    \end{tabular}
\end{center}

\begin{center}
    \begin{tabular}{@{}r C{1.9cm} L{5cm} L{5cm}@{}}
        \cabeceraSolFormEN{Hydrogen with a nonmetal of groups 16 and 17}
        65 & \ch{H2S} & \rojo{Dihydrogen sulfide} & \rojo{Hydrogen sulfide} \\[2pt]
        66 & \ch{H2Se} & \rojo{Dihydrogen selenide} & \rojo{Hydrogen selenide} \\[2pt]
        67 & \ch{H2Te} & \rojo{Dihydrogen telluride} & \rojo{Hydrogen telluride} \\[2pt]
        68 & \ch{HF} & \rojo{Hydrogen fluoride} & \rojo{Hydrogen fluoride} \\[2pt]
        69 & \ch{HCl} & \rojo{Hydrogen chloride} & \rojo{Hydrogen chloride} \\[2pt]
        70 & \ch{HBr} & \rojo{Hydrogen bromide} & \rojo{Hydrogen bromide} \\[2pt]
        71 & \ch{HI} & \rojo{Hydrogen iodide} & \rojo{Hydrogen iodide} \\
        \bottomrule
    \end{tabular}
\end{center}

\noindent In 58-71 the table leaves the cells blank: the names are given as extra information. The hydrides of groups 13, 14 and 15 (and water) have their own names accepted by IUPAC. Those of groups 16 and 17 are named «hydrogen …-ide»; since hydrogen only has the +1 state, the second name is the same without prefixes.

\begin{center}
    \begin{tabular}{@{}r C{1.9cm} L{5cm} L{5cm}@{}}
        \cabeceraSolFormEN{Binary salts}
        72 & \ch{NaCl} & \rojo{Sodium chloride} & \rojo{Sodium chloride} \\[2pt]
        73 & \ch{MgI2} & \rojo{Magnesium diiodide} & \rojo{Magnesium iodide} \\[2pt]
        74 & \ch{BeBr2} & \rojo{Beryllium dibromide} & \rojo{Beryllium bromide} \\[2pt]
        75 & \ch{K2S} & \rojo{Dipotassium sulfide} & \rojo{Potassium sulfide} \\[2pt]
        76 & \ch{FeF2} & \rojo{Iron difluoride} & \rojo{Iron (II) fluoride} \\[2pt]
        77 & \ch{PbCl4} & \rojo{Lead tetrachloride} & \rojo{Lead (IV) chloride} \\[2pt]
        78 & \ch{Au2S3} & \rojo{Digold trisulfide} & \rojo{Gold (III) sulfide} \\[2pt]
        79 & \ch{SnS} & \rojo{Tin sulfide (tin monosulfide)} & \rojo{Tin (II) sulfide} \\[2pt]
        80 & \ch{ZnBr2} & \rojo{Zinc dibromide} & \rojo{Zinc bromide} \\[2pt]
        81 & \ch{HgTe} & \rojo{Mercury telluride} & \rojo{Mercury (II) telluride} \\[2pt]
        82 & \ch{CoI2} & \rojo{Cobalt diiodide} & \rojo{Cobalt (II) iodide} \\[2pt]
        83 & \ch{CaCl2} & \rojo{Calcium dichloride} & \rojo{Calcium chloride} \\[2pt]
        84 & \ch{Al2S3} & \rojo{Dialuminum trisulfide} & \rojo{Aluminum sulfide} \\[2pt]
        85 & \ch{Cr2S3} & \rojo{Dichromium trisulfide} & \rojo{Chromium (III) sulfide} \\[2pt]
        86 & \ch{MnBr2} & \rojo{Manganese dibromide} & \rojo{Manganese (II) bromide} \\[2pt]
        87 & \ch{Na2Te} & \rojo{Disodium telluride} & \rojo{Sodium telluride} \\[2pt]
        88 & \ch{ZnCl2} & \rojo{Zinc dichloride} & \rojo{Zinc chloride} \\[2pt]
        89 & \ch{Ag2S} & \rojo{Disilver sulfide} & \rojo{Silver sulfide} \\[2pt]
        90 & \ch{CI4} & \rojo{Carbon tetraiodide} & \rojo{Carbon (IV) iodide} \\[2pt]
        91 & \ch{NF3} & \rojo{Nitrogen trifluoride} & \rojo{Nitrogen (III) fluoride} \\[2pt]
        92 & \ch{IBr} & \rojo{Iodine bromide} & \rojo{Iodine (I) bromide} \\[2pt]
        93 & \ch{SeBr4} & \rojo{Selenium tetrabromide} & \rojo{Selenium (IV) bromide} \\[2pt]
        94 & \ch{SeCl6} & \rojo{Selenium hexachloride} & \rojo{Selenium (VI) chloride} \\[2pt]
        95 & \ch{AsF5} & \rojo{Arsenic pentafluoride} & \rojo{Arsenic (V) fluoride} \\[2pt]
        96 & \ch{CCl4} & \rojo{Carbon tetrachloride} & \rojo{Carbon (IV) chloride} \\[2pt]
        97 & \ch{SiF4} & \rojo{Silicon tetrafluoride} & \rojo{Silicon fluoride} \\
        \bottomrule
    \end{tabular}
\end{center}

\begin{center}
    \begin{tabular}{@{}r C{1.9cm} L{5cm} L{5cm}@{}}
        \cabeceraSolFormEN{Binary salts (continued)}
        98 & \rojo{\ch{LiCl}} & \rojo{Lithium chloride} & Lithium chloride \\[2pt]
        99 & \rojo{\ch{MgI2}} & Magnesium diiodide & \rojo{Magnesium iodide} \\[2pt]
        100 & \rojo{\ch{BeBr2}} & \rojo{Beryllium dibromide} & Beryllium bromide \\[2pt]
        101 & \rojo{\ch{K2S}} & Dipotassium sulfide & \rojo{Potassium sulfide} \\[2pt]
        102 & \rojo{\ch{FeF2}} & \rojo{Iron difluoride} & Iron (II) fluoride \\[2pt]
        103 & \rojo{\ch{PbCl4}} & Lead tetrachloride & \rojo{Lead (IV) chloride} \\[2pt]
        104 & \rojo{\ch{Au2S3}} & \rojo{Digold trisulfide} & Gold (III) sulfide \\[2pt]
        105 & \rojo{\ch{SnS}} & \rojo{Tin sulfide (tin monosulfide)} & Tin (II) sulfide \\[2pt]
        106 & \rojo{\ch{ZnBr2}} & Zinc dibromide & \rojo{Zinc bromide} \\[2pt]
        107 & \rojo{\ch{HgTe}} & \rojo{Mercury telluride} & Mercury (II) telluride \\[2pt]
        108 & \rojo{\ch{CoI2}} & \rojo{Cobalt diiodide} & Cobalt (II) iodide \\[2pt]
        109 & \rojo{\ch{CaCl2}} & \rojo{Calcium dichloride} & Calcium chloride \\[2pt]
        110 & \rojo{\ch{Al2S3}} & Dialuminum trisulfide & \rojo{Aluminum sulfide} \\[2pt]
        111 & \rojo{\ch{Cr2S3}} & Dichromium trisulfide & \rojo{Chromium (III) sulfide} \\[2pt]
        112 & \rojo{\ch{MnBr2}} & \rojo{Manganese dibromide} & Manganese (II) bromide \\[2pt]
        113 & \rojo{\ch{Na2Te}} & Disodium telluride & \rojo{Sodium telluride} \\[2pt]
        114 & \rojo{\ch{ZnCl2}} & \rojo{Zinc dichloride} & Zinc chloride \\[2pt]
        115 & \rojo{\ch{Ag2S}} & Disilver sulfide & \rojo{Silver sulfide} \\[2pt]
        116 & \rojo{\ch{BrF3}} & Bromine trifluoride & \rojo{Bromine (III) fluoride} \\[2pt]
        117 & \rojo{\ch{BrCl}} & \rojo{Bromine chloride} & Bromine (I) chloride \\[2pt]
        118 & \rojo{\ch{SF6}} & Sulfur hexafluoride & \rojo{Sulfur (VI) fluoride} \\[2pt]
        119 & \rojo{\ch{CS2}} & \rojo{Carbon disulfide} & Carbon (IV) sulfide \\
        \bottomrule
    \end{tabular}
\end{center}

\begin{center}
    \begin{tabular}{@{}r C{1.9cm} L{5cm} L{5cm}@{}}
        \cabeceraSolFormEN{Hydroxides}
        120 & \ch{CsOH} & \rojo{Cesium hydroxide} & \rojo{Cesium hydroxide} \\[2pt]
        121 & \ch{KOH} & \rojo{Potassium hydroxide} & \rojo{Potassium hydroxide} \\[2pt]
        122 & \ch{Be(OH)2} & \rojo{Beryllium dihydroxide} & \rojo{Beryllium hydroxide} \\[2pt]
        123 & \ch{Fe(OH)3} & \rojo{Iron trihydroxide} & \rojo{Iron (III) hydroxide} \\[2pt]
        124 & \ch{AgOH} & \rojo{Silver hydroxide} & \rojo{Silver hydroxide} \\[2pt]
        125 & \ch{Al(OH)3} & \rojo{Aluminum trihydroxide} & \rojo{Aluminum hydroxide} \\[2pt]
        126 & \ch{NH4OH} & \rojo{Ammonium hydroxide} & \rojo{Ammonium hydroxide} \\[2pt]
        127 & \ch{Cd(OH)2} & \rojo{Cadmium dihydroxide} & \rojo{Cadmium hydroxide} \\[2pt]
        128 & \ch{Au(OH)3} & \rojo{Gold trihydroxide} & \rojo{Gold (III) hydroxide} \\[2pt]
        129 & \ch{Mn(OH)3} & \rojo{Manganese trihydroxide} & \rojo{Manganese (III) hydroxide} \\[2pt]
        130 & \ch{Pt(OH)2} & \rojo{Platinum dihydroxide} & \rojo{Platinum (II) hydroxide} \\[2pt]
        131 & \rojo{\ch{Ni(OH)2}} & Nickel dihydroxide & \rojo{Nickel (II) hydroxide} \\[2pt]
        132 & \rojo{\ch{Ni(OH)3}} & \rojo{Nickel trihydroxide} & Nickel (III) hydroxide \\[2pt]
        133 & \rojo{\ch{NaOH}} & \rojo{Sodium hydroxide} & Sodium hydroxide \\[2pt]
        134 & \rojo{\ch{Ca(OH)2}} & \rojo{Calcium dihydroxide} & Calcium hydroxide \\[2pt]
        135 & \rojo{\ch{Mg(OH)2}} & Magnesium dihydroxide & \rojo{Magnesium hydroxide} \\[2pt]
        136 & \rojo{\ch{Hg(OH)2}} & \rojo{Mercury dihydroxide} & Mercury (II) hydroxide \\[2pt]
        137 & \rojo{\ch{Fe(OH)2}} & Iron dihydroxide & \rojo{Iron (II) hydroxide} \\[2pt]
        138 & \rojo{\ch{Zn(OH)2}} & \rojo{Zinc dihydroxide} & Zinc hydroxide \\[2pt]
        139 & \rojo{\ch{Cu(OH)2}} & Copper dihydroxide & \rojo{Copper (II) hydroxide} \\[2pt]
        140 & \rojo{\ch{Pb(OH)4}} & \rojo{Lead tetrahydroxide} & Lead (IV) hydroxide \\[2pt]
        141 & \rojo{\ch{Sn(OH)2}} & \rojo{Tin dihydroxide} & Tin (II) hydroxide \\[2pt]
        142 & \rojo{\ch{Ni(OH)3}} & \rojo{Nickel trihydroxide} & Nickel (III) hydroxide \\[2pt]
        143 & \rojo{\ch{Cr(OH)3}} & Chromium trihydroxide & \rojo{Chromium (III) hydroxide} \\[2pt]
        144 & \rojo{\ch{Cr(OH)2}} & \rojo{Chromium dihydroxide} & Chromium (II) hydroxide \\
        \bottomrule
    \end{tabular}
\end{center}

\begin{center}
    \begin{tabular}{@{}r C{1.9cm} L{5cm} L{5cm}@{}}
        \cabeceraSolFormEN{Peroxides}
        145 & \ch{Na2O2} & \rojo{Sodium peroxide} & -- \\[2pt]
        146 & \ch{K2O2} & \rojo{Potassium peroxide} & -- \\[2pt]
        147 & \ch{SrO2} & \rojo{Strontium peroxide} & -- \\[2pt]
        148 & \ch{CaO2} & \rojo{Calcium peroxide} & -- \\[2pt]
        149 & \ch{Rb2O2} & \rojo{Rubidium peroxide} & -- \\[2pt]
        150 & \ch{CuO2} & \rojo{Copper (II) peroxide} & -- \\[2pt]
        151 & \ch{MgO2} & \rojo{Magnesium peroxide} & -- \\[2pt]
        152 & \ch{BaO2} & \rojo{Barium peroxide} & -- \\
        \midrule
        \multicolumn{4}{c}{\textbf{Hydrides from groups 16 and 17 in aqueous solution}} \\
        \midrule
        153 & \ch{HCl(aq)} & \rojo{Hydrochloric acid} & -- \\[2pt]
        154 & \ch{HBr(aq)} & \rojo{Hydrobromic acid} & -- \\[2pt]
        155 & \ch{HI(aq)} & \rojo{Hydroiodic acid} & -- \\[2pt]
        156 & \ch{HF(aq)} & \rojo{Hydrofluoric acid} & -- \\[2pt]
        157 & \ch{H2S(aq)} & \rojo{Hydrosulfuric acid} & -- \\[2pt]
        158 & \ch{H2Te(aq)} & \rojo{Hydrotelluric acid} & -- \\[2pt]
        159 & \ch{H2Se(aq)} & \rojo{Hydroselenic acid} & -- \\
        \bottomrule
    \end{tabular}
\end{center}
}
\else% Soluciones de la tabla «Compuestos binarios» del capítulo de formulación (inorganica_apuntes_3_ESO_claude_Libro.tex).
% Se carga desde el fichero de soluciones (ejercicio de la tabla, entorno solu) en el libro en español.
% En negro, los datos del enunciado; en rojo, las respuestas. Óxidos, hidróxidos y peróxidos coinciden con la tabla «Binary compounds».

\providecommand{\cabeceraSolFormES}[1]{%
    \toprule
    \textbf{Nº} & \textbf{Fórmula} & \textbf{Multiplicadores} & \textbf{Estado de oxidación} \\
    \midrule
    \multicolumn{4}{c}{\textbf{#1}} \\
    \midrule
}

{\small
\begin{center}
    \begin{tabular}{@{}r C{1.9cm} L{5cm} L{5cm}@{}}
        \cabeceraSolFormES{Óxidos}
        1 & \ch{Na2O} & \rojo{Óxido de disodio} & \rojo{Óxido de sodio} \\[2pt]
        2 & \ch{Co2O3} & \rojo{Trióxido de dicobalto} & \rojo{Óxido de cobalto (III)} \\[2pt]
        3 & \ch{CuO} & \rojo{Óxido de cobre (monóxido de cobre)} & \rojo{Óxido de cobre (II)} \\[2pt]
        4 & \ch{Ag2O} & \rojo{Óxido de diplata} & \rojo{Óxido de plata} \\[2pt]
        5 & \ch{CdO} & \rojo{Óxido de cadmio} & \rojo{Óxido de cadmio} \\[2pt]
        6 & \ch{NO} & \rojo{Monóxido de nitrógeno} & \rojo{Óxido de nitrógeno (II)} \\[2pt]
        7 & \ch{CO2} & \rojo{Dióxido de carbono} & \rojo{Óxido de carbono (IV)} \\[2pt]
        8 & \ch{As2O5} & \rojo{Pentaóxido de diarsénico} & \rojo{Óxido de arsénico (V)} \\[2pt]
        9 & \ch{SO3} & \rojo{Trióxido de azufre} & \rojo{Óxido de azufre (VI)} \\[2pt]
        10 & \ch{O5I2} & \rojo{Diyoduro de pentaoxígeno} & -- \\[2pt]
        11 & \ch{O7Br2} & \rojo{Dibromuro de heptaoxígeno} & -- \\[2pt]
        12 & \rojo{\ch{Li2O}} & Óxido de dilitio & \rojo{Óxido de litio} \\[2pt]
        13 & \rojo{\ch{MgO}} & \rojo{Óxido de magnesio} & Óxido de magnesio \\[2pt]
        14 & \rojo{\ch{K2O}} & \rojo{Óxido de dipotasio} & Óxido de potasio \\[2pt]
        15 & \rojo{\ch{SrO}} & \rojo{Óxido de estroncio} & Óxido de estroncio \\[2pt]
        16 & \rojo{\ch{FeO}} & \rojo{Óxido de hierro (monóxido de hierro)} & Óxido de hierro (II) \\[2pt]
        17 & \rojo{\ch{PbO2}} & Dióxido de plomo & \rojo{Óxido de plomo (IV)} \\[2pt]
        18 & \rojo{\ch{Au2O3}} & \rojo{Trióxido de dioro} & Óxido de oro (III) \\[2pt]
        19 & \rojo{\ch{Al2O3}} & Trióxido de dialuminio & \rojo{Óxido de aluminio} \\[2pt]
        20 & \rojo{\ch{SnO}} & \rojo{Óxido de estaño (monóxido de estaño)} & Óxido de estaño (II) \\[2pt]
        21 & \rojo{\ch{ZnO}} & \rojo{Óxido de cinc} & Óxido de cinc \\[2pt]
        22 & \rojo{\ch{HgO}} & \rojo{Óxido de mercurio (monóxido de mercurio)} & Óxido de mercurio (II) \\[2pt]
        23 & \rojo{\ch{N2O4}} & Tetraóxido de dinitrógeno & \rojo{Óxido de nitrógeno (IV)} \\[2pt]
        24 & \rojo{\ch{Co2O3}} & \rojo{Trióxido de dicobalto} & Óxido de cobalto (III) \\[2pt]
        25 & \rojo{\ch{N2O5}} & Pentaóxido de dinitrógeno & \rojo{Óxido de nitrógeno (V)} \\[2pt]
        26 & \rojo{\ch{BaO}} & \rojo{Óxido de bario} & Óxido de bario \\[2pt]
        27 & \rojo{\ch{SO3}} & Trióxido de azufre & \rojo{Óxido de azufre (VI)} \\[2pt]
        28 & \rojo{\ch{Cr2O3}} & Trióxido de dicromo & \rojo{Óxido de cromo (III)} \\[2pt]
        29 & \rojo{\ch{MnO2}} & \rojo{Dióxido de manganeso} & Óxido de manganeso (IV) \\[2pt]
        30 & \rojo{\ch{N2O}} & Monóxido de dinitrógeno & \rojo{Óxido de nitrógeno (I)} \\[2pt]
        31 & \rojo{\ch{PtO2}} & \rojo{Dióxido de platino} & Óxido de platino (IV) \\[2pt]
        32 & \rojo{\ch{O3Br2}} & Dibromuro de trioxígeno & -- \\
        \bottomrule
    \end{tabular}
\end{center}

\noindent En 10, 11 y 32 el halógeno va a la derecha (es más electronegativo que el oxígeno), así que solo se nombran con multiplicadores.

\begin{center}
    \begin{tabular}{@{}r C{1.9cm} L{5cm} L{5cm}@{}}
        \cabeceraSolFormES{Hidruros}
        33 & \ch{NaH} & \rojo{Hidruro de sodio} & \rojo{Hidruro de sodio} \\[2pt]
        34 & \ch{CaH2} & \rojo{Dihidruro de calcio} & \rojo{Hidruro de calcio} \\[2pt]
        35 & \ch{CoH2} & \rojo{Dihidruro de cobalto} & \rojo{Hidruro de cobalto (II)} \\[2pt]
        36 & \ch{CuH2} & \rojo{Dihidruro de cobre} & \rojo{Hidruro de cobre (II)} \\[2pt]
        37 & \ch{AgH} & \rojo{Hidruro de plata} & \rojo{Hidruro de plata} \\[2pt]
        38 & \ch{AlH3} & \rojo{Trihidruro de aluminio} & \rojo{Hidruro de aluminio} \\[2pt]
        39 & \ch{NH3} & \rojo{Amoniaco (trihidruro de nitrógeno)} & -- \\[2pt]
        40 & \ch{CdH2} & \rojo{Dihidruro de cadmio} & \rojo{Hidruro de cadmio} \\[2pt]
        41 & \ch{AuH3} & \rojo{Trihidruro de oro} & \rojo{Hidruro de oro (III)} \\[2pt]
        42 & \ch{MnH3} & \rojo{Trihidruro de manganeso} & \rojo{Hidruro de manganeso (III)} \\[2pt]
        43 & \ch{PtH2} & \rojo{Dihidruro de platino} & \rojo{Hidruro de platino (II)} \\[2pt]
        44 & \rojo{\ch{NiH2}} & Dihidruro de níquel & \rojo{Hidruro de níquel (II)} \\[2pt]
        45 & \rojo{\ch{NiH3}} & \rojo{Trihidruro de níquel} & Hidruro de níquel (III) \\[2pt]
        46 & \rojo{\ch{NaH}} & \rojo{Hidruro de sodio} & Hidruro de sodio \\[2pt]
        47 & \rojo{\ch{CaH2}} & \rojo{Dihidruro de calcio} & Hidruro de calcio \\[2pt]
        48 & \rojo{\ch{MgH2}} & Dihidruro de magnesio & \rojo{Hidruro de magnesio} \\[2pt]
        49 & \rojo{\ch{HgH2}} & \rojo{Dihidruro de mercurio} & Hidruro de mercurio (II) \\[2pt]
        50 & \rojo{\ch{FeH2}} & Dihidruro de hierro & \rojo{Hidruro de hierro (II)} \\[2pt]
        51 & \rojo{\ch{ZnH2}} & \rojo{Dihidruro de cinc} & Hidruro de cinc \\[2pt]
        52 & \rojo{\ch{CuH2}} & Dihidruro de cobre & \rojo{Hidruro de cobre (II)} \\[2pt]
        53 & \rojo{\ch{PbH4}} & \rojo{Tetrahidruro de plomo} & Hidruro de plomo (IV) \\[2pt]
        54 & \rojo{\ch{SnH2}} & \rojo{Dihidruro de estaño} & Hidruro de estaño (II) \\[2pt]
        55 & \rojo{\ch{NiH3}} & \rojo{Trihidruro de níquel} & Hidruro de níquel (III) \\[2pt]
        56 & \rojo{\ch{CrH3}} & Trihidruro de cromo & \rojo{Hidruro de cromo (III)} \\[2pt]
        57 & \rojo{\ch{CrH2}} & \rojo{Dihidruro de cromo} & Hidruro de cromo (II) \\
        \midrule
        \multicolumn{4}{c}{\textbf{Hidruros de los grupos 13, 14 y 15 y agua (nombres aceptados por la IUPAC)}} \\
        \midrule
        58 & \ch{PH3} & \rojo{Fosfano (trihidruro de fósforo)} & -- \\[2pt]
        59 & \ch{AsH3} & \rojo{Arsano (trihidruro de arsénico)} & -- \\[2pt]
        60 & \ch{SbH3} & \rojo{Estibano (trihidruro de antimonio)} & -- \\[2pt]
        61 & \ch{BH3} & \rojo{Borano (trihidruro de boro)} & -- \\[2pt]
        62 & \ch{SiH4} & \rojo{Silano (tetrahidruro de silicio)} & -- \\[2pt]
        63 & \ch{CH4} & \rojo{Metano (tetrahidruro de carbono)} & -- \\[2pt]
        64 & \ch{H2O} & \rojo{Agua (oxidano)} & -- \\
        \bottomrule
    \end{tabular}
\end{center}

\begin{center}
    \begin{tabular}{@{}r C{1.9cm} L{5cm} L{5cm}@{}}
        \cabeceraSolFormES{Hidruros no metálicos (hidrógeno con un no metal de los grupos 16 y 17)}
        65 & \ch{H2S} & \rojo{Sulfuro de dihidrógeno} & \rojo{Sulfuro de hidrógeno} \\[2pt]
        66 & \ch{H2Se} & \rojo{Seleniuro de dihidrógeno} & \rojo{Seleniuro de hidrógeno} \\[2pt]
        67 & \ch{H2Te} & \rojo{Telururo de dihidrógeno} & \rojo{Telururo de hidrógeno} \\[2pt]
        68 & \ch{HF} & \rojo{Fluoruro de hidrógeno} & \rojo{Fluoruro de hidrógeno} \\[2pt]
        69 & \ch{HCl} & \rojo{Cloruro de hidrógeno} & \rojo{Cloruro de hidrógeno} \\[2pt]
        70 & \ch{HBr} & \rojo{Bromuro de hidrógeno} & \rojo{Bromuro de hidrógeno} \\[2pt]
        71 & \ch{HI} & \rojo{Yoduro de hidrógeno} & \rojo{Yoduro de hidrógeno} \\
        \bottomrule
    \end{tabular}
\end{center}

\noindent En 58-71 el enunciado deja las casillas en blanco: los nombres se dan como ampliación. Los hidruros de los grupos 13, 14 y 15 (y el agua) tienen nombres propios aceptados por la IUPAC. Los de los grupos 16 y 17 se nombran como «-uro de hidrógeno»; como el hidrógeno solo tiene el estado +1, el segundo nombre es el mismo sin prefijos.

\begin{center}
    \begin{tabular}{@{}r C{1.9cm} L{5cm} L{5cm}@{}}
        \cabeceraSolFormES{Sales binarias}
        72 & \ch{NaCl} & \rojo{Cloruro de sodio} & \rojo{Cloruro de sodio} \\[2pt]
        73 & \ch{MgI2} & \rojo{Diyoduro de magnesio} & \rojo{Yoduro de magnesio} \\[2pt]
        74 & \ch{BeBr2} & \rojo{Dibromuro de berilio} & \rojo{Bromuro de berilio} \\[2pt]
        75 & \ch{K2S} & \rojo{Sulfuro de dipotasio} & \rojo{Sulfuro de potasio} \\[2pt]
        76 & \ch{FeF2} & \rojo{Difluoruro de hierro} & \rojo{Fluoruro de hierro (II)} \\[2pt]
        77 & \ch{PbCl4} & \rojo{Tetracloruro de plomo} & \rojo{Cloruro de plomo (IV)} \\[2pt]
        78 & \ch{Au2S3} & \rojo{Trisulfuro de dioro} & \rojo{Sulfuro de oro (III)} \\[2pt]
        79 & \ch{SnS} & \rojo{Sulfuro de estaño (monosulfuro de estaño)} & \rojo{Sulfuro de estaño (II)} \\[2pt]
        80 & \ch{ZnBr2} & \rojo{Dibromuro de cinc} & \rojo{Bromuro de cinc} \\[2pt]
        81 & \ch{HgTe} & \rojo{Telururo de mercurio} & \rojo{Telururo de mercurio (II)} \\[2pt]
        82 & \ch{CoI2} & \rojo{Diyoduro de cobalto} & \rojo{Yoduro de cobalto (II)} \\[2pt]
        83 & \ch{CaCl2} & \rojo{Dicloruro de calcio} & \rojo{Cloruro de calcio} \\[2pt]
        84 & \ch{Al2S3} & \rojo{Trisulfuro de dialuminio} & \rojo{Sulfuro de aluminio} \\[2pt]
        85 & \ch{Cr2S3} & \rojo{Trisulfuro de dicromo} & \rojo{Sulfuro de cromo (III)} \\[2pt]
        86 & \ch{MnBr2} & \rojo{Dibromuro de manganeso} & \rojo{Bromuro de manganeso (II)} \\[2pt]
        87 & \ch{Na2Te} & \rojo{Telururo de disodio} & \rojo{Telururo de sodio} \\[2pt]
        88 & \ch{ZnCl2} & \rojo{Dicloruro de cinc} & \rojo{Cloruro de cinc} \\[2pt]
        89 & \ch{Ag2S} & \rojo{Sulfuro de diplata} & \rojo{Sulfuro de plata} \\[2pt]
        90 & \ch{CI4} & \rojo{Tetrayoduro de carbono} & \rojo{Yoduro de carbono (IV)} \\[2pt]
        91 & \ch{NF3} & \rojo{Trifluoruro de nitrógeno} & \rojo{Fluoruro de nitrógeno (III)} \\[2pt]
        92 & \ch{IBr} & \rojo{Bromuro de yodo} & \rojo{Bromuro de yodo (I)} \\[2pt]
        93 & \ch{SeBr4} & \rojo{Tetrabromuro de selenio} & \rojo{Bromuro de selenio (IV)} \\[2pt]
        94 & \ch{SeCl6} & \rojo{Hexacloruro de selenio} & \rojo{Cloruro de selenio (VI)} \\[2pt]
        95 & \ch{AsF5} & \rojo{Pentafluoruro de arsénico} & \rojo{Fluoruro de arsénico (V)} \\[2pt]
        96 & \ch{CCl4} & \rojo{Tetracloruro de carbono} & \rojo{Cloruro de carbono (IV)} \\[2pt]
        97 & \ch{SiF4} & \rojo{Tetrafluoruro de silicio} & \rojo{Fluoruro de silicio} \\
        \bottomrule
    \end{tabular}
\end{center}

\begin{center}
    \begin{tabular}{@{}r C{1.9cm} L{5cm} L{5cm}@{}}
        \cabeceraSolFormES{Sales binarias (continuación)}
        98 & \rojo{\ch{LiCl}} & \rojo{Cloruro de litio} & Cloruro de litio \\[2pt]
        99 & \rojo{\ch{MgI2}} & Diyoduro de magnesio & \rojo{Yoduro de magnesio} \\[2pt]
        100 & \rojo{\ch{BeBr2}} & \rojo{Dibromuro de berilio} & Bromuro de berilio \\[2pt]
        101 & \rojo{\ch{K2S}} & Sulfuro de dipotasio & \rojo{Sulfuro de potasio} \\[2pt]
        102 & \rojo{\ch{FeF2}} & \rojo{Difluoruro de hierro} & Fluoruro de hierro (II) \\[2pt]
        103 & \rojo{\ch{PbCl4}} & Tetracloruro de plomo & \rojo{Cloruro de plomo (IV)} \\[2pt]
        104 & \rojo{\ch{Au2S3}} & \rojo{Trisulfuro de dioro} & Sulfuro de oro (III) \\[2pt]
        105 & \rojo{\ch{SnS}} & \rojo{Sulfuro de estaño (monosulfuro de estaño)} & Sulfuro de estaño (II) \\[2pt]
        106 & \rojo{\ch{ZnBr2}} & Dibromuro de cinc & \rojo{Bromuro de cinc} \\[2pt]
        107 & \rojo{\ch{HgTe}} & \rojo{Telururo de mercurio} & Telururo de mercurio (II) \\[2pt]
        108 & \rojo{\ch{CoI2}} & \rojo{Diyoduro de cobalto} & Yoduro de cobalto (II) \\[2pt]
        109 & \rojo{\ch{CaCl2}} & \rojo{Dicloruro de calcio} & Cloruro de calcio \\[2pt]
        110 & \rojo{\ch{Al2S3}} & Trisulfuro de dialuminio & \rojo{Sulfuro de aluminio} \\[2pt]
        111 & \rojo{\ch{Cr2S3}} & Trisulfuro de dicromo & \rojo{Sulfuro de cromo (III)} \\[2pt]
        112 & \rojo{\ch{MnBr2}} & \rojo{Dibromuro de manganeso} & Bromuro de manganeso (II) \\[2pt]
        113 & \rojo{\ch{Na2Te}} & Telururo de disodio & \rojo{Telururo de sodio} \\[2pt]
        114 & \rojo{\ch{ZnCl2}} & \rojo{Dicloruro de cinc} & Cloruro de cinc \\[2pt]
        115 & \rojo{\ch{Ag2S}} & Sulfuro de diplata & \rojo{Sulfuro de plata} \\[2pt]
        116 & \rojo{\ch{BrF3}} & Trifluoruro de bromo & \rojo{Fluoruro de bromo (III)} \\[2pt]
        117 & \rojo{\ch{BrCl}} & \rojo{Cloruro de bromo} & Cloruro de bromo (I) \\[2pt]
        118 & \rojo{\ch{SF6}} & Hexafluoruro de azufre & \rojo{Fluoruro de azufre (VI)} \\[2pt]
        119 & \rojo{\ch{CS2}} & \rojo{Disulfuro de carbono} & Sulfuro de carbono (IV) \\
        \bottomrule
    \end{tabular}
\end{center}

\begin{center}
    \begin{tabular}{@{}r C{1.9cm} L{5cm} L{5cm}@{}}
        \cabeceraSolFormES{Hidróxidos}
        120 & \ch{CsOH} & \rojo{Hidróxido de cesio} & \rojo{Hidróxido de cesio} \\[2pt]
        121 & \ch{KOH} & \rojo{Hidróxido de potasio} & \rojo{Hidróxido de potasio} \\[2pt]
        122 & \ch{Be(OH)2} & \rojo{Dihidróxido de berilio} & \rojo{Hidróxido de berilio} \\[2pt]
        123 & \ch{Fe(OH)3} & \rojo{Trihidróxido de hierro} & \rojo{Hidróxido de hierro (III)} \\[2pt]
        124 & \ch{AgOH} & \rojo{Hidróxido de plata} & \rojo{Hidróxido de plata} \\[2pt]
        125 & \ch{Al(OH)3} & \rojo{Trihidróxido de aluminio} & \rojo{Hidróxido de aluminio} \\[2pt]
        126 & \ch{NH4OH} & \rojo{Hidróxido de amonio} & \rojo{Hidróxido de amonio} \\[2pt]
        127 & \ch{Cd(OH)2} & \rojo{Dihidróxido de cadmio} & \rojo{Hidróxido de cadmio} \\[2pt]
        128 & \ch{Au(OH)3} & \rojo{Trihidróxido de oro} & \rojo{Hidróxido de oro (III)} \\[2pt]
        129 & \ch{Mn(OH)3} & \rojo{Trihidróxido de manganeso} & \rojo{Hidróxido de manganeso (III)} \\[2pt]
        130 & \ch{Pt(OH)2} & \rojo{Dihidróxido de platino} & \rojo{Hidróxido de platino (II)} \\[2pt]
        131 & \rojo{\ch{Ni(OH)2}} & Dihidróxido de níquel & \rojo{Hidróxido de níquel (II)} \\[2pt]
        132 & \rojo{\ch{Ni(OH)3}} & \rojo{Trihidróxido de níquel} & Hidróxido de níquel (III) \\[2pt]
        133 & \rojo{\ch{NaOH}} & \rojo{Hidróxido de sodio} & Hidróxido de sodio \\[2pt]
        134 & \rojo{\ch{Ca(OH)2}} & \rojo{Dihidróxido de calcio} & Hidróxido de calcio \\[2pt]
        135 & \rojo{\ch{Mg(OH)2}} & Dihidróxido de magnesio & \rojo{Hidróxido de magnesio} \\[2pt]
        136 & \rojo{\ch{Hg(OH)2}} & \rojo{Dihidróxido de mercurio} & Hidróxido de mercurio (II) \\[2pt]
        137 & \rojo{\ch{Fe(OH)2}} & Dihidróxido de hierro & \rojo{Hidróxido de hierro (II)} \\[2pt]
        138 & \rojo{\ch{Zn(OH)2}} & \rojo{Dihidróxido de cinc} & Hidróxido de cinc \\[2pt]
        139 & \rojo{\ch{Cu(OH)2}} & Dihidróxido de cobre & \rojo{Hidróxido de cobre (II)} \\[2pt]
        140 & \rojo{\ch{Pb(OH)4}} & \rojo{Tetrahidróxido de plomo} & Hidróxido de plomo (IV) \\[2pt]
        141 & \rojo{\ch{Sn(OH)2}} & \rojo{Dihidróxido de estaño} & Hidróxido de estaño (II) \\[2pt]
        142 & \rojo{\ch{Ni(OH)3}} & \rojo{Trihidróxido de níquel} & Hidróxido de níquel (III) \\[2pt]
        143 & \rojo{\ch{Cr(OH)3}} & Trihidróxido de cromo & \rojo{Hidróxido de cromo (III)} \\[2pt]
        144 & \rojo{\ch{Cr(OH)2}} & \rojo{Dihidróxido de cromo} & Hidróxido de cromo (II) \\
        \bottomrule
    \end{tabular}
\end{center}

\begin{center}
    \begin{tabular}{@{}r C{1.9cm} L{5cm} L{5cm}@{}}
        \cabeceraSolFormES{Peróxidos}
        145 & \ch{Na2O2} & \rojo{Peróxido de sodio} & -- \\[2pt]
        146 & \ch{K2O2} & \rojo{Peróxido de potasio} & -- \\[2pt]
        147 & \ch{SrO2} & \rojo{Peróxido de estroncio} & -- \\[2pt]
        148 & \ch{CaO2} & \rojo{Peróxido de calcio} & -- \\[2pt]
        149 & \ch{Rb2O2} & \rojo{Peróxido de rubidio} & -- \\[2pt]
        150 & \ch{CuO2} & \rojo{Peróxido de cobre (II)} & -- \\[2pt]
        151 & \ch{MgO2} & \rojo{Peróxido de magnesio} & -- \\[2pt]
        152 & \ch{BaO2} & \rojo{Peróxido de bario} & -- \\
        \midrule
        \multicolumn{4}{c}{\textbf{Hidruros de los grupos 16 y 17 en disolución acuosa}} \\
        \midrule
        153 & \ch{HCl(ac)} & \rojo{Ácido clorhídrico} & -- \\[2pt]
        154 & \ch{HBr(ac)} & \rojo{Ácido bromhídrico} & -- \\[2pt]
        155 & \ch{HI(ac)} & \rojo{Ácido yodhídrico} & -- \\[2pt]
        156 & \ch{HF(ac)} & \rojo{Ácido fluorhídrico} & -- \\[2pt]
        157 & \ch{H2S(ac)} & \rojo{Ácido sulfhídrico} & -- \\[2pt]
        158 & \ch{H2Te(ac)} & \rojo{Ácido telurhídrico} & -- \\[2pt]
        159 & \ch{H2Se(ac)} & \rojo{Ácido selenhídrico} & -- \\
        \bottomrule
    \end{tabular}
\end{center}
}
\fi
            \end{solu}
        \end{exercise}
        \setcounter{numeroformula}{0}
        \begin{longtable}{@{}c C{2cm} C{12cm} C{12cm}@{}}
            \caption{Compuestos binarios}\label{tab:compuestos_binarios} \\
            \toprule
            \textbf{nº} & \textbf{Fórmula} & \textbf{Multiplicadores} & \textbf{Estado de Oxidación} \\
            \midrule
            \endfirsthead
            \toprule
            \textbf{nº} & \textbf{Fórmula} & \textbf{Multiplicadores} & \textbf{Estado de Oxidación} \\
            \midrule
            \endhead
            \bottomrule
            \endlastfoot
            \multicolumn{4}{c}{\textbf{Óxidos}} \\[3pt]
            \myorder{} & \ch{Na2O} & \rojo{} & \rojo{} \\[3pt]
            \myorder{} & \ch{Co2O3} & \rojo{} & \rojo{} \\[3pt]
            \myorder{} & \ch{CuO} & \rojo{} & \rojo{} \\[3pt]
            \myorder{} & \ch{Ag2O} & \rojo{} & \rojo{} \\[3pt]
            \myorder{} & \ch{CdO} & \rojo{} & \rojo{} \\[3pt]
            \myorder{} & \ch{NO} & \rojo{} & \rojo{} \\[3pt]
            \myorder{} & \ch{CO2} & \rojo{} & \rojo{} \\[3pt]
            \myorder{} & \ch{As2O5} & \rojo{} & \rojo{} \\[3pt]
            \myorder{} & \ch{SO3} & \rojo{} & \rojo{} \\[3pt]
            \myorder{} & \ch{O5I2} & \rojo{} & \rojo{} \\[3pt]
            \myorder{} & \ch{O7Br2} & \rojo{} & \rojo{} \\[3pt]
            \myorder{} & \rojo{\ch{}} & Óxido de dilitio & \rojo{} \\[3pt]
            \myorder{} & \rojo{\ch{}} & \rojo{} & Óxido de magnesio \\[3pt]
            \myorder{} & \rojo{\ch{}} & \rojo{} & Óxido de potasio \\[3pt]
            \myorder{} & \rojo{\ch{}} & \rojo{} & Óxido de estroncio \\[3pt]
            \myorder{} & \rojo{\ch{}} & \rojo{} & óxido de hierro (II) \\[3pt]
            \myorder{} & \rojo{\ch{}} & Dióxido de plomo & \rojo{} \\[3pt]
            \myorder{} & \rojo{\ch{}} & \rojo{} & óxido de oro (III) \\[3pt]
            \myorder{} & \rojo{\ch{}} & trióxido de dialuminio & \rojo{} \\[3pt]
            \myorder{} & \rojo{\ch{}} & \rojo{} & óxido de estaño (II) \\[3pt]
            \myorder{} & \rojo{\ch{}} & \rojo{} & óxido de cinc \\[3pt]
            \myorder{} & \rojo{\ch{}} & \rojo{} & óxido de mercurio (II) \\[3pt]
            \myorder{} & \rojo{\ch{}} & tetraóxido de dinitrógeno & \rojo{} \\[3pt]
            \myorder{} & \rojo{\ch{}} & \rojo{} & óxido de cobalto (III) \\[3pt]
            \myorder{} & \rojo{\ch{}} & Pentóxido de dinitrógeno & \rojo{} \\[3pt]
            \myorder{} & \rojo{\ch{}} & \rojo{} & óxido de bario \\[3pt]
            \myorder{} & \rojo{\ch{}} & trióxido de azufre & \rojo{} \\[3pt]
            \myorder{} & \rojo{\ch{}} & trióxido de dicromo & \rojo{} \\[3pt]
            \myorder{} & \rojo{\ch{}} & \rojo{} & óxido de manganeso (IV) \\[3pt]
            \myorder{} & \rojo{\ch{}} & monóxido de dinitrógeno & \rojo{} \\[3pt]
            \myorder{} & \rojo{\ch{}} & \rojo{} & óxido de platino (IV) \\[3pt]
            \myorder{} & \rojo{\ch{}} & dibromuro de trioxígeno & \rojo{} \\[3pt]
            \multicolumn{4}{c}{\textbf{Hidruros}} \\[3pt]
            \myorder{} & \ch{NaH} & \rojo{} & \rojo{} \\[3pt]
            \myorder{} & \ch{CaH2} & \rojo{} & \rojo{} \\[3pt]
            \myorder{} & \ch{CoH2} & \rojo{} & \rojo{} \\[3pt]
            \myorder{} & \ch{CuH2} & \rojo{} & \rojo{} \\[3pt]
            \myorder{} & \ch{AgH} & \rojo{} & \rojo{} \\[3pt]
            \myorder{} & \ch{AlH3} & \rojo{} & \rojo{} \\[3pt]
            \myorder{} & \ch{NH3} & \rojo{} &  \\[3pt]
            \myorder{} & \ch{CdH2} & \rojo{} & \rojo{} \\[3pt]
            \myorder{} & \ch{AuH3} & \rojo{} & \rojo{} \\[3pt]
            \myorder{} & \ch{MnH3} & \rojo{} & \rojo{} \\[3pt]
            \myorder{} & \ch{PtH2} & \rojo{} & \rojo{} \\[3pt]
            \myorder{} & \rojo{\ch{}} & dihidruro de níquel & \rojo{} \\[3pt]
            \myorder{} & \rojo{\ch{}} & \rojo{} & hidruro de níquel  (III) \\[3pt]
            \myorder{} & \rojo{\ch{}} & \rojo{} & hidruro de sodio \\[3pt]
            \myorder{} & \rojo{\ch{}} & \rojo{} & hidruro de calcio \\[3pt]
            \myorder{} & \rojo{\ch{}} & dihidruro de magnesio & \rojo{} \\[3pt]
            \myorder{} & \rojo{\ch{}} & \rojo{} & hidruro de mercurio (II) \\[3pt]
            \myorder{} & \rojo{\ch{}} & dihidruro de hierro & \rojo{} \\[3pt]
            \myorder{} & \rojo{\ch{}} & \rojo{} & hidruro de cinc \\[3pt]
            \myorder{} & \rojo{\ch{}} & dihidruro de cobre & \rojo{} \\[3pt]
            \myorder{} & \rojo{\ch{}} & \rojo{} & hidruro de plomo (IV) \\[3pt]
            \myorder{} & \rojo{\ch{}} & \rojo{} & hidruro de estaño (II) \\[3pt]
            \myorder{} & \rojo{\ch{}} & \rojo{} & hidruro de níquel (III) \\[3pt]
            \myorder{} & \rojo{\ch{}} & trihidruro de cromo & \rojo{} \\[3pt]
            \myorder{} & \rojo{\ch{}} & \rojo{} & hidruro de cromo (II) \\[3pt]
            \myorder{} & \ch{PH3} &  &  \\[3pt]
            \myorder{} & \ch{AsH3} &  &  \\[3pt]
            \myorder{} & \ch{SbH3} &  &  \\[3pt]
            \myorder{} & \ch{BH3} &  &  \\[3pt]
            \myorder{} & \ch{SiH4} &  &  \\[3pt]
            \myorder{} & \ch{CH4} &  &  \\[3pt]
            \myorder{} & \ch{H2O} &  &  \\[3pt]
            \myorder{} & \ch{H2S} &  &  \\[3pt]
            \myorder{} & \ch{H2Se} &  &  \\[3pt]
            \myorder{} & \ch{H2Te} &  &  \\[3pt]
            \myorder{} & \ch{HF} &  &  \\[3pt]
            \myorder{} & \ch{HCl} &  &  \\[3pt]
            \myorder{} & \ch{HBr} &  &  \\[3pt]
            \myorder{} & \ch{HI} &  &  \\[3pt]
            \multicolumn{4}{c}{\textbf{Sales binarias}} \\[3pt]
            \myorder{} & \ch{NaCl} &  & \rojo{} \\[3pt]
            \myorder{} & \ch{MgI2} &  & \rojo{} \\[3pt]
            \myorder{} & \ch{BeBr2} &  & \rojo{} \\[3pt]
            \myorder{} & \ch{K2S} &  & \rojo{} \\[3pt]
            \myorder{} & \ch{FeF2} &  & \rojo{} \\[3pt]
            \myorder{} & \ch{PbCl4} &  & \rojo{} \\[3pt]
            \myorder{} & \ch{Au2S3} &  & \rojo{} \\[3pt]
            \myorder{} & \ch{SnS} &  & \rojo{} \\[3pt]
            \myorder{} & \ch{ZnBr2} &  & \rojo{} \\[3pt]
            \myorder{} & \ch{HgTe} &  & \rojo{} \\[3pt]
            \myorder{} & \ch{CoI2} &  & \rojo{} \\[3pt]
            \myorder{} & \ch{CaCl2} &  & \rojo{} \\[3pt]
            \myorder{} & \ch{Al2S3} &  & \rojo{} \\[3pt]
            \myorder{} & \ch{Cr2S3} &  & \rojo{} \\[3pt]
            \myorder{} & \ch{MnBr2} &  & \rojo{} \\[3pt]
            \myorder{} & \ch{Na2Te} &  & \rojo{} \\[3pt]
            \myorder{} & \ch{ZnCl2} &  & \rojo{} \\[3pt]
            \myorder{} & \ch{Ag2S} &  & \rojo{} \\[3pt]
            \myorder{} & \ch{CI4} &  & \rojo{} \\[3pt]
            \myorder{} & \ch{NF3} &  & \rojo{} \\[3pt]
            \myorder{} & \ch{IBr} & \rojo{} &  \\[3pt]
            \myorder{} & \ch{SeBr4} &  & \rojo{} \\[3pt]
            \myorder{} & \ch{SeCl6} &  & \rojo{} \\[3pt]
            \myorder{} & \ch{AsF5} &  &  \\[3pt]
            \myorder{} & \ch{CCl4} &  & \rojo{} \\[3pt]
            \myorder{} & \ch{SiF4} &  & \rojo{} \\[3pt]
            \myorder{} & \rojo{\ch{}} &  & cloruro de litio \\[3pt]
            \myorder{} & \rojo{\ch{}} & diyoduro de magnesio & \rojo{} \\[3pt]
            \myorder{} & \rojo{\ch{}} & \rojo{} & bromuro de berilio \\[3pt]
            \myorder{} & \rojo{\ch{}} & sulfuro de dipotasio & \rojo{} \\[3pt]
            \myorder{} & \rojo{\ch{}} & \rojo{} & fluoruro de hierro (II) \\[3pt]
            \myorder{} & \rojo{\ch{}} & tetracloruro de plomo & \rojo{} \\[3pt]
            \myorder{} & \rojo{\ch{}} & \rojo{} & sulfuro de oro (III) \\[3pt]
            \myorder{} & \rojo{\ch{}} & \rojo{} & sulfuro de estaño (II) \\[3pt]
            \myorder{} & \rojo{\ch{}} & dibromuro de cinc & \rojo{} \\[3pt]
            \myorder{} & \rojo{\ch{}} & \rojo{} & telururo de mercurio (II) \\[3pt]
            \myorder{} & \rojo{\ch{}} & \rojo{} & yoduro de cobalto (II) \\[3pt]
            \myorder{} & \rojo{\ch{}} & \rojo{} & cloruro de calcio \\[3pt]
            \myorder{} & \rojo{\ch{}} & trisulfuro de dialuminio & \rojo{} \\[3pt]
            \myorder{} & \rojo{\ch{}} & trisulfuro de dicromo & \rojo{} \\[3pt]
            \myorder{} & \rojo{\ch{}} & \rojo{} & bromuro de manganeso (II) \\[3pt]
            \myorder{} & \rojo{\ch{}} & telururo de disodio & \rojo{} \\[3pt]
            \myorder{} & \rojo{\ch{}} & \rojo{} & cloruro de cinc \\[3pt]
            \myorder{} & \rojo{\ch{}} & sulfuro de diplata & \rojo{} \\[3pt]
            \myorder{} & \rojo{\ch{}} & trifluoruro de bromo & \rojo{} \\[3pt]
            \myorder{} & \rojo{\ch{}} & \rojo{} & cloruro de bromo (I) \\[3pt]
            \myorder{} & \rojo{\ch{}} & hexafluoruro de azufre & \rojo{} \\[3pt]
            \myorder{} & \rojo{\ch{}} & \rojo{} & sulfuro de carbono (IV) \\[3pt]
            \multicolumn{4}{c}{\textbf{Hidróxidos}} \\[3pt]
            \myorder{} & \ch{CsOH} & \rojo{} & \rojo{} \\[3pt]
            \myorder{} & \ch{KOH} & \rojo{} & \rojo{} \\[3pt]
            \myorder{} & \ch{Be(OH)2} & \rojo{} & \rojo{} \\[3pt]
            \myorder{} & \ch{Fe(OH)3} & \rojo{} & \rojo{} \\[3pt]
            \myorder{} & \ch{AgOH} & \rojo{} & \rojo{} \\[3pt]
            \myorder{} & \ch{Al(OH)3} & \rojo{} & \rojo{} \\[3pt]
            \myorder{} & \ch{NH4OH} & \rojo{} & \rojo{} \\[3pt]
            \myorder{} & \ch{Cd(OH)2} & \rojo{} & \rojo{} \\[3pt]
            \myorder{} & \ch{Au(OH)3} & \rojo{} & \rojo{} \\[3pt]
            \myorder{} & \ch{Mn(OH)3} & \rojo{} & \rojo{} \\[3pt]
            \myorder{} & \ch{Pt(OH)2} & \rojo{} & \rojo{} \\[3pt]
            \myorder{} & \rojo{\ch{}} & dihidróxido de níquel & \rojo{} \\[3pt]
            \myorder{} & \rojo{\ch{}} & \rojo{} & hidróxido de níquel  (III) \\[3pt]
            \myorder{} & \rojo{\ch{}} & \rojo{} & hidróxido de sodio \\[3pt]
            \myorder{} & \rojo{\ch{}} & \rojo{} & hidróxido de calcio \\[3pt]
            \myorder{} & \rojo{\ch{}} & dihidróxido de magnesio & \rojo{} \\[3pt]
            \myorder{} & \rojo{\ch{}} & \rojo{} & hidróxido de mercurio (II) \\[3pt]
            \myorder{} & \rojo{\ch{}} & dihidróxido de hierro & \rojo{} \\[3pt]
            \myorder{} & \rojo{\ch{}} & \rojo{} & hidróxido de cinc \\[3pt]
            \myorder{} & \rojo{\ch{}} & dihidróxido de cobre & \rojo{} \\[3pt]
            \myorder{} & \rojo{\ch{}} & \rojo{} & hidróxido de plomo (IV) \\[3pt]
            \myorder{} & \rojo{\ch{}} & \rojo{} & hidróxido de estaño (II) \\[3pt]
            \myorder{} & \rojo{\ch{}} & \rojo{} & hidróxido de níquel (III) \\[3pt]
            \myorder{} & \rojo{\ch{}} & trihidróxido de cromo & \rojo{} \\[3pt]
            \myorder{} & \rojo{\ch{}} & \rojo{} & hidróxido de cromo (II) \\[3pt]
            \multicolumn{4}{c}{\makecell{\textbf{Peróxidos}. \\ En este caso solo hemos dado la nomenclatura de los peróxidos.       \\ Podéis poner su nombre en la casilla de multiplicadores dejando el resto en blanco}} \\[3pt]
            \myorder{} & \ch{Na2O2} & \rojo{} & -- \\[3pt]
            \myorder{} & \ch{K2O2} & \rojo{} & -- \\[3pt]
            \myorder{} & \ch{SrO2} & \rojo{} & -- \\[3pt]
            \myorder{} & \ch{CaO2} & \rojo{} & -- \\[3pt]
            \myorder{} & \ch{Rb2O2} & \rojo{} & -- \\[3pt]
            \myorder{} & \ch{CuO2} & \rojo{} & -- \\[3pt]
            \myorder{} & \ch{MgO2} & \rojo{} & -- \\[3pt]
            \myorder{} & \ch{BaO2} & \rojo{} & -- \\[3pt]
            \multicolumn{4}{c}{\makecell{\textbf{Hidruros grupos 16 y 17 en disolución acuosa}. Solo hemos dado una nomenclatura.       \\ Podéis poner su nombre en la casilla de multiplicadores dejando el resto en blanco}} \\[3pt]
            \myorder{} & \ch{HCl_{(aq)}} & \rojo{} &  \\[3pt]
            \myorder{} & \ch{HBr_\acu} & \rojo{} &  \\[3pt]
            \myorder{} & \ch{HI_\acu} & \rojo{} &  \\[3pt]
            \myorder{} & \ch{HF_\acu} & \rojo{} &  \\[3pt]
            \myorder{} & \ch{H2S_\acu} & \rojo{} &  \\[3pt]
            \myorder{} & \ch{H2Te_\acu} & \rojo{} &  \\[3pt]
            \myorder{} & \ch{H2Se_\acu} & \rojo{} &  \\
        \end{longtable}

    \end{landscape}

\newpage
\DefTblrTemplate{contfoot-text}{normal}{Continued on next page}
    \SetTblrTemplate{contfoot-text}{normal}
    \DefTblrTemplate{conthead-text}{normal}{(Continued)}
    \SetTblrTemplate{conthead-text}{normal}

    \setcounter{numeroformula}{0}

    \begin{landscape}

        % --- tabla antigua ---
        % \begin{longtblr}[caption=Binary compounds]{hlines,vlines, colspec={Q[c,m] Q[c,m, 2cm] Q[c,m,12cm] Q[c,m,12cm] }, rowsep=5pt, rowhead=1}
            % \setcounter{numeroformula}{0}
            % \SetRow{gray!50} No. & Formula         & Multipliers            & Oxidation State            \\
            % \SetCell[c=4]{c} \textbf{Oxides} \\
            % \myorder{}          & \ch{Na2O}       &    \rojo{}                       &             \rojo{}                            \\
            % \myorder{}          & \ch{Co2O3}      &           \rojo{}                 &   \rojo{}                                      \\
            % \myorder{}          & \ch{CuO}        &            \rojo{}                &   \rojo{}                                     \\
            % \myorder{}          & \ch{Ag2O}       &       \rojo{}                     &          \rojo{}                               \\
            % \myorder{}          & \ch{CdO}        &              \rojo{}              &       \rojo{}                                  \\
            % \myorder{}          & \ch{NO}         &           \rojo{}                 &               \rojo{}                        \\
            % \myorder{}          & \ch{CO2}        &              \rojo{}              &      \rojo{}                                   \\
            % \myorder{}          & \ch{As2O5}      &            \rojo{}                &         \rojo{}                                \\
            % \myorder{}          & \ch{SO3}        &               \rojo{}             &        \rojo{}                                 \\
            % \myorder{}          & \ch{O5I2}       &               \rojo{}             &     \rojo{}                                    \\
            % \myorder{}          & \ch{O7Br2}      &                    \rojo{}        &        \rojo{}                                 \\
            % \myorder{}          &   \rojo{\ch{}}          & Dilithium oxide           &         \rojo{}                                \\
            % \myorder{}          &   \rojo{\ch{}}              &        \rojo{}               & Magnesium oxide                       \\
            % \myorder{}          &   \rojo{\ch{}}              &        \rojo{}                    & Potassium oxide                        \\
            % \myorder{}          &   \rojo{\ch{}}              &        \rojo{}                   & Strontium oxide                      \\
            % \myorder{}          &   \rojo{\ch{}}              &        \rojo{}                   & Iron(II) oxide                    \\
            % \myorder{}          &    \rojo{\ch{}}             & Lead dioxide           &         \rojo{}                                \\
            % \myorder{}          &    \rojo{\ch{}}             &       \rojo{}                     & Gold(III) oxide                      \\
            % \myorder{}          &    \rojo{\ch{}}             & Dialuminum trioxide     &           \rojo{}                             \\
            % \myorder{}          &    \rojo{\ch{}}             &       \rojo{}                     & Tin(II) oxide                    \\
            % \myorder{}          &    \rojo{\ch{}}             &       \rojo{}                     & Zinc oxide                           \\
            % \myorder{}          &    \rojo{\ch{}}             &          \rojo{}                  & Mercury(II) oxide                  \\
            % \myorder{}          &    \rojo{\ch{}}             & Dinitrogen tetraoxide  &       \rojo{}                                  \\
            % \myorder{}          &    \rojo{\ch{}}             &               \rojo{}             & Cobalt(III) oxide                  \\
            % \myorder{}          &     \rojo{\ch{}}            & Dinitrogen pentoxide   &          \rojo{}                               \\
            % \myorder{}          &     \rojo{\ch{}}            &                 \rojo{}          & Barium oxide                          \\
            % \myorder{}          &     \rojo{\ch{}}            & Sulfur trioxide         &      \rojo{}                                   \\
            % \myorder{}          &     \rojo{\ch{}}            & Dichromium trioxide        &       \rojo{}                                  \\
            % \myorder{}          &     \rojo{\ch{}}            &                 \rojo{}          & Manganese(IV) oxide                 \\
            % \myorder{}          &     \rojo{\ch{}}            & Dinitrogen monoxide    &        \rojo{}                                 \\
            % \myorder{}          &     \rojo{\ch{}}            &                   \rojo{}        & Platinum(IV) oxide                   \\
            % \myorder{}          &     \rojo{\ch{}}            & Trioxygen dibromide    &       \rojo{}                                 \\
            % \SetCell[c=4]{c} \textbf{Hydrides} \\
            % \myorder{}          & \ch{NaH}        &              \rojo{}              &      \rojo{}                                   \\
            % \myorder{}          & \ch{CaH2}       &                   \rojo{}         &     \rojo{}                                    \\
            % \myorder{}          & \ch{CoH2}       &                          \rojo{}  &      \rojo{}                                   \\
            % \myorder{}          & \ch{CuH2}       &                 \rojo{}           &  \rojo{}                                       \\
            % \myorder{}          & \ch{SnH4}       &         \rojo{}                       &        \rojo{}                                \\
            % \myorder{}          & \ch{LiH}        &                \rojo{}                &     \rojo{}                                    \\
            % \myorder{}          & \ch{HgH2}       &                      \rojo{}      &       \rojo{}                                  \\
            % \myorder{}          & \ch{NiH3}       &                \rojo{}                &      \rojo{}                                   \\
            % \myorder{}          & \ch{AlH3}       &               \rojo{}                 &        \rojo{}                                \\
            % \myorder{}          & \ch{H4Pb}       &                   \rojo{}           &     \rojo{}                                    \\
            % \myorder{}          & \ch{H3Mn}       &                  \rojo{}            &      \rojo{}                                   \\
            % \myorder{}          &       \rojo{\ch{}}          &         \rojo{}                   & Magnesium hydride                     \\
            % \myorder{}          &      \rojo{\ch{}}          &        \rojo{}                    & Sodium hydride                        \\
            % \myorder{}          &        \rojo{\ch{}}         &      \rojo{}                      & Strontium hydride                    \\
            % \myorder{}          &         \rojo{\ch{}}        &      \rojo{}                      & Chromium(II) hydride                   \\
            % \myorder{}          &       \rojo{\ch{}}          & Iron trihydride       &               \rojo{}                          \\
            % \myorder{}          &        \rojo{\ch{}}         &      \rojo{}                      & Lead(IV) hydride                   \\
            % \myorder{}          &          \rojo{\ch{}}       & Zinc dihydride          &            \rojo{}                            \\
            % \myorder{}          &          \rojo{\ch{}}       &     \rojo{}                      & Cadmium hydride                       \\
            % \myorder{}          &         \rojo{\ch{}}        &     \rojo{}                       & Copper(I) hydride                    \\
            % \myorder{}          &          \rojo{\ch{}}       &     \rojo{}                       & Gold(III) hydride                    \\
            % \myorder{}          &         \rojo{\ch{}}        &     \rojo{}                       & Tin(II) hydride                  \\
            % \myorder{}          &          \rojo{\ch{}}       & Potassium hydride         &           \rojo{}                             \\
            % \myorder{}          &           \rojo{\ch{}}      &    \rojo{}                        & Cobalt(II) hydride                 \\
            % \myorder{}          &           \rojo{\ch{}}      &    \rojo{}                        & Platinum(IV) hydride                 \\
            % \myorder{}          &           \rojo{\ch{}}      & Barium dihydride         &         \rojo{}                               \\
            % \myorder{}          &          \rojo{\ch{}}       & Disilver hydride         &         \rojo{}                               \\
            % \myorder{}          &          \rojo{\ch{}}       &      \rojo{}                      & Gold(I) hydride                      \\
            % \SetCell[c=4]{c} \textbf{Acids} \\
            % \myorder{}          & \ch{HCl}        &               \rojo{}             &        \rojo{}                                 \\
            % \myorder{}          & \ch{HBr}        &                 \rojo{}           &      \rojo{}                                   \\
            % \myorder{}          & \ch{HI}         &                      \rojo{}      &         \rojo{}                                \\
            % \myorder{}          & \ch{HF}         &                   \rojo{}         &        \rojo{}                                 \\
            % \myorder{}          & \ch{H2S}        &              \rojo{}              &      \rojo{}                                   \\
            % \myorder{}          & \ch{H2Te}       &             \rojo{}                 &       \rojo{}                                  \\
            % \myorder{}          & \ch{H2Se}       &              \rojo{}                &      \rojo{}                                   \\
            % \myorder{}          &        \rojo{\ch{}}         & Hydrogen chloride       &                                         \\
            % \myorder{}          &        \rojo{\ch{}}         &     \rojo{}                       & Hydrogen bromide                    \\
            % \myorder{}          &         \rojo{\ch{}}        &     \rojo{}                       & Hydrogen iodide                     \\
            % \myorder{}          &         \rojo{\ch{}}        &     \rojo{}                       & Hydrogen fluoride                   \\
            % \myorder{}          &        \rojo{\ch{}}         & Dihydrogen sulfide     &             \rojo{}                            \\
            % \myorder{}          &       \rojo{\ch{}}          &    \rojo{}                        & Hydrogen telluride                     \\
            % \myorder{}          &       \rojo{\ch{}}          &    \rojo{}                        & Hydrogen selenide                  \\
            % \SetCell[c=4]{c} \textbf{Binary salts} \\
            % \myorder{}          & \ch{MgCl2}      &            \rojo{}                &       \rojo{}                                  \\
            % \myorder{}          & \ch{BaBr2}      &                      \rojo{}      &       \rojo{}                                  \\
            % \myorder{}          & \ch{Na3N}       &                 \rojo{}           &        \rojo{}                                 \\
            % \myorder{}          & \ch{K2S}        &               \rojo{}             &        \rojo{}                                 \\
            % \myorder{}          & \ch{CaI2}       &                   \rojo{}         &        \rojo{}                                 \\
            % \myorder{}          & \ch{FeBr3}      &                  \rojo{}          &       \rojo{}                                  \\
            % \myorder{}          & \ch{PbS2}       &               \rojo{}               &       \rojo{}                                  \\
            % \myorder{}          & \ch{Au2S3}      &                 \rojo{}           &       \rojo{}                                  \\
            % \myorder{}          & \ch{HgF2}       &                \rojo{}              &        \rojo{}                                 \\
            % \myorder{}          & \ch{AlCl3}      &                 \rojo{}           &        \rojo{}                                 \\
            % \myorder{}          & \ch{SnBr2}      &                  \rojo{}          &       \rojo{}                                  \\
            % \myorder{}          & \ch{CuI}        &                  \rojo{}            &       \rojo{}                                  \\
            % \myorder{}          & \ch{CI4}        &                  \rojo{}          &             \rojo{}                            \\
            % \myorder{}          & \ch{NF3}        &                 \rojo{}           &               \rojo{}                          \\
            % \myorder{}          & \ch{IBr}        &                      \rojo{}      &                                         \\
            % \myorder{}          & \ch{SeBr4}      &               \rojo{}             &               \rojo{}                          \\
            % \myorder{}          & \ch{SeCl6}      &                    \rojo{}        &               \rojo{}                          \\
            % \myorder{}          & \ch{AsF5}       &                      \rojo{}      &                                         \\
            % \myorder{}          & \ch{CCl4}       &                        \rojo{}    &                \rojo{}                         \\
            % \myorder{}          & \ch{SiF4}       &                        \rojo{}    &                \rojo{}                         \\
            % \myorder{}          &       \rojo{\ch{}}          &       \rojo{}                     & Lithium chloride                        \\
            % \myorder{}          &      \rojo{\ch{}}          & Magnesium diiodide       &               \rojo{}                          \\
            % \myorder{}          &        \rojo{\ch{}}         &      \rojo{}                      & Beryllium bromide                      \\
            % \myorder{}          &       \rojo{\ch{}}          & Dipotassium sulfide       &            \rojo{}                             \\
            % \myorder{}          &        \rojo{\ch{}}         &     \rojo{}                       & Iron (II) fluoride                  \\
            % \myorder{}          &         \rojo{\ch{}}        & Lead tetrachloride      &               \rojo{}                          \\
            % \myorder{}          &         \rojo{\ch{}}        &    \rojo{}                        & Gold (III) sulfide                    \\
            % \myorder{}          &         \rojo{\ch{}}        &    \rojo{}                        & Tin (II) sulfide                  \\
            % \myorder{}          &          \rojo{\ch{}}       & Zinc dibromide          &                  \rojo{}                       \\
            % \myorder{}          &           \rojo{\ch{}}      &       \rojo{}                     & Mercury (II) telluride                 \\
            % \myorder{}          &           \rojo{\ch{}}      &      \rojo{}                      & Cobalt (II) iodide                  \\
            % \myorder{}          &            \rojo{\ch{}}     &      \rojo{}                      & Calcium chloride                       \\
            % \myorder{}          &           \rojo{\ch{}}      & Dialuminum trisulfide   &               \rojo{}                          \\
            % \myorder{}          &           \rojo{\ch{}}      & Dichromium trisulfide      &           \rojo{}                              \\
            % \myorder{}          &          \rojo{\ch{}}       &      \rojo{}                     & Manganese (II) bromide               \\
            % \myorder{}          &            \rojo{\ch{}}     & Disodium telluride          &           \rojo{}                              \\
            % \myorder{}          &           \rojo{\ch{}}     &      \rojo{}                      & Zinc chloride                         \\
            % \myorder{}          &            \rojo{\ch{}}     & Disilver sulfide         &              \rojo{}                           \\
            % \myorder{}          &            \rojo{\ch{}}     & Bromine trifluoride       &            \rojo{}                             \\
            % \myorder{}          &           \rojo{\ch{}}      &     \rojo{}                       & Bromine (I) chloride                    \\
            % \myorder{}          &          \rojo{\ch{}}      & Sulfur hexafluoride     &                 \rojo{}                        \\
            % \myorder{}          &          \rojo{\ch{}}       &     \rojo{}                       & Carbon (IV) sulfide                 \\
            % \SetCell[c=4]{c} \textbf{Hydroxides} \\
            % \myorder{}          & \ch{CsOH}       &            \rojo{}                &               \rojo{}                          \\
            % \myorder{}          & \ch{KOH}        &            \rojo{}                &             \rojo{}                            \\
            % \myorder{}          & \ch{Be(OH)2}    &           \rojo{}                 &               \rojo{}                          \\
            % \myorder{}          & \ch{Fe(OH)3}    &           \rojo{}                 &               \rojo{}                          \\
            % \myorder{}          & \ch{AgOH}       &           \rojo{}                 &               \rojo{}                          \\
            % \myorder{}       & \ch{Al(OH)3}    &             \rojo{}               &                 \rojo{}                        \\
            % \myorder{}       & \ch{NH4OH}      &             \rojo{}               &                  \rojo{}                       \\
            % \myorder{}       & \ch{Cd(OH)2}    &            \rojo{}                &                \rojo{}                         \\
            % \myorder{}       & \ch{Au(OH)3}    &             \rojo{}               &                 \rojo{}                       \\
            % \myorder{}       & \ch{Mn(OH)3}    &            \rojo{}                &                 \rojo{}                        \\
            % \myorder{}       & \ch{Pt(OH)2}    &             \rojo{}               &                  \rojo{}                       \\
            % \myorder{}       &        \rojo{\ch{}}         & Nickel dihydroxide  &                                         \\
            % \myorder{}       &       \rojo{\ch{}}          &     \rojo{}                       & Nickel (III) hydroxide                   \\
            % \myorder{}       &         \rojo{\ch{}}        &     \rojo{}                       & Sodium hydroxide                      \\
            % \myorder{}       &        \rojo{\ch{}}         &     \rojo{}                       & Calcium hydroxide                     \\
            % \myorder{}       &          \rojo{\ch{}}       & Magnesium dihydroxide    &                 \rojo{}                        \\
            % \myorder{}       &        \rojo{\ch{}}         &      \rojo{}                      & Mercury (II) hydroxide              \\
            % \myorder{}       &            \rojo{\ch{}}     & Iron dihydroxide      &                  \rojo{}                       \\
            % \myorder{}       &             \rojo{\ch{}}    &     \rojo{}                       & Zinc hydroxide                       \\
            % \myorder{}       &           \rojo{\ch{}}      & Copper dihydroxide       &               \rojo{}                          \\
            % \myorder{}       &            \rojo{\ch{}}     &     \rojo{}                       & Lead (IV) hydroxide                 \\
            % \myorder{}       &            \rojo{\ch{}}     &     \rojo{}                       & Tin (II) hydroxide                \\
            % \myorder{}       &           \rojo{\ch{}}     &     \rojo{}                       & Nickel (III) hydroxide               \\
            % \myorder{}       &            \rojo{\ch{}}     & Chromium trihydroxide      &           \rojo{}                              \\
            % \myorder{}       &           \rojo{\ch{}}      &    \rojo{}                        & Chromium (II) hydroxide                 \\
            % \SetCell[c=4]{c} { \textbf{Peroxides}. \\ In this case we have only given the nomenclature for peroxides.       \\ You can write their name in the multipliers column leaving the rest blank}       \\
            % \myorder{}       & \ch{Na2O2}       &         \rojo{}                   & --                                      \\
            % \myorder{}       & \ch{K2O2}       &     \rojo{}                       & --                                      \\
            % \myorder{}       & \ch{SrO2}      &      \rojo{}                      & --                                      \\
            % \myorder{}       & \ch{CaO2}       &      \rojo{}                      & --                                     \\
            % \myorder{}       & \ch{Rb2O2}       &     \rojo{}                       & --                                    \\
            % \myorder{} & \ch{CuO2} & & --  \\
            % \myorder{}           & \ch{MgO2} & & -- \\
            % \myorder{}       & \ch{BaO2} & --                            \\
            % \SetCell[c=4]{c} {\textbf{Hydrides from groups 16 and 17 in aqueous solution}. Only one nomenclature has been given.       \\ You can write their name in the multipliers column leaving the rest blank }       \\
            % \myorder{}       & \ch{HCl_{(aq)}}    & \rojo{}            &                \\
            % \myorder{}       & \ch{HBr_\acu}   & \rojo{}               &             \\
            % \myorder{}       & \ch{HI_\acu}  & \rojo{}                &            \\
            % \myorder{}       & \ch{HF_\acu}  & \rojo{}                &              \\
            % \myorder{}       & \ch{H2S_\acu}  & \rojo{}                &                \\
            % \myorder{}       & \ch{H2Te_\acu}  & \rojo{}                &             \\
            % \myorder{}       & \ch{H2Se_\acu}  & \rojo{}                &
        % \end{longtblr}
        % --- tabla nueva ---
        % Ejercicio con las soluciones de la tabla: cada libro carga las suyas
        % (\LibroENG se define en Libro_3_ESO_ENG.tex)
        \begin{exercise}
            \ifdefined\LibroENG
                Complete Table~\ref{tab:binary_compounds}: write the formula or the missing names (multipliers and oxidation state) of each compound.
            \else
                Completa la Tabla~\ref{tab:binary_compounds}: escribe la fórmula o los nombres que faltan (multiplicadores y estado de oxidación) de cada compuesto.
            \fi
            \begin{solu}
                \ifdefined\LibroENG% Soluciones de la tabla «Binary compounds» del capítulo de formulación (inorganica_apuntes_3_ESO_claude_Libro.tex).
% Se carga desde el fichero de soluciones (\begin{solu} del ejercicio de la tabla) en el libro en inglés.
% En negro, los datos del enunciado; en rojo, las respuestas.
% Versión en español: Soluciones_formulacion_3ESO_ES.tex (mismos compuestos y números).

\providecommand{\cabeceraSolFormEN}[1]{%
    \toprule
    \textbf{No.} & \textbf{Formula} & \textbf{Multipliers} & \textbf{Oxidation state} \\
    \midrule
    \multicolumn{4}{c}{\textbf{#1}} \\
    \midrule
}

{\small
\begin{center}
    \begin{tabular}{@{}r C{1.9cm} L{5cm} L{5cm}@{}}
        \cabeceraSolFormEN{Oxides}
        1 & \ch{Na2O} & \rojo{Disodium oxide} & \rojo{Sodium oxide} \\[2pt]
        2 & \ch{Co2O3} & \rojo{Dicobalt trioxide} & \rojo{Cobalt (III) oxide} \\[2pt]
        3 & \ch{CuO} & \rojo{Copper oxide (copper monoxide)} & \rojo{Copper (II) oxide} \\[2pt]
        4 & \ch{Ag2O} & \rojo{Disilver oxide} & \rojo{Silver oxide} \\[2pt]
        5 & \ch{CdO} & \rojo{Cadmium oxide} & \rojo{Cadmium oxide} \\[2pt]
        6 & \ch{NO} & \rojo{Nitrogen monoxide} & \rojo{Nitrogen (II) oxide} \\[2pt]
        7 & \ch{CO2} & \rojo{Carbon dioxide} & \rojo{Carbon (IV) oxide} \\[2pt]
        8 & \ch{As2O5} & \rojo{Diarsenic pentoxide} & \rojo{Arsenic (V) oxide} \\[2pt]
        9 & \ch{SO3} & \rojo{Sulfur trioxide} & \rojo{Sulfur (VI) oxide} \\[2pt]
        10 & \ch{O5I2} & \rojo{Pentaoxygen diiodide} & -- \\[2pt]
        11 & \ch{O7Br2} & \rojo{Heptaoxygen dibromide} & -- \\[2pt]
        12 & \rojo{\ch{Li2O}} & Dilithium oxide & \rojo{Lithium oxide} \\[2pt]
        13 & \rojo{\ch{MgO}} & \rojo{Magnesium oxide} & Magnesium oxide \\[2pt]
        14 & \rojo{\ch{K2O}} & \rojo{Dipotassium oxide} & Potassium oxide \\[2pt]
        15 & \rojo{\ch{SrO}} & \rojo{Strontium oxide} & Strontium oxide \\[2pt]
        16 & \rojo{\ch{FeO}} & \rojo{Iron oxide (iron monoxide)} & Iron (II) oxide \\[2pt]
        17 & \rojo{\ch{PbO2}} & Lead dioxide & \rojo{Lead (IV) oxide} \\[2pt]
        18 & \rojo{\ch{Au2O3}} & \rojo{Digold trioxide} & Gold (III) oxide \\[2pt]
        19 & \rojo{\ch{Al2O3}} & Dialuminum trioxide & \rojo{Aluminum oxide} \\[2pt]
        20 & \rojo{\ch{SnO}} & \rojo{Tin oxide (tin monoxide)} & Tin (II) oxide \\[2pt]
        21 & \rojo{\ch{ZnO}} & \rojo{Zinc oxide} & Zinc oxide \\[2pt]
        22 & \rojo{\ch{HgO}} & \rojo{Mercury oxide (mercury monoxide)} & Mercury (II) oxide \\[2pt]
        23 & \rojo{\ch{N2O4}} & Dinitrogen tetraoxide & \rojo{Nitrogen (IV) oxide} \\[2pt]
        24 & \rojo{\ch{Co2O3}} & \rojo{Dicobalt trioxide} & Cobalt (III) oxide \\[2pt]
        25 & \rojo{\ch{N2O5}} & Dinitrogen pentoxide & \rojo{Nitrogen (V) oxide} \\[2pt]
        26 & \rojo{\ch{BaO}} & \rojo{Barium oxide} & Barium oxide \\[2pt]
        27 & \rojo{\ch{SO3}} & Sulfur trioxide & \rojo{Sulfur (VI) oxide} \\[2pt]
        28 & \rojo{\ch{Cr2O3}} & Dichromium trioxide & \rojo{Chromium (III) oxide} \\[2pt]
        29 & \rojo{\ch{MnO2}} & \rojo{Manganese dioxide} & Manganese (IV) oxide \\[2pt]
        30 & \rojo{\ch{N2O}} & Dinitrogen monoxide & \rojo{Nitrogen (I) oxide} \\[2pt]
        31 & \rojo{\ch{PtO2}} & \rojo{Platinum dioxide} & Platinum (IV) oxide \\[2pt]
        32 & \rojo{\ch{O3Br2}} & Trioxygen dibromide & -- \\
        \bottomrule
    \end{tabular}
\end{center}

\noindent In 10, 11 and 32 the halogen is written on the right (it is more electronegative than oxygen), so they are named only with multipliers.

\begin{center}
    \begin{tabular}{@{}r C{1.9cm} L{5cm} L{5cm}@{}}
        \cabeceraSolFormEN{Hydrides}
        33 & \ch{NaH} & \rojo{Sodium hydride} & \rojo{Sodium hydride} \\[2pt]
        34 & \ch{CaH2} & \rojo{Calcium dihydride} & \rojo{Calcium hydride} \\[2pt]
        35 & \ch{CoH2} & \rojo{Cobalt dihydride} & \rojo{Cobalt (II) hydride} \\[2pt]
        36 & \ch{CuH2} & \rojo{Copper dihydride} & \rojo{Copper (II) hydride} \\[2pt]
        37 & \ch{SnH4} & \rojo{Tin tetrahydride} & \rojo{Tin (IV) hydride} \\[2pt]
        38 & \ch{LiH} & \rojo{Lithium hydride} & \rojo{Lithium hydride} \\[2pt]
        39 & \ch{HgH2} & \rojo{Mercury dihydride} & \rojo{Mercury (II) hydride} \\[2pt]
        40 & \ch{NiH3} & \rojo{Nickel trihydride} & \rojo{Nickel (III) hydride} \\[2pt]
        41 & \ch{AlH3} & \rojo{Aluminum trihydride} & \rojo{Aluminum hydride} \\[2pt]
        42 & \ch{H4Pb} & \rojo{Lead tetrahydride} & \rojo{Lead (IV) hydride} \\[2pt]
        43 & \ch{H3Mn} & \rojo{Manganese trihydride} & \rojo{Manganese (III) hydride} \\[2pt]
        44 & \rojo{\ch{MgH2}} & \rojo{Magnesium dihydride} & Magnesium hydride \\[2pt]
        45 & \rojo{\ch{NaH}} & \rojo{Sodium hydride} & Sodium hydride \\[2pt]
        46 & \rojo{\ch{SrH2}} & \rojo{Strontium dihydride} & Strontium hydride \\[2pt]
        47 & \rojo{\ch{CrH2}} & \rojo{Chromium dihydride} & Chromium (II) hydride \\[2pt]
        48 & \rojo{\ch{FeH3}} & Iron trihydride & \rojo{Iron (III) hydride} \\[2pt]
        49 & \rojo{\ch{PbH4}} & \rojo{Lead tetrahydride} & Lead (IV) hydride \\[2pt]
        50 & \rojo{\ch{ZnH2}} & Zinc dihydride & \rojo{Zinc hydride} \\[2pt]
        51 & \rojo{\ch{CdH2}} & \rojo{Cadmium dihydride} & Cadmium hydride \\[2pt]
        52 & \rojo{\ch{CuH}} & \rojo{Copper hydride} & Copper (I) hydride \\[2pt]
        53 & \rojo{\ch{AuH3}} & \rojo{Gold trihydride} & Gold (III) hydride \\[2pt]
        54 & \rojo{\ch{SnH2}} & \rojo{Tin dihydride} & Tin (II) hydride \\[2pt]
        55 & \rojo{\ch{KH}} & Potassium hydride & \rojo{Potassium hydride} \\[2pt]
        56 & \rojo{\ch{CoH2}} & \rojo{Cobalt dihydride} & Cobalt (II) hydride \\[2pt]
        57 & \rojo{\ch{PtH4}} & \rojo{Platinum tetrahydride} & Platinum (IV) hydride \\[2pt]
        58 & \rojo{\ch{BaH2}} & Barium dihydride & \rojo{Barium hydride} \\[2pt]
        59 & \rojo{\ch{AgH}} & Silver hydride & \rojo{Silver hydride} \\[2pt]
        60 & \rojo{\ch{AuH}} & \rojo{Gold hydride} & Gold (I) hydride \\
        \bottomrule
    \end{tabular}
\end{center}

\noindent In 42 and 43 the usual way to write the formula is with the metal first: \ch{PbH4} and \ch{MnH3}.

\begin{center}
    \begin{tabular}{@{}r C{1.9cm} L{5cm} L{5cm}@{}}
        \cabeceraSolFormEN{Acids (hydrogen with a nonmetal of groups 16 and 17)}
        61 & \ch{HCl} & \rojo{Hydrogen chloride} & \rojo{Hydrogen chloride} \\[2pt]
        62 & \ch{HBr} & \rojo{Hydrogen bromide} & \rojo{Hydrogen bromide} \\[2pt]
        63 & \ch{HI} & \rojo{Hydrogen iodide} & \rojo{Hydrogen iodide} \\[2pt]
        64 & \ch{HF} & \rojo{Hydrogen fluoride} & \rojo{Hydrogen fluoride} \\[2pt]
        65 & \ch{H2S} & \rojo{Dihydrogen sulfide} & \rojo{Hydrogen sulfide} \\[2pt]
        66 & \ch{H2Te} & \rojo{Dihydrogen telluride} & \rojo{Hydrogen telluride} \\[2pt]
        67 & \ch{H2Se} & \rojo{Dihydrogen selenide} & \rojo{Hydrogen selenide} \\[2pt]
        68 & \rojo{\ch{HCl}} & Hydrogen chloride & \rojo{Hydrogen chloride} \\[2pt]
        69 & \rojo{\ch{HBr}} & \rojo{Hydrogen bromide} & Hydrogen bromide \\[2pt]
        70 & \rojo{\ch{HI}} & \rojo{Hydrogen iodide} & Hydrogen iodide \\[2pt]
        71 & \rojo{\ch{HF}} & \rojo{Hydrogen fluoride} & Hydrogen fluoride \\[2pt]
        72 & \rojo{\ch{H2S}} & Dihydrogen sulfide & \rojo{Hydrogen sulfide} \\[2pt]
        73 & \rojo{\ch{H2Te}} & \rojo{Dihydrogen telluride} & Hydrogen telluride \\[2pt]
        74 & \rojo{\ch{H2Se}} & \rojo{Dihydrogen selenide} & Hydrogen selenide \\
        \bottomrule
    \end{tabular}
\end{center}

\noindent Hydrogen has only one oxidation state (+1), so the second name is the same without prefixes.

\begin{center}
    \begin{tabular}{@{}r C{1.9cm} L{5cm} L{5cm}@{}}
        \cabeceraSolFormEN{Binary salts}
        75 & \ch{MgCl2} & \rojo{Magnesium dichloride} & \rojo{Magnesium chloride} \\[2pt]
        76 & \ch{BaBr2} & \rojo{Barium dibromide} & \rojo{Barium bromide} \\[2pt]
        77 & \ch{Na3N} & \rojo{Trisodium nitride} & \rojo{Sodium nitride} \\[2pt]
        78 & \ch{K2S} & \rojo{Dipotassium sulfide} & \rojo{Potassium sulfide} \\[2pt]
        79 & \ch{CaI2} & \rojo{Calcium diiodide} & \rojo{Calcium iodide} \\[2pt]
        80 & \ch{FeBr3} & \rojo{Iron tribromide} & \rojo{Iron (III) bromide} \\[2pt]
        81 & \ch{PbS2} & \rojo{Lead disulfide} & \rojo{Lead (IV) sulfide} \\[2pt]
        82 & \ch{Au2S3} & \rojo{Digold trisulfide} & \rojo{Gold (III) sulfide} \\[2pt]
        83 & \ch{HgF2} & \rojo{Mercury difluoride} & \rojo{Mercury (II) fluoride} \\[2pt]
        84 & \ch{AlCl3} & \rojo{Aluminum trichloride} & \rojo{Aluminum chloride} \\[2pt]
        85 & \ch{SnBr2} & \rojo{Tin dibromide} & \rojo{Tin (II) bromide} \\[2pt]
        86 & \ch{CuI} & \rojo{Copper iodide} & \rojo{Copper (I) iodide} \\[2pt]
        87 & \ch{CI4} & \rojo{Carbon tetraiodide} & \rojo{Carbon (IV) iodide} \\[2pt]
        88 & \ch{NF3} & \rojo{Nitrogen trifluoride} & \rojo{Nitrogen (III) fluoride} \\[2pt]
        89 & \ch{IBr} & \rojo{Iodine bromide} & \rojo{Iodine (I) bromide} \\[2pt]
        90 & \ch{SeBr4} & \rojo{Selenium tetrabromide} & \rojo{Selenium (IV) bromide} \\[2pt]
        91 & \ch{SeCl6} & \rojo{Selenium hexachloride} & \rojo{Selenium (VI) chloride} \\[2pt]
        92 & \ch{AsF5} & \rojo{Arsenic pentafluoride} & \rojo{Arsenic (V) fluoride} \\[2pt]
        93 & \ch{CCl4} & \rojo{Carbon tetrachloride} & \rojo{Carbon (IV) chloride} \\[2pt]
        94 & \ch{SiF4} & \rojo{Silicon tetrafluoride} & \rojo{Silicon fluoride} \\
        \bottomrule
    \end{tabular}
\end{center}

\begin{center}
    \begin{tabular}{@{}r C{1.9cm} L{5cm} L{5cm}@{}}
        \cabeceraSolFormEN{Binary salts (continued)}
        95 & \rojo{\ch{LiCl}} & \rojo{Lithium chloride} & Lithium chloride \\[2pt]
        96 & \rojo{\ch{MgI2}} & Magnesium diiodide & \rojo{Magnesium iodide} \\[2pt]
        97 & \rojo{\ch{BeBr2}} & \rojo{Beryllium dibromide} & Beryllium bromide \\[2pt]
        98 & \rojo{\ch{K2S}} & Dipotassium sulfide & \rojo{Potassium sulfide} \\[2pt]
        99 & \rojo{\ch{FeF2}} & \rojo{Iron difluoride} & Iron (II) fluoride \\[2pt]
        100 & \rojo{\ch{PbCl4}} & Lead tetrachloride & \rojo{Lead (IV) chloride} \\[2pt]
        101 & \rojo{\ch{Au2S3}} & \rojo{Digold trisulfide} & Gold (III) sulfide \\[2pt]
        102 & \rojo{\ch{SnS}} & \rojo{Tin sulfide (tin monosulfide)} & Tin (II) sulfide \\[2pt]
        103 & \rojo{\ch{ZnBr2}} & Zinc dibromide & \rojo{Zinc bromide} \\[2pt]
        104 & \rojo{\ch{HgTe}} & \rojo{Mercury telluride} & Mercury (II) telluride \\[2pt]
        105 & \rojo{\ch{CoI2}} & \rojo{Cobalt diiodide} & Cobalt (II) iodide \\[2pt]
        106 & \rojo{\ch{CaCl2}} & \rojo{Calcium dichloride} & Calcium chloride \\[2pt]
        107 & \rojo{\ch{Al2S3}} & Dialuminum trisulfide & \rojo{Aluminum sulfide} \\[2pt]
        108 & \rojo{\ch{Cr2S3}} & Dichromium trisulfide & \rojo{Chromium (III) sulfide} \\[2pt]
        109 & \rojo{\ch{MnBr2}} & \rojo{Manganese dibromide} & Manganese (II) bromide \\[2pt]
        110 & \rojo{\ch{Na2Te}} & Disodium telluride & \rojo{Sodium telluride} \\[2pt]
        111 & \rojo{\ch{ZnCl2}} & \rojo{Zinc dichloride} & Zinc chloride \\[2pt]
        112 & \rojo{\ch{Ag2S}} & Disilver sulfide & \rojo{Silver sulfide} \\[2pt]
        113 & \rojo{\ch{BrF3}} & Bromine trifluoride & \rojo{Bromine (III) fluoride} \\[2pt]
        114 & \rojo{\ch{BrCl}} & \rojo{Bromine chloride} & Bromine (I) chloride \\[2pt]
        115 & \rojo{\ch{SF6}} & Sulfur hexafluoride & \rojo{Sulfur (VI) fluoride} \\[2pt]
        116 & \rojo{\ch{CS2}} & \rojo{Carbon disulfide} & Carbon (IV) sulfide \\
        \bottomrule
    \end{tabular}
\end{center}

\begin{center}
    \begin{tabular}{@{}r C{1.9cm} L{5cm} L{5cm}@{}}
        \cabeceraSolFormEN{Hydroxides}
        117 & \ch{CsOH} & \rojo{Cesium hydroxide} & \rojo{Cesium hydroxide} \\[2pt]
        118 & \ch{KOH} & \rojo{Potassium hydroxide} & \rojo{Potassium hydroxide} \\[2pt]
        119 & \ch{Be(OH)2} & \rojo{Beryllium dihydroxide} & \rojo{Beryllium hydroxide} \\[2pt]
        120 & \ch{Fe(OH)3} & \rojo{Iron trihydroxide} & \rojo{Iron (III) hydroxide} \\[2pt]
        121 & \ch{AgOH} & \rojo{Silver hydroxide} & \rojo{Silver hydroxide} \\[2pt]
        122 & \ch{Al(OH)3} & \rojo{Aluminum trihydroxide} & \rojo{Aluminum hydroxide} \\[2pt]
        123 & \ch{NH4OH} & \rojo{Ammonium hydroxide} & \rojo{Ammonium hydroxide} \\[2pt]
        124 & \ch{Cd(OH)2} & \rojo{Cadmium dihydroxide} & \rojo{Cadmium hydroxide} \\[2pt]
        125 & \ch{Au(OH)3} & \rojo{Gold trihydroxide} & \rojo{Gold (III) hydroxide} \\[2pt]
        126 & \ch{Mn(OH)3} & \rojo{Manganese trihydroxide} & \rojo{Manganese (III) hydroxide} \\[2pt]
        127 & \ch{Pt(OH)2} & \rojo{Platinum dihydroxide} & \rojo{Platinum (II) hydroxide} \\[2pt]
        128 & \rojo{\ch{Ni(OH)2}} & Nickel dihydroxide & \rojo{Nickel (II) hydroxide} \\[2pt]
        129 & \rojo{\ch{Ni(OH)3}} & \rojo{Nickel trihydroxide} & Nickel (III) hydroxide \\[2pt]
        130 & \rojo{\ch{NaOH}} & \rojo{Sodium hydroxide} & Sodium hydroxide \\[2pt]
        131 & \rojo{\ch{Ca(OH)2}} & \rojo{Calcium dihydroxide} & Calcium hydroxide \\[2pt]
        132 & \rojo{\ch{Mg(OH)2}} & Magnesium dihydroxide & \rojo{Magnesium hydroxide} \\[2pt]
        133 & \rojo{\ch{Hg(OH)2}} & \rojo{Mercury dihydroxide} & Mercury (II) hydroxide \\[2pt]
        134 & \rojo{\ch{Fe(OH)2}} & Iron dihydroxide & \rojo{Iron (II) hydroxide} \\[2pt]
        135 & \rojo{\ch{Zn(OH)2}} & \rojo{Zinc dihydroxide} & Zinc hydroxide \\[2pt]
        136 & \rojo{\ch{Cu(OH)2}} & Copper dihydroxide & \rojo{Copper (II) hydroxide} \\[2pt]
        137 & \rojo{\ch{Pb(OH)4}} & \rojo{Lead tetrahydroxide} & Lead (IV) hydroxide \\[2pt]
        138 & \rojo{\ch{Sn(OH)2}} & \rojo{Tin dihydroxide} & Tin (II) hydroxide \\[2pt]
        139 & \rojo{\ch{Ni(OH)3}} & \rojo{Nickel trihydroxide} & Nickel (III) hydroxide \\[2pt]
        140 & \rojo{\ch{Cr(OH)3}} & Chromium trihydroxide & \rojo{Chromium (III) hydroxide} \\[2pt]
        141 & \rojo{\ch{Cr(OH)2}} & \rojo{Chromium dihydroxide} & Chromium (II) hydroxide \\
        \bottomrule
    \end{tabular}
\end{center}

\begin{center}
    \begin{tabular}{@{}r C{1.9cm} L{5cm} L{5cm}@{}}
        \cabeceraSolFormEN{Peroxides}
        142 & \ch{Na2O2} & \rojo{Sodium peroxide} & -- \\[2pt]
        143 & \ch{K2O2} & \rojo{Potassium peroxide} & -- \\[2pt]
        144 & \ch{SrO2} & \rojo{Strontium peroxide} & -- \\[2pt]
        145 & \ch{CaO2} & \rojo{Calcium peroxide} & -- \\[2pt]
        146 & \ch{Rb2O2} & \rojo{Rubidium peroxide} & -- \\[2pt]
        147 & \ch{CuO2} & \rojo{Copper (II) peroxide} & -- \\[2pt]
        148 & \ch{MgO2} & \rojo{Magnesium peroxide} & -- \\[2pt]
        149 & \ch{BaO2} & \rojo{Barium peroxide} & -- \\
        \midrule
        \multicolumn{4}{c}{\textbf{Hydrides from groups 16 and 17 in aqueous solution}} \\
        \midrule
        150 & \ch{HCl(aq)} & \rojo{Hydrochloric acid} & -- \\[2pt]
        151 & \ch{HBr(aq)} & \rojo{Hydrobromic acid} & -- \\[2pt]
        152 & \ch{HI(aq)} & \rojo{Hydroiodic acid} & -- \\[2pt]
        153 & \ch{HF(aq)} & \rojo{Hydrofluoric acid} & -- \\[2pt]
        154 & \ch{H2S(aq)} & \rojo{Hydrosulfuric acid} & -- \\[2pt]
        155 & \ch{H2Te(aq)} & \rojo{Hydrotelluric acid} & -- \\[2pt]
        156 & \ch{H2Se(aq)} & \rojo{Hydroselenic acid} & -- \\
        \bottomrule
    \end{tabular}
\end{center}
}
\else% Soluciones de la tabla «Binary compounds» del capítulo de formulación (inorganica_apuntes_3_ESO_claude_Libro.tex).
% Se carga desde el fichero de soluciones (\begin{solu} del ejercicio de la tabla) en el libro en español.
% En negro, los datos del enunciado; en rojo, las respuestas.
% Traducción de Soluciones_formulacion_3ESO_ENG.tex (mismos compuestos y números).

\providecommand{\cabeceraSolFormES}[1]{%
    \toprule
    \textbf{Nº} & \textbf{Fórmula} & \textbf{Multiplicadores} & \textbf{Estado de oxidación} \\
    \midrule
    \multicolumn{4}{c}{\textbf{#1}} \\
    \midrule
}

{\small
\begin{center}
    \begin{tabular}{@{}r C{1.9cm} L{5cm} L{5cm}@{}}
        \cabeceraSolFormES{Óxidos}
        1 & \ch{Na2O} & \rojo{Óxido de disodio} & \rojo{Óxido de sodio} \\[2pt]
        2 & \ch{Co2O3} & \rojo{Trióxido de dicobalto} & \rojo{Óxido de cobalto (III)} \\[2pt]
        3 & \ch{CuO} & \rojo{Óxido de cobre (monóxido de cobre)} & \rojo{Óxido de cobre (II)} \\[2pt]
        4 & \ch{Ag2O} & \rojo{Óxido de diplata} & \rojo{Óxido de plata} \\[2pt]
        5 & \ch{CdO} & \rojo{Óxido de cadmio} & \rojo{Óxido de cadmio} \\[2pt]
        6 & \ch{NO} & \rojo{Monóxido de nitrógeno} & \rojo{Óxido de nitrógeno (II)} \\[2pt]
        7 & \ch{CO2} & \rojo{Dióxido de carbono} & \rojo{Óxido de carbono (IV)} \\[2pt]
        8 & \ch{As2O5} & \rojo{Pentaóxido de diarsénico} & \rojo{Óxido de arsénico (V)} \\[2pt]
        9 & \ch{SO3} & \rojo{Trióxido de azufre} & \rojo{Óxido de azufre (VI)} \\[2pt]
        10 & \ch{O5I2} & \rojo{Diyoduro de pentaoxígeno} & -- \\[2pt]
        11 & \ch{O7Br2} & \rojo{Dibromuro de heptaoxígeno} & -- \\[2pt]
        12 & \rojo{\ch{Li2O}} & Óxido de dilitio & \rojo{Óxido de litio} \\[2pt]
        13 & \rojo{\ch{MgO}} & \rojo{Óxido de magnesio} & Óxido de magnesio \\[2pt]
        14 & \rojo{\ch{K2O}} & \rojo{Óxido de dipotasio} & Óxido de potasio \\[2pt]
        15 & \rojo{\ch{SrO}} & \rojo{Óxido de estroncio} & Óxido de estroncio \\[2pt]
        16 & \rojo{\ch{FeO}} & \rojo{Óxido de hierro (monóxido de hierro)} & Óxido de hierro (II) \\[2pt]
        17 & \rojo{\ch{PbO2}} & Dióxido de plomo & \rojo{Óxido de plomo (IV)} \\[2pt]
        18 & \rojo{\ch{Au2O3}} & \rojo{Trióxido de dioro} & Óxido de oro (III) \\[2pt]
        19 & \rojo{\ch{Al2O3}} & Trióxido de dialuminio & \rojo{Óxido de aluminio} \\[2pt]
        20 & \rojo{\ch{SnO}} & \rojo{Óxido de estaño (monóxido de estaño)} & Óxido de estaño (II) \\[2pt]
        21 & \rojo{\ch{ZnO}} & \rojo{Óxido de cinc} & Óxido de cinc \\[2pt]
        22 & \rojo{\ch{HgO}} & \rojo{Óxido de mercurio (monóxido de mercurio)} & Óxido de mercurio (II) \\[2pt]
        23 & \rojo{\ch{N2O4}} & Tetraóxido de dinitrógeno & \rojo{Óxido de nitrógeno (IV)} \\[2pt]
        24 & \rojo{\ch{Co2O3}} & \rojo{Trióxido de dicobalto} & Óxido de cobalto (III) \\[2pt]
        25 & \rojo{\ch{N2O5}} & Pentaóxido de dinitrógeno & \rojo{Óxido de nitrógeno (V)} \\[2pt]
        26 & \rojo{\ch{BaO}} & \rojo{Óxido de bario} & Óxido de bario \\[2pt]
        27 & \rojo{\ch{SO3}} & Trióxido de azufre & \rojo{Óxido de azufre (VI)} \\[2pt]
        28 & \rojo{\ch{Cr2O3}} & Trióxido de dicromo & \rojo{Óxido de cromo (III)} \\[2pt]
        29 & \rojo{\ch{MnO2}} & \rojo{Dióxido de manganeso} & Óxido de manganeso (IV) \\[2pt]
        30 & \rojo{\ch{N2O}} & Monóxido de dinitrógeno & \rojo{Óxido de nitrógeno (I)} \\[2pt]
        31 & \rojo{\ch{PtO2}} & \rojo{Dióxido de platino} & Óxido de platino (IV) \\[2pt]
        32 & \rojo{\ch{O3Br2}} & Dibromuro de trioxígeno & -- \\
        \bottomrule
    \end{tabular}
\end{center}

\noindent En 10, 11 y 32 el halógeno va a la derecha (es más electronegativo que el oxígeno), así que solo se nombran con multiplicadores.

\begin{center}
    \begin{tabular}{@{}r C{1.9cm} L{5cm} L{5cm}@{}}
        \cabeceraSolFormES{Hidruros}
        33 & \ch{NaH} & \rojo{Hidruro de sodio} & \rojo{Hidruro de sodio} \\[2pt]
        34 & \ch{CaH2} & \rojo{Dihidruro de calcio} & \rojo{Hidruro de calcio} \\[2pt]
        35 & \ch{CoH2} & \rojo{Dihidruro de cobalto} & \rojo{Hidruro de cobalto (II)} \\[2pt]
        36 & \ch{CuH2} & \rojo{Dihidruro de cobre} & \rojo{Hidruro de cobre (II)} \\[2pt]
        37 & \ch{SnH4} & \rojo{Tetrahidruro de estaño} & \rojo{Hidruro de estaño (IV)} \\[2pt]
        38 & \ch{LiH} & \rojo{Hidruro de litio} & \rojo{Hidruro de litio} \\[2pt]
        39 & \ch{HgH2} & \rojo{Dihidruro de mercurio} & \rojo{Hidruro de mercurio (II)} \\[2pt]
        40 & \ch{NiH3} & \rojo{Trihidruro de níquel} & \rojo{Hidruro de níquel (III)} \\[2pt]
        41 & \ch{AlH3} & \rojo{Trihidruro de aluminio} & \rojo{Hidruro de aluminio} \\[2pt]
        42 & \ch{H4Pb} & \rojo{Tetrahidruro de plomo} & \rojo{Hidruro de plomo (IV)} \\[2pt]
        43 & \ch{H3Mn} & \rojo{Trihidruro de manganeso} & \rojo{Hidruro de manganeso (III)} \\[2pt]
        44 & \rojo{\ch{MgH2}} & \rojo{Dihidruro de magnesio} & Hidruro de magnesio \\[2pt]
        45 & \rojo{\ch{NaH}} & \rojo{Hidruro de sodio} & Hidruro de sodio \\[2pt]
        46 & \rojo{\ch{SrH2}} & \rojo{Dihidruro de estroncio} & Hidruro de estroncio \\[2pt]
        47 & \rojo{\ch{CrH2}} & \rojo{Dihidruro de cromo} & Hidruro de cromo (II) \\[2pt]
        48 & \rojo{\ch{FeH3}} & Trihidruro de hierro & \rojo{Hidruro de hierro (III)} \\[2pt]
        49 & \rojo{\ch{PbH4}} & \rojo{Tetrahidruro de plomo} & Hidruro de plomo (IV) \\[2pt]
        50 & \rojo{\ch{ZnH2}} & Dihidruro de cinc & \rojo{Hidruro de cinc} \\[2pt]
        51 & \rojo{\ch{CdH2}} & \rojo{Dihidruro de cadmio} & Hidruro de cadmio \\[2pt]
        52 & \rojo{\ch{CuH}} & \rojo{Hidruro de cobre} & Hidruro de cobre (I) \\[2pt]
        53 & \rojo{\ch{AuH3}} & \rojo{Trihidruro de oro} & Hidruro de oro (III) \\[2pt]
        54 & \rojo{\ch{SnH2}} & \rojo{Dihidruro de estaño} & Hidruro de estaño (II) \\[2pt]
        55 & \rojo{\ch{KH}} & Hidruro de potasio & \rojo{Hidruro de potasio} \\[2pt]
        56 & \rojo{\ch{CoH2}} & \rojo{Dihidruro de cobalto} & Hidruro de cobalto (II) \\[2pt]
        57 & \rojo{\ch{PtH4}} & \rojo{Tetrahidruro de platino} & Hidruro de platino (IV) \\[2pt]
        58 & \rojo{\ch{BaH2}} & Dihidruro de bario & \rojo{Hidruro de bario} \\[2pt]
        59 & \rojo{\ch{AgH}} & Hidruro de plata & \rojo{Hidruro de plata} \\[2pt]
        60 & \rojo{\ch{AuH}} & \rojo{Hidruro de oro} & Hidruro de oro (I) \\
        \bottomrule
    \end{tabular}
\end{center}

\noindent En 42 y 43 lo habitual es escribir primero el metal: \ch{PbH4} y \ch{MnH3}.

\begin{center}
    \begin{tabular}{@{}r C{1.9cm} L{5cm} L{5cm}@{}}
        \cabeceraSolFormES{Hidruros no metálicos (hidrógeno con un no metal de los grupos 16 y 17)}
        61 & \ch{HCl} & \rojo{Cloruro de hidrógeno} & \rojo{Cloruro de hidrógeno} \\[2pt]
        62 & \ch{HBr} & \rojo{Bromuro de hidrógeno} & \rojo{Bromuro de hidrógeno} \\[2pt]
        63 & \ch{HI} & \rojo{Yoduro de hidrógeno} & \rojo{Yoduro de hidrógeno} \\[2pt]
        64 & \ch{HF} & \rojo{Fluoruro de hidrógeno} & \rojo{Fluoruro de hidrógeno} \\[2pt]
        65 & \ch{H2S} & \rojo{Sulfuro de dihidrógeno} & \rojo{Sulfuro de hidrógeno} \\[2pt]
        66 & \ch{H2Te} & \rojo{Telururo de dihidrógeno} & \rojo{Telururo de hidrógeno} \\[2pt]
        67 & \ch{H2Se} & \rojo{Seleniuro de dihidrógeno} & \rojo{Seleniuro de hidrógeno} \\[2pt]
        68 & \rojo{\ch{HCl}} & Cloruro de hidrógeno & \rojo{Cloruro de hidrógeno} \\[2pt]
        69 & \rojo{\ch{HBr}} & \rojo{Bromuro de hidrógeno} & Bromuro de hidrógeno \\[2pt]
        70 & \rojo{\ch{HI}} & \rojo{Yoduro de hidrógeno} & Yoduro de hidrógeno \\[2pt]
        71 & \rojo{\ch{HF}} & \rojo{Fluoruro de hidrógeno} & Fluoruro de hidrógeno \\[2pt]
        72 & \rojo{\ch{H2S}} & Sulfuro de dihidrógeno & \rojo{Sulfuro de hidrógeno} \\[2pt]
        73 & \rojo{\ch{H2Te}} & \rojo{Telururo de dihidrógeno} & Telururo de hidrógeno \\[2pt]
        74 & \rojo{\ch{H2Se}} & \rojo{Seleniuro de dihidrógeno} & Seleniuro de hidrógeno \\
        \bottomrule
    \end{tabular}
\end{center}

\noindent El hidrógeno solo tiene un estado de oxidación (+1), así que el segundo nombre es el mismo sin prefijos.

\begin{center}
    \begin{tabular}{@{}r C{1.9cm} L{5cm} L{5cm}@{}}
        \cabeceraSolFormES{Sales binarias}
        75 & \ch{MgCl2} & \rojo{Dicloruro de magnesio} & \rojo{Cloruro de magnesio} \\[2pt]
        76 & \ch{BaBr2} & \rojo{Dibromuro de bario} & \rojo{Bromuro de bario} \\[2pt]
        77 & \ch{Na3N} & \rojo{Nitruro de trisodio} & \rojo{Nitruro de sodio} \\[2pt]
        78 & \ch{K2S} & \rojo{Sulfuro de dipotasio} & \rojo{Sulfuro de potasio} \\[2pt]
        79 & \ch{CaI2} & \rojo{Diyoduro de calcio} & \rojo{Yoduro de calcio} \\[2pt]
        80 & \ch{FeBr3} & \rojo{Tribromuro de hierro} & \rojo{Bromuro de hierro (III)} \\[2pt]
        81 & \ch{PbS2} & \rojo{Disulfuro de plomo} & \rojo{Sulfuro de plomo (IV)} \\[2pt]
        82 & \ch{Au2S3} & \rojo{Trisulfuro de dioro} & \rojo{Sulfuro de oro (III)} \\[2pt]
        83 & \ch{HgF2} & \rojo{Difluoruro de mercurio} & \rojo{Fluoruro de mercurio (II)} \\[2pt]
        84 & \ch{AlCl3} & \rojo{Tricloruro de aluminio} & \rojo{Cloruro de aluminio} \\[2pt]
        85 & \ch{SnBr2} & \rojo{Dibromuro de estaño} & \rojo{Bromuro de estaño (II)} \\[2pt]
        86 & \ch{CuI} & \rojo{Yoduro de cobre} & \rojo{Yoduro de cobre (I)} \\[2pt]
        87 & \ch{CI4} & \rojo{Tetrayoduro de carbono} & \rojo{Yoduro de carbono (IV)} \\[2pt]
        88 & \ch{NF3} & \rojo{Trifluoruro de nitrógeno} & \rojo{Fluoruro de nitrógeno (III)} \\[2pt]
        89 & \ch{IBr} & \rojo{Bromuro de yodo} & \rojo{Bromuro de yodo (I)} \\[2pt]
        90 & \ch{SeBr4} & \rojo{Tetrabromuro de selenio} & \rojo{Bromuro de selenio (IV)} \\[2pt]
        91 & \ch{SeCl6} & \rojo{Hexacloruro de selenio} & \rojo{Cloruro de selenio (VI)} \\[2pt]
        92 & \ch{AsF5} & \rojo{Pentafluoruro de arsénico} & \rojo{Fluoruro de arsénico (V)} \\[2pt]
        93 & \ch{CCl4} & \rojo{Tetracloruro de carbono} & \rojo{Cloruro de carbono (IV)} \\[2pt]
        94 & \ch{SiF4} & \rojo{Tetrafluoruro de silicio} & \rojo{Fluoruro de silicio} \\
        \bottomrule
    \end{tabular}
\end{center}

\begin{center}
    \begin{tabular}{@{}r C{1.9cm} L{5cm} L{5cm}@{}}
        \cabeceraSolFormES{Sales binarias (continuación)}
        95 & \rojo{\ch{LiCl}} & \rojo{Cloruro de litio} & Cloruro de litio \\[2pt]
        96 & \rojo{\ch{MgI2}} & Diyoduro de magnesio & \rojo{Yoduro de magnesio} \\[2pt]
        97 & \rojo{\ch{BeBr2}} & \rojo{Dibromuro de berilio} & Bromuro de berilio \\[2pt]
        98 & \rojo{\ch{K2S}} & Sulfuro de dipotasio & \rojo{Sulfuro de potasio} \\[2pt]
        99 & \rojo{\ch{FeF2}} & \rojo{Difluoruro de hierro} & Fluoruro de hierro (II) \\[2pt]
        100 & \rojo{\ch{PbCl4}} & Tetracloruro de plomo & \rojo{Cloruro de plomo (IV)} \\[2pt]
        101 & \rojo{\ch{Au2S3}} & \rojo{Trisulfuro de dioro} & Sulfuro de oro (III) \\[2pt]
        102 & \rojo{\ch{SnS}} & \rojo{Sulfuro de estaño (monosulfuro de estaño)} & Sulfuro de estaño (II) \\[2pt]
        103 & \rojo{\ch{ZnBr2}} & Dibromuro de cinc & \rojo{Bromuro de cinc} \\[2pt]
        104 & \rojo{\ch{HgTe}} & \rojo{Telururo de mercurio} & Telururo de mercurio (II) \\[2pt]
        105 & \rojo{\ch{CoI2}} & \rojo{Diyoduro de cobalto} & Yoduro de cobalto (II) \\[2pt]
        106 & \rojo{\ch{CaCl2}} & \rojo{Dicloruro de calcio} & Cloruro de calcio \\[2pt]
        107 & \rojo{\ch{Al2S3}} & Trisulfuro de dialuminio & \rojo{Sulfuro de aluminio} \\[2pt]
        108 & \rojo{\ch{Cr2S3}} & Trisulfuro de dicromo & \rojo{Sulfuro de cromo (III)} \\[2pt]
        109 & \rojo{\ch{MnBr2}} & \rojo{Dibromuro de manganeso} & Bromuro de manganeso (II) \\[2pt]
        110 & \rojo{\ch{Na2Te}} & Telururo de disodio & \rojo{Telururo de sodio} \\[2pt]
        111 & \rojo{\ch{ZnCl2}} & \rojo{Dicloruro de cinc} & Cloruro de cinc \\[2pt]
        112 & \rojo{\ch{Ag2S}} & Sulfuro de diplata & \rojo{Sulfuro de plata} \\[2pt]
        113 & \rojo{\ch{BrF3}} & Trifluoruro de bromo & \rojo{Fluoruro de bromo (III)} \\[2pt]
        114 & \rojo{\ch{BrCl}} & \rojo{Cloruro de bromo} & Cloruro de bromo (I) \\[2pt]
        115 & \rojo{\ch{SF6}} & Hexafluoruro de azufre & \rojo{Fluoruro de azufre (VI)} \\[2pt]
        116 & \rojo{\ch{CS2}} & \rojo{Disulfuro de carbono} & Sulfuro de carbono (IV) \\
        \bottomrule
    \end{tabular}
\end{center}

\begin{center}
    \begin{tabular}{@{}r C{1.9cm} L{5cm} L{5cm}@{}}
        \cabeceraSolFormES{Hidróxidos}
        117 & \ch{CsOH} & \rojo{Hidróxido de cesio} & \rojo{Hidróxido de cesio} \\[2pt]
        118 & \ch{KOH} & \rojo{Hidróxido de potasio} & \rojo{Hidróxido de potasio} \\[2pt]
        119 & \ch{Be(OH)2} & \rojo{Dihidróxido de berilio} & \rojo{Hidróxido de berilio} \\[2pt]
        120 & \ch{Fe(OH)3} & \rojo{Trihidróxido de hierro} & \rojo{Hidróxido de hierro (III)} \\[2pt]
        121 & \ch{AgOH} & \rojo{Hidróxido de plata} & \rojo{Hidróxido de plata} \\[2pt]
        122 & \ch{Al(OH)3} & \rojo{Trihidróxido de aluminio} & \rojo{Hidróxido de aluminio} \\[2pt]
        123 & \ch{NH4OH} & \rojo{Hidróxido de amonio} & \rojo{Hidróxido de amonio} \\[2pt]
        124 & \ch{Cd(OH)2} & \rojo{Dihidróxido de cadmio} & \rojo{Hidróxido de cadmio} \\[2pt]
        125 & \ch{Au(OH)3} & \rojo{Trihidróxido de oro} & \rojo{Hidróxido de oro (III)} \\[2pt]
        126 & \ch{Mn(OH)3} & \rojo{Trihidróxido de manganeso} & \rojo{Hidróxido de manganeso (III)} \\[2pt]
        127 & \ch{Pt(OH)2} & \rojo{Dihidróxido de platino} & \rojo{Hidróxido de platino (II)} \\[2pt]
        128 & \rojo{\ch{Ni(OH)2}} & Dihidróxido de níquel & \rojo{Hidróxido de níquel (II)} \\[2pt]
        129 & \rojo{\ch{Ni(OH)3}} & \rojo{Trihidróxido de níquel} & Hidróxido de níquel (III) \\[2pt]
        130 & \rojo{\ch{NaOH}} & \rojo{Hidróxido de sodio} & Hidróxido de sodio \\[2pt]
        131 & \rojo{\ch{Ca(OH)2}} & \rojo{Dihidróxido de calcio} & Hidróxido de calcio \\[2pt]
        132 & \rojo{\ch{Mg(OH)2}} & Dihidróxido de magnesio & \rojo{Hidróxido de magnesio} \\[2pt]
        133 & \rojo{\ch{Hg(OH)2}} & \rojo{Dihidróxido de mercurio} & Hidróxido de mercurio (II) \\[2pt]
        134 & \rojo{\ch{Fe(OH)2}} & Dihidróxido de hierro & \rojo{Hidróxido de hierro (II)} \\[2pt]
        135 & \rojo{\ch{Zn(OH)2}} & \rojo{Dihidróxido de cinc} & Hidróxido de cinc \\[2pt]
        136 & \rojo{\ch{Cu(OH)2}} & Dihidróxido de cobre & \rojo{Hidróxido de cobre (II)} \\[2pt]
        137 & \rojo{\ch{Pb(OH)4}} & \rojo{Tetrahidróxido de plomo} & Hidróxido de plomo (IV) \\[2pt]
        138 & \rojo{\ch{Sn(OH)2}} & \rojo{Dihidróxido de estaño} & Hidróxido de estaño (II) \\[2pt]
        139 & \rojo{\ch{Ni(OH)3}} & \rojo{Trihidróxido de níquel} & Hidróxido de níquel (III) \\[2pt]
        140 & \rojo{\ch{Cr(OH)3}} & Trihidróxido de cromo & \rojo{Hidróxido de cromo (III)} \\[2pt]
        141 & \rojo{\ch{Cr(OH)2}} & \rojo{Dihidróxido de cromo} & Hidróxido de cromo (II) \\
        \bottomrule
    \end{tabular}
\end{center}

\begin{center}
    \begin{tabular}{@{}r C{1.9cm} L{5cm} L{5cm}@{}}
        \cabeceraSolFormES{Peróxidos}
        142 & \ch{Na2O2} & \rojo{Peróxido de sodio} & -- \\[2pt]
        143 & \ch{K2O2} & \rojo{Peróxido de potasio} & -- \\[2pt]
        144 & \ch{SrO2} & \rojo{Peróxido de estroncio} & -- \\[2pt]
        145 & \ch{CaO2} & \rojo{Peróxido de calcio} & -- \\[2pt]
        146 & \ch{Rb2O2} & \rojo{Peróxido de rubidio} & -- \\[2pt]
        147 & \ch{CuO2} & \rojo{Peróxido de cobre (II)} & -- \\[2pt]
        148 & \ch{MgO2} & \rojo{Peróxido de magnesio} & -- \\[2pt]
        149 & \ch{BaO2} & \rojo{Peróxido de bario} & -- \\
        \midrule
        \multicolumn{4}{c}{\textbf{Hidruros de los grupos 16 y 17 en disolución acuosa}} \\
        \midrule
        150 & \ch{HCl(ac)} & \rojo{Ácido clorhídrico} & -- \\[2pt]
        151 & \ch{HBr(ac)} & \rojo{Ácido bromhídrico} & -- \\[2pt]
        152 & \ch{HI(ac)} & \rojo{Ácido yodhídrico} & -- \\[2pt]
        153 & \ch{HF(ac)} & \rojo{Ácido fluorhídrico} & -- \\[2pt]
        154 & \ch{H2S(ac)} & \rojo{Ácido sulfhídrico} & -- \\[2pt]
        155 & \ch{H2Te(ac)} & \rojo{Ácido telurhídrico} & -- \\[2pt]
        156 & \ch{H2Se(ac)} & \rojo{Ácido selenhídrico} & -- \\
        \bottomrule
    \end{tabular}
\end{center}
}
\fi
            \end{solu}
        \end{exercise}
        \setcounter{numeroformula}{0}
        \begin{longtable}{@{}c C{2cm} C{12cm} C{12cm}@{}}
            \caption{Binary compounds}\label{tab:binary_compounds} \\
            \toprule
            \textbf{No.} & \textbf{Formula} & \textbf{Multipliers} & \textbf{Oxidation State} \\
            \midrule
            \endfirsthead
            \toprule
            \textbf{No.} & \textbf{Formula} & \textbf{Multipliers} & \textbf{Oxidation State} \\
            \midrule
            \endhead
            \bottomrule
            \endlastfoot
            \multicolumn{4}{c}{\textbf{Oxides}} \\[3pt]
            \myorder{} & \ch{Na2O} & \rojo{} & \rojo{} \\[3pt]
            \myorder{} & \ch{Co2O3} & \rojo{} & \rojo{} \\[3pt]
            \myorder{} & \ch{CuO} & \rojo{} & \rojo{} \\[3pt]
            \myorder{} & \ch{Ag2O} & \rojo{} & \rojo{} \\[3pt]
            \myorder{} & \ch{CdO} & \rojo{} & \rojo{} \\[3pt]
            \myorder{} & \ch{NO} & \rojo{} & \rojo{} \\[3pt]
            \myorder{} & \ch{CO2} & \rojo{} & \rojo{} \\[3pt]
            \myorder{} & \ch{As2O5} & \rojo{} & \rojo{} \\[3pt]
            \myorder{} & \ch{SO3} & \rojo{} & \rojo{} \\[3pt]
            \myorder{} & \ch{O5I2} & \rojo{} & \rojo{} \\[3pt]
            \myorder{} & \ch{O7Br2} & \rojo{} & \rojo{} \\[3pt]
            \myorder{} & \rojo{\ch{}} & Dilithium oxide & \rojo{} \\[3pt]
            \myorder{} & \rojo{\ch{}} & \rojo{} & Magnesium oxide \\[3pt]
            \myorder{} & \rojo{\ch{}} & \rojo{} & Potassium oxide \\[3pt]
            \myorder{} & \rojo{\ch{}} & \rojo{} & Strontium oxide \\[3pt]
            \myorder{} & \rojo{\ch{}} & \rojo{} & Iron(II) oxide \\[3pt]
            \myorder{} & \rojo{\ch{}} & Lead dioxide & \rojo{} \\[3pt]
            \myorder{} & \rojo{\ch{}} & \rojo{} & Gold(III) oxide \\[3pt]
            \myorder{} & \rojo{\ch{}} & Dialuminum trioxide & \rojo{} \\[3pt]
            \myorder{} & \rojo{\ch{}} & \rojo{} & Tin(II) oxide \\[3pt]
            \myorder{} & \rojo{\ch{}} & \rojo{} & Zinc oxide \\[3pt]
            \myorder{} & \rojo{\ch{}} & \rojo{} & Mercury(II) oxide \\[3pt]
            \myorder{} & \rojo{\ch{}} & Dinitrogen tetraoxide & \rojo{} \\[3pt]
            \myorder{} & \rojo{\ch{}} & \rojo{} & Cobalt(III) oxide \\[3pt]
            \myorder{} & \rojo{\ch{}} & Dinitrogen pentoxide & \rojo{} \\[3pt]
            \myorder{} & \rojo{\ch{}} & \rojo{} & Barium oxide \\[3pt]
            \myorder{} & \rojo{\ch{}} & Sulfur trioxide & \rojo{} \\[3pt]
            \myorder{} & \rojo{\ch{}} & Dichromium trioxide & \rojo{} \\[3pt]
            \myorder{} & \rojo{\ch{}} & \rojo{} & Manganese(IV) oxide \\[3pt]
            \myorder{} & \rojo{\ch{}} & Dinitrogen monoxide & \rojo{} \\[3pt]
            \myorder{} & \rojo{\ch{}} & \rojo{} & Platinum(IV) oxide \\[3pt]
            \myorder{} & \rojo{\ch{}} & Trioxygen dibromide & \rojo{} \\[3pt]
            \multicolumn{4}{c}{\textbf{Hydrides}} \\[3pt]
            \myorder{} & \ch{NaH} & \rojo{} & \rojo{} \\[3pt]
            \myorder{} & \ch{CaH2} & \rojo{} & \rojo{} \\[3pt]
            \myorder{} & \ch{CoH2} & \rojo{} & \rojo{} \\[3pt]
            \myorder{} & \ch{CuH2} & \rojo{} & \rojo{} \\[3pt]
            \myorder{} & \ch{SnH4} & \rojo{} & \rojo{} \\[3pt]
            \myorder{} & \ch{LiH} & \rojo{} & \rojo{} \\[3pt]
            \myorder{} & \ch{HgH2} & \rojo{} & \rojo{} \\[3pt]
            \myorder{} & \ch{NiH3} & \rojo{} & \rojo{} \\[3pt]
            \myorder{} & \ch{AlH3} & \rojo{} & \rojo{} \\[3pt]
            \myorder{} & \ch{H4Pb} & \rojo{} & \rojo{} \\[3pt]
            \myorder{} & \ch{H3Mn} & \rojo{} & \rojo{} \\[3pt]
            \myorder{} & \rojo{\ch{}} & \rojo{} & Magnesium hydride \\[3pt]
            \myorder{} & \rojo{\ch{}} & \rojo{} & Sodium hydride \\[3pt]
            \myorder{} & \rojo{\ch{}} & \rojo{} & Strontium hydride \\[3pt]
            \myorder{} & \rojo{\ch{}} & \rojo{} & Chromium(II) hydride \\[3pt]
            \myorder{} & \rojo{\ch{}} & Iron trihydride & \rojo{} \\[3pt]
            \myorder{} & \rojo{\ch{}} & \rojo{} & Lead(IV) hydride \\[3pt]
            \myorder{} & \rojo{\ch{}} & Zinc dihydride & \rojo{} \\[3pt]
            \myorder{} & \rojo{\ch{}} & \rojo{} & Cadmium hydride \\[3pt]
            \myorder{} & \rojo{\ch{}} & \rojo{} & Copper(I) hydride \\[3pt]
            \myorder{} & \rojo{\ch{}} & \rojo{} & Gold(III) hydride \\[3pt]
            \myorder{} & \rojo{\ch{}} & \rojo{} & Tin(II) hydride \\[3pt]
            \myorder{} & \rojo{\ch{}} & Potassium hydride & \rojo{} \\[3pt]
            \myorder{} & \rojo{\ch{}} & \rojo{} & Cobalt(II) hydride \\[3pt]
            \myorder{} & \rojo{\ch{}} & \rojo{} & Platinum(IV) hydride \\[3pt]
            \myorder{} & \rojo{\ch{}} & Barium dihydride & \rojo{} \\[3pt]
            \myorder{} & \rojo{\ch{}} & Silver hydride & \rojo{} \\[3pt]
            \myorder{} & \rojo{\ch{}} & \rojo{} & Gold(I) hydride \\[3pt]
            \multicolumn{4}{c}{\textbf{Acids}} \\[3pt]
            \myorder{} & \ch{HCl} & \rojo{} & \rojo{} \\[3pt]
            \myorder{} & \ch{HBr} & \rojo{} & \rojo{} \\[3pt]
            \myorder{} & \ch{HI} & \rojo{} & \rojo{} \\[3pt]
            \myorder{} & \ch{HF} & \rojo{} & \rojo{} \\[3pt]
            \myorder{} & \ch{H2S} & \rojo{} & \rojo{} \\[3pt]
            \myorder{} & \ch{H2Te} & \rojo{} & \rojo{} \\[3pt]
            \myorder{} & \ch{H2Se} & \rojo{} & \rojo{} \\[3pt]
            \myorder{} & \rojo{\ch{}} & Hydrogen chloride &  \\[3pt]
            \myorder{} & \rojo{\ch{}} & \rojo{} & Hydrogen bromide \\[3pt]
            \myorder{} & \rojo{\ch{}} & \rojo{} & Hydrogen iodide \\[3pt]
            \myorder{} & \rojo{\ch{}} & \rojo{} & Hydrogen fluoride \\[3pt]
            \myorder{} & \rojo{\ch{}} & Dihydrogen sulfide & \rojo{} \\[3pt]
            \myorder{} & \rojo{\ch{}} & \rojo{} & Hydrogen telluride \\[3pt]
            \myorder{} & \rojo{\ch{}} & \rojo{} & Hydrogen selenide \\[3pt]
            \multicolumn{4}{c}{\textbf{Binary salts}} \\[3pt]
            \myorder{} & \ch{MgCl2} & \rojo{} & \rojo{} \\[3pt]
            \myorder{} & \ch{BaBr2} & \rojo{} & \rojo{} \\[3pt]
            \myorder{} & \ch{Na3N} & \rojo{} & \rojo{} \\[3pt]
            \myorder{} & \ch{K2S} & \rojo{} & \rojo{} \\[3pt]
            \myorder{} & \ch{CaI2} & \rojo{} & \rojo{} \\[3pt]
            \myorder{} & \ch{FeBr3} & \rojo{} & \rojo{} \\[3pt]
            \myorder{} & \ch{PbS2} & \rojo{} & \rojo{} \\[3pt]
            \myorder{} & \ch{Au2S3} & \rojo{} & \rojo{} \\[3pt]
            \myorder{} & \ch{HgF2} & \rojo{} & \rojo{} \\[3pt]
            \myorder{} & \ch{AlCl3} & \rojo{} & \rojo{} \\[3pt]
            \myorder{} & \ch{SnBr2} & \rojo{} & \rojo{} \\[3pt]
            \myorder{} & \ch{CuI} & \rojo{} & \rojo{} \\[3pt]
            \myorder{} & \ch{CI4} & \rojo{} & \rojo{} \\[3pt]
            \myorder{} & \ch{NF3} & \rojo{} & \rojo{} \\[3pt]
            \myorder{} & \ch{IBr} & \rojo{} &  \\[3pt]
            \myorder{} & \ch{SeBr4} & \rojo{} & \rojo{} \\[3pt]
            \myorder{} & \ch{SeCl6} & \rojo{} & \rojo{} \\[3pt]
            \myorder{} & \ch{AsF5} & \rojo{} &  \\[3pt]
            \myorder{} & \ch{CCl4} & \rojo{} & \rojo{} \\[3pt]
            \myorder{} & \ch{SiF4} & \rojo{} & \rojo{} \\[3pt]
            \myorder{} & \rojo{\ch{}} & \rojo{} & Lithium chloride \\[3pt]
            \myorder{} & \rojo{\ch{}} & Magnesium diiodide & \rojo{} \\[3pt]
            \myorder{} & \rojo{\ch{}} & \rojo{} & Beryllium bromide \\[3pt]
            \myorder{} & \rojo{\ch{}} & Dipotassium sulfide & \rojo{} \\[3pt]
            \myorder{} & \rojo{\ch{}} & \rojo{} & Iron (II) fluoride \\[3pt]
            \myorder{} & \rojo{\ch{}} & Lead tetrachloride & \rojo{} \\[3pt]
            \myorder{} & \rojo{\ch{}} & \rojo{} & Gold (III) sulfide \\[3pt]
            \myorder{} & \rojo{\ch{}} & \rojo{} & Tin (II) sulfide \\[3pt]
            \myorder{} & \rojo{\ch{}} & Zinc dibromide & \rojo{} \\[3pt]
            \myorder{} & \rojo{\ch{}} & \rojo{} & Mercury (II) telluride \\[3pt]
            \myorder{} & \rojo{\ch{}} & \rojo{} & Cobalt (II) iodide \\[3pt]
            \myorder{} & \rojo{\ch{}} & \rojo{} & Calcium chloride \\[3pt]
            \myorder{} & \rojo{\ch{}} & Dialuminum trisulfide & \rojo{} \\[3pt]
            \myorder{} & \rojo{\ch{}} & Dichromium trisulfide & \rojo{} \\[3pt]
            \myorder{} & \rojo{\ch{}} & \rojo{} & Manganese (II) bromide \\[3pt]
            \myorder{} & \rojo{\ch{}} & Disodium telluride & \rojo{} \\[3pt]
            \myorder{} & \rojo{\ch{}} & \rojo{} & Zinc chloride \\[3pt]
            \myorder{} & \rojo{\ch{}} & Disilver sulfide & \rojo{} \\[3pt]
            \myorder{} & \rojo{\ch{}} & Bromine trifluoride & \rojo{} \\[3pt]
            \myorder{} & \rojo{\ch{}} & \rojo{} & Bromine (I) chloride \\[3pt]
            \myorder{} & \rojo{\ch{}} & Sulfur hexafluoride & \rojo{} \\[3pt]
            \myorder{} & \rojo{\ch{}} & \rojo{} & Carbon (IV) sulfide \\[3pt]
            \multicolumn{4}{c}{\textbf{Hydroxides}} \\[3pt]
            \myorder{} & \ch{CsOH} & \rojo{} & \rojo{} \\[3pt]
            \myorder{} & \ch{KOH} & \rojo{} & \rojo{} \\[3pt]
            \myorder{} & \ch{Be(OH)2} & \rojo{} & \rojo{} \\[3pt]
            \myorder{} & \ch{Fe(OH)3} & \rojo{} & \rojo{} \\[3pt]
            \myorder{} & \ch{AgOH} & \rojo{} & \rojo{} \\[3pt]
            \myorder{} & \ch{Al(OH)3} & \rojo{} & \rojo{} \\[3pt]
            \myorder{} & \ch{NH4OH} & \rojo{} & \rojo{} \\[3pt]
            \myorder{} & \ch{Cd(OH)2} & \rojo{} & \rojo{} \\[3pt]
            \myorder{} & \ch{Au(OH)3} & \rojo{} & \rojo{} \\[3pt]
            \myorder{} & \ch{Mn(OH)3} & \rojo{} & \rojo{} \\[3pt]
            \myorder{} & \ch{Pt(OH)2} & \rojo{} & \rojo{} \\[3pt]
            \myorder{} & \rojo{\ch{}} & Nickel dihydroxide &  \\[3pt]
            \myorder{} & \rojo{\ch{}} & \rojo{} & Nickel (III) hydroxide \\[3pt]
            \myorder{} & \rojo{\ch{}} & \rojo{} & Sodium hydroxide \\[3pt]
            \myorder{} & \rojo{\ch{}} & \rojo{} & Calcium hydroxide \\[3pt]
            \myorder{} & \rojo{\ch{}} & Magnesium dihydroxide & \rojo{} \\[3pt]
            \myorder{} & \rojo{\ch{}} & \rojo{} & Mercury (II) hydroxide \\[3pt]
            \myorder{} & \rojo{\ch{}} & Iron dihydroxide & \rojo{} \\[3pt]
            \myorder{} & \rojo{\ch{}} & \rojo{} & Zinc hydroxide \\[3pt]
            \myorder{} & \rojo{\ch{}} & Copper dihydroxide & \rojo{} \\[3pt]
            \myorder{} & \rojo{\ch{}} & \rojo{} & Lead (IV) hydroxide \\[3pt]
            \myorder{} & \rojo{\ch{}} & \rojo{} & Tin (II) hydroxide \\[3pt]
            \myorder{} & \rojo{\ch{}} & \rojo{} & Nickel (III) hydroxide \\[3pt]
            \myorder{} & \rojo{\ch{}} & Chromium trihydroxide & \rojo{} \\[3pt]
            \myorder{} & \rojo{\ch{}} & \rojo{} & Chromium (II) hydroxide \\[3pt]
            \multicolumn{4}{c}{\makecell{\textbf{Peroxides}. \\ In this case we have only given the nomenclature for peroxides.       \\ You can write their name in the multipliers column leaving the rest blank}} \\[3pt]
            \myorder{} & \ch{Na2O2} & \rojo{} & -- \\[3pt]
            \myorder{} & \ch{K2O2} & \rojo{} & -- \\[3pt]
            \myorder{} & \ch{SrO2} & \rojo{} & -- \\[3pt]
            \myorder{} & \ch{CaO2} & \rojo{} & -- \\[3pt]
            \myorder{} & \ch{Rb2O2} & \rojo{} & -- \\[3pt]
            \myorder{} & \ch{CuO2} &  & -- \\[3pt]
            \myorder{} & \ch{MgO2} &  & -- \\[3pt]
            \myorder{} & \ch{BaO2} & -- \\[3pt]
            \multicolumn{4}{c}{\makecell{\textbf{Hydrides from groups 16 and 17 in aqueous solution}. Only one nomenclature has been given.       \\ You can write their name in the multipliers column leaving the rest blank}} \\[3pt]
            \myorder{} & \ch{HCl_{(aq)}} & \rojo{} &  \\[3pt]
            \myorder{} & \ch{HBr_\acu} & \rojo{} &  \\[3pt]
            \myorder{} & \ch{HI_\acu} & \rojo{} &  \\[3pt]
            \myorder{} & \ch{HF_\acu} & \rojo{} &  \\[3pt]
            \myorder{} & \ch{H2S_\acu} & \rojo{} &  \\[3pt]
            \myorder{} & \ch{H2Te_\acu} & \rojo{} &  \\[3pt]
            \myorder{} & \ch{H2Se_\acu} & \rojo{} &  \\
        \end{longtable}
    \end{landscape}
% Restaurar los márgenes originales
\restoregeometry

\pagestyle{fancy}


%\section{Ejercicios de autoevaluación}
%
%\subsection{Nivel 1: Identificar el tipo de compuesto}
%
%Indica qué tipo de compuesto es cada uno:
%
%\begin{enumerate}
%    \item \ch{CaO} → \_\_\_\_\_\_\_\_\_\_\_\_\_\_
%    \item \ch{NaH} → \_\_\_\_\_\_\_\_\_\_\_\_\_\_
%    \item \ch{HCl(ac)} → \_\_\_\_\_\_\_\_\_\_\_\_\_\_
%    \item \ch{KOH} → \_\_\_\_\_\_\_\_\_\_\_\_\_\_
%    \item \ch{Na2O2} → \_\_\_\_\_\_\_\_\_\_\_\_\_\_
%    \item \ch{FeCl3} → \_\_\_\_\_\_\_\_\_\_\_\_\_\_
%\end{enumerate}
%
%\subsection{Nivel 2: Nomenclatura (de fórmula a nombre)}
%
%Nombra los siguientes compuestos usando \textbf{nomenclatura de multiplicadores}:
%
%\begin{enumerate}
%    \item \ch{MgO} → \_\_\_\_\_\_\_\_\_\_\_\_\_\_\_\_\_\_\_\_\_\_
%    \item \ch{Al2O3} → \_\_\_\_\_\_\_\_\_\_\_\_\_\_\_\_\_\_\_\_\_\_
%    \item \ch{LiH} → \_\_\_\_\_\_\_\_\_\_\_\_\_\_\_\_\_\_\_\_\_\_
%    \item \ch{H2S(ac)} → \_\_\_\_\_\_\_\_\_\_\_\_\_\_\_\_\_\_\_\_\_\_
%    \item \ch{CaCl2} → \_\_\_\_\_\_\_\_\_\_\_\_\_\_\_\_\_\_\_\_\_\_
%    \item \ch{Fe(OH)3} → \_\_\_\_\_\_\_\_\_\_\_\_\_\_\_\_\_\_\_\_\_\_
%\end{enumerate}
%
%Nombra los siguientes compuestos usando \textbf{nomenclatura con estados de oxidación}:
%
%\begin{enumerate}[resume]
%        \item \ch{CuO} → \_\_\_\_\_\_\_\_\_\_\_\_\_\_\_\_\_\_\_\_\_\_
%        \item \ch{Cu2O} → \_\_\_\_\_\_\_\_\_\_\_\_\_\_\_\_\_\_\_\_\_\_
%        \item \ch{FeH2} → \_\_\_\_\_\_\_\_\_\_\_\_\_\_\_\_\_\_\_\_\_\_
%        \item \ch{FeH3} → \_\_\_\_\_\_\_\_\_\_\_\_\_\_\_\_\_\_\_\_\_\_
%\end{enumerate}
%
%\subsection{Nivel 3: Formulación (de nombre a fórmula)}
%
%Escribe la fórmula de los siguientes compuestos:
%
%\begin{enumerate}
%    \item Óxido de sodio → \_\_\_\_\_\_\_\_\_\_
%    \item Trióxido de dihierro → \_\_\_\_\_\_\_\_\_\_
%    \item Dihidruro de calcio → \_\_\_\_\_\_\_\_\_\_
%    \item Ácido bromhídrico → \_\_\_\_\_\_\_\_\_\_
%    \item Cloruro de aluminio → \_\_\_\_\_\_\_\_\_\_
%    \item Trihidróxido de aluminio → \_\_\_\_\_\_\_\_\_\_
%    \item Peróxido de potasio → \_\_\_\_\_\_\_\_\_\_
%    \item Óxido de hierro (III) → \_\_\_\_\_\_\_\_\_\_
%    \item Hidruro de cobre (I) → \_\_\_\_\_\_\_\_\_\_
%    \item Sulfuro de hierro (II) → \_\_\_\_\_\_\_\_\_\_
%\end{enumerate}
%
%\subsection{Nivel 4: Ejercicios mixtos}
%
%Completa la tabla:
%
%\begin{center}
%    \begin{tabular}{|l|c|c|}
%        \hline
%        \textbf{Fórmula} & \textbf{Nombre (multiplicadores)} & \textbf{Nombre (est. oxidación)} \\
%        \hline
%        \ch{SO3}         &                                   &                                  \\
%        \hline
%        & Dióxido de carbono                &                                  \\
%        \hline
%        &                                   & Óxido de plomo (IV)              \\
%        \hline
%        \ch{SnH4}        &                                   &                                  \\
%        \hline
%        & Trihidruro de hierro              &                                  \\
%        \hline
%    \end{tabular}
%\end{center}
%
%\subsection{Soluciones a los ejercicios}
%
%\subsection{Soluciones Nivel 1}
%
%\begin{enumerate}
%    \item \ch{CaO} → Óxido
%    \item \ch{NaH} → Hidruro metálico
%    \item \ch{HCl(ac)} → Hidrácido
%    \item \ch{KOH} → Hidróxido
%    \item \ch{Na2O2} → Peróxido
%    \item \ch{FeCl3} → Sal binaria
%\end{enumerate}
%
%\subsection{Soluciones Nivel 2}
%
%\textbf{Nomenclatura de multiplicadores}:
%\begin{enumerate}
%    \item \ch{MgO} → Óxido de magnesio (o monóxido de magnesio)
%    \item \ch{Al2O3} → Trióxido de dialuminio
%    \item \ch{LiH} → Hidruro de litio
%    \item \ch{H2S(ac)} → Ácido sulfhídrico
%    \item \ch{CaCl2} → Dicloruro de calcio
%    \item \ch{Fe(OH)3} → Trihidróxido de hierro
%\end{enumerate}
%
%\textbf{Nomenclatura con estados de oxidación}:
%\begin{enumerate}[resume]
%        \item \ch{CuO} → Óxido de cobre (II)
%        \item \ch{Cu2O} → Óxido de cobre (I)
%        \item \ch{FeH2} → Hidruro de hierro (II)
%        \item \ch{FeH3} → Hidruro de hierro (III)
%\end{enumerate}
%
%\subsection{Soluciones Nivel 3}
%
%\begin{enumerate}
%    \item Óxido de sodio → \ch{Na2O}
%    \item Trióxido de dihierro → \ch{Fe2O3}
%    \item Dihidruro de calcio → \ch{CaH2}
%    \item Ácido bromhídrico → \ch{HBr(ac)}
%    \item Cloruro de aluminio → \ch{AlCl3}
%    \item Trihidróxido de aluminio → \ch{Al(OH)3}
%    \item Peróxido de potasio → \ch{K2O2}
%    \item Óxido de hierro (III) → \ch{Fe2O3}
%    \item Hidruro de cobre (I) → \ch{CuH}
%    \item Sulfuro de hierro (II) → \ch{FeS}
%\end{enumerate}
%
%\subsection{Soluciones Nivel 4}
%
%\begin{center}
%    \begin{tabular}{|l|l|l|}
%        \hline
%        \textbf{Fórmula} & \textbf{Nombre (multiplicadores)} & \textbf{Nombre (est. oxidación)} \\
%        \hline
%        \ch{SO3}         & Trióxido de azufre                & Óxido de azufre (VI)             \\
%        \hline
%        \ch{CO2}         & Dióxido de carbono                & Óxido de carbono (IV)            \\
%        \hline
%        \ch{PbO2}        & Dióxido de plomo                  & Óxido de plomo (IV)              \\
%        \hline
%        \ch{SnH4}        & Tetrahidruro de estaño            & Hidruro de estaño (IV)           \\
%        \hline
%        \ch{FeH3}        & Trihidruro de hierro              & Hidruro de hierro (III)          \\
%        \hline
%    \end{tabular}
%\end{center}
%
%\newpage
%\pagestyle{empty}
%\begin{landscape}
%
%    \begin{figure}[H]
%        \centering
%        \includegraphics[scale=0.51]{images/2_3_ESO_TP_ES}
%        \caption{}
%        \label{fig:}
%    \end{figure}
%
%\end{landscape}
%
%\newpage
%
%\begin{landscape}
%
%    \begin{figure}[H]
%        \centering
%        \includegraphics[scale=0.51]{images/2_3_ESO_TP_EN}
%        \caption{}
%        \label{fig:}
%    \end{figure}
%
%\end{landscape}
%
%\newpage
%
%\begin{landscape}
%
%    \begin{figure}[H]
%        \centering
%        \includegraphics[scale=0.51]{images/TP_Valencias}
%        \caption{}
%        \label{fig:}
%    \end{figure}
%\end{landscape}
%\pagestyle{fancy}
%
